\documentclass[10pt,a4paper]{article}
\usepackage[utf8]{inputenc}
\usepackage[T1]{fontenc}
\usepackage[a4paper,top=20mm,bottom=20mm,left=16mm,right=16mm,headheight=14pt]{geometry}
\usepackage{fancyhdr}
\usepackage{booktabs}
\usepackage{longtable}
\usepackage{array}
\usepackage{enumitem}
\usepackage{amssymb}
\usepackage{xcolor}
\usepackage[hyphens]{url}
\usepackage[colorlinks=true,urlcolor=blue!60!black,linkcolor=blue!60!black]{hyperref}

\setlength{\parindent}{0pt}
\setlength{\parskip}{3pt}
\renewcommand{\arraystretch}{1.28}

\pagestyle{fancy}
\fancyhf{}
\fancyhead[L]{\small LANGUAGE \& SOCIETY}
\fancyhead[R]{\small TRANSCRIPT ANALYSIS WORKSHEET}
\fancyfoot[C]{\small\thepage}
\renewcommand{\headrulewidth}{0.4pt}

\begin{document}

\begin{center}
    {\large\bfseries LANGUAGE \& SOCIETY PROJECT}\\[3pt]
    {\LARGE\bfseries Transcript Analysis Worksheet}\\[3pt]
    {\normalsize Qualitative analysis of political news discourse}
\end{center}
\vspace{2pt}
\noindent\rule{\textwidth}{0.4pt}
\vspace{4pt}


\section*{Creator 1: Sarthak Goswami (\textit{Sunday Show} --- Independent Satirical Infotainment)}
\begin{itemize}[leftmargin=1.5em, itemsep=3pt, topsep=4pt]
  \item \textbf{Analyst:} Baani Solanki (\texttt{2025114016})
  \item \textbf{Channel / Outlet:} Sarthak Goswami (\texttt{@SundaySarthak}, YouTube)
  \item \textbf{Ideological / Editorial Positioning:} Centrist-liberal satirical infotainment; critiques ruling-party majoritarianism ("dharma ka dhandha"), opposition opportunism ("sasti BJP"), bureaucratic corruption, and mainstream "godi media" polarization while foregrounding civic constitutionalism and youth agency.
  \item \textbf{Transcript Files Referenced:}
  \begin{itemize}[leftmargin=1.4em, itemsep=2pt, topsep=2pt]
    \item Event 1: \texttt{transcripts/event1\_ram\_mandir/01\_sarthak\_goswami\_lb7gZnyZH0c.txt}
    \item Event 2: \texttt{transcripts/event2\_neet\_protest/01\_sarthak\_goswami\_EM1GeZufO6E.txt}
  \end{itemize}
\end{itemize}
\vspace{4pt}\noindent\rule{\textwidth}{0.4pt}\vspace{4pt}

\section*{Guidelines: Five features to look for in every event}
\begin{longtable}{@{}>{\raggedright\arraybackslash}p{0.17\textwidth} >{\raggedright\arraybackslash}p{0.20\textwidth} >{\raggedright\arraybackslash}p{0.24\textwidth} >{\raggedright\arraybackslash}p{0.32\textwidth}@{}}
\toprule
\textbf{Feature} & \textbf{What to look for} & \textbf{Illustrative example} & \textbf{Questions to answer} \\
\midrule
\endhead
\textbf{1. Naming} & What are people, groups, actions, or events called? & "Students", "protesters", "agitators"; "irregularity", "scam", "controversy". & What exact label is used? Who or what does it refer to? Who uses it? What alternative label could have been used? \\
\addlinespace[4pt]
\textbf{2. Presupposition} & What is treated as already true or taken for granted? & "Why are students creating chaos again?" may presuppose earlier disruption. & What does the wording assume? Is that assumption supported, disputed, or left unexamined? Does it appear in a question or words such as "again", "still", or "finally"? \\
\addlinespace[4pt]
\textbf{3. Agency} & Who is shown doing something? Who is acted upon? Who disappears? & "Police injured protesters" vs. "Protesters were injured." & Who is the actor? Who experiences the action? Is responsibility named, backgrounded, or omitted? \\
\addlinespace[4pt]
\textbf{4. Antagonisation} & Who is positioned against whom? & "Ordinary citizens" vs. "political troublemakers". & Which groups are contrasted? Is one presented as more legitimate, moral, patriotic, dangerous, or disloyal? Is an "us vs. them" division implied? \\
\addlinespace[4pt]
\textbf{5. Intensification / softening} & What is exaggerated, minimised, softened, or euphemised? & "Violent mob" vs. "peaceful demonstrators"; "incident" instead of a direct description. & Which words heighten emotion or seriousness? Which minimise it? Are metaphors, dramatic adjectives, hedges, or euphemisms used? Does the wording match the evidence? \\
\bottomrule
\end{longtable}
\textit{Reminder: A label or implication is not automatically false or manipulative. Examine the surrounding context and evidence.}

\vspace{4pt}\noindent\rule{\textwidth}{0.4pt}\vspace{4pt}

\section*{Event1. Ram Mandir donation controversy}
\textbf{Focus:} How is the controversy framed---as financial accountability, religious sentiment, political motivation, or something else? Be precise about the specific allegation or episode covered.

\begin{itemize}[leftmargin=1.5em, itemsep=3pt, topsep=4pt]
  \item \textbf{Outlet:} Sarthak Goswami (\textit{The Sunday Show}, YouTube)
  \item \textbf{Date:} \texttt{2026-07-05}
  \item \textbf{Video title:} \textit{Ram Mandir Scam, Gadkari's Ethanol, Kejriwal's Temple \& More | Sunday Show} (\url{https://youtu.be/lb7gZnyZH0c?si=QvgBq0ACOhcx0pdO})
  \item \textbf{Timestamp(s):} Relevant Event 1 segment \texttt{[00:00--06:30]} (with brief callbacks at \texttt{[27:55--28:15]} and \texttt{[29:15--29:20]}; total video duration \texttt{45:41}; unrelated weekly news segments \texttt{[06:31--45:41]} and Odoo POS sponsor segment \texttt{[07:56--10:28]} excluded per project guidelines. Contextual follow-up in 26 July episode \texttt{EM1GeZufO6E} at \texttt{[27:17--29:45]} also noted).
  \item \textbf{Analyst:} Baani Solanki (\texttt{2025114016})
  \item \textbf{Language:} Hindi / Hinglish (Native Hindi audio, Romanized script; non-AI-dubbed)
\end{itemize}

\subsection*{Event-specific questions (Event 1)}
\begin{enumerate}[leftmargin=1.6em, itemsep=4pt, topsep=4pt]
  \item \textbf{What terms are used: "scam", "alleged irregularity", "corruption", "political controversy", or something else?}
  \begin{itemize}[leftmargin=1.4em, itemsep=2pt, topsep=2pt]
    \item Sarthak uses blunt vernacular terms of financial crime---\texttt{ghotala} [scam] (\texttt{[00:00]}, \texttt{[00:19]}, \texttt{[00:46]}), \texttt{paisa chori} / \texttt{chanda chori} [donation theft] (\texttt{[00:00]}, \texttt{[00:19]}, \texttt{[00:29]}), \texttt{ghapla} [bungling/fraud] (\texttt{[02:21]}), and \texttt{skaimon} [scams] (\texttt{[01:31]})---before reframing the controversy metaphorically as \texttt{asli ghotala} [the real scam] (\texttt{[00:10]}, \texttt{[04:04]}, \texttt{[04:30]}, \texttt{[05:22]}, \texttt{[06:16]}) and \texttt{Ram ka charitra chori} [theft of Ram's character/ethos] (\texttt{[06:26]}). He also coins the satirical political hybrid \texttt{"kangresiphikeshan of BJP politiks"} [Congressification of BJP politics] (\texttt{[00:29]}).
  \end{itemize}
  \item \textbf{How are the temple, trust, donors, and people making allegations described?}
  \begin{itemize}[leftmargin=1.4em, itemsep=2pt, topsep=2pt]
    \item \textbf{Trust \& Accused Staff:} Senior trustee Champat Rai is labelled \texttt{"RSS VHP ka aadmi"} [an RSS-VHP man] (\texttt{[00:46]}) who claims innocence by blaming his driver \texttt{"salman khaan se inspireshan le rahe hain"} [taking inspiration from Salman Khan's hit-and-run defense] (\texttt{[01:31--01:40]}). Driver Ram Shankar "Tinnu" Yadav is simultaneously cited as the alleged \texttt{"mastarmaind"} (\texttt{[00:37]}) and satirized as \texttt{"bali ka bakra"} [scapegoat] (\texttt{[00:46]}). Counting staff like Lavkush Mishra are mocked for treating temple donations as \texttt{"hamare baap ka maal"} [our father's property] (\texttt{[01:14]}).
    \item \textbf{Donors / Public:} Devotees are depicted with biting second-person irony---both as \texttt{"bhakt bechare kamphyooj ho gae hain"} [poor devotees are confused about how to accuse one 'god' of stealing from another God] (\texttt{[01:48--01:56]}) and as complicit citizens (\texttt{"aap bhi bhagidar ho usmen"}, \texttt{[02:29]}; \texttt{"dharma chalane ka bachane ka theka toh tum hi ne diya tha"}, \texttt{[02:46--02:54]}) who outsourced their religion to political contractors. Online hyper-partisans defending the theft are mocked for arguing \texttt{"Hindu ka paisa agar kisi Hindu ne le bhi liya... galat kya hai?"} (\texttt{[02:54--03:12]}).
    \item \textbf{People Making Allegations:} Local institutional actors like the Ayodhya \texttt{"baar esosieshan"} [Bar Association] threatening a city blockade (\texttt{[01:48]}) are cited alongside \texttt{"apojishan ne alag baval kaat diya hai"} [the opposition has raised its own ruckus] (\texttt{[01:48]}).
  \end{itemize}
  \item \textbf{Is the issue framed as financial accountability, an attack on Hindu faith, a political dispute, or a matter of public trust?}
  \begin{itemize}[leftmargin=1.4em, itemsep=2pt, topsep=2pt]
    \item Primarily framed as \textbf{moral-civic accountability and the commercialization/politicization of religion (\texttt{"dharma ka dhandha"}, \texttt{[02:21]})}. Rather than treating the controversy as an external attack on Hindu faith, Sarthak argues that Hindu faith was hollowed out \textit{from within} by self-styled religious gatekeepers (\texttt{"adharmi log jab dharma ka prachar karne lage wo tha asli ghotala"}, \texttt{[04:04]}).
  \end{itemize}
  \item \textbf{Are questions about donations presented as legitimate scrutiny or as an attempt to malign the temple or its supporters?}
  \begin{itemize}[leftmargin=1.4em, itemsep=2pt, topsep=2pt]
    \item Presented unequivocally as \textbf{legitimate, overdue scrutiny}. In fact, Sarthak ridicules partisan apologists who gaslight questioners (\texttt{"tumhare dharma ka theka leke tumko hi loot liya aur ab tum sawal karoge"}, \texttt{[03:12--03:20]}).
  \end{itemize}
  \item \textbf{Are allegations attributed to specific people and supported with evidence, or presented as established facts?}
  \begin{itemize}[leftmargin=1.4em, itemsep=2pt, topsep=2pt]
    \item Hybrid: Financial recoveries and arrests are anchored to the SIT investigation (\texttt{"SIT ki jaanch baithi jinhonne apni report dayar ki jiske baad Trust ke aath logon ko arest kiya gaya"}, \texttt{[00:29--00:37]}) with concrete figures (\texttt{Rs.\,20 lakh} at Avinash Shukla's yoga center, \texttt{Rs.\,12,000} salary vs \texttt{Rs.\,25 lakh} car and \texttt{Rs.\,12 lakh} stable for Lavkush Mishra at \texttt{[01:14--01:23]}). However, his broader claim that senior figures used the driver as a \texttt{"bali ka bakra"} (\texttt{[00:46]}) is presented through satirical inference (\texttt{"ab aap thoda Google karo pata lag jaega"}, \texttt{[00:46]}) rather than adjudicated proof.
  \end{itemize}
  \item \textbf{How are denials and counterarguments represented: explained, dismissed, challenged, or omitted?}
  \begin{itemize}[leftmargin=1.4em, itemsep=2pt, topsep=2pt]
    \item \textbf{Mocked and challenged through satire.} Champat Rai's denial (\texttt{"champat raay ji ka kahna hai ki ji main toh beksoor hoon. mere draivar ne mujhe dhokha diya hai"}, \texttt{[01:31--01:40]}) is immediately ridiculed via the Salman Khan driver-blame trope (\texttt{[01:40]}). Partisan online counterarguments ("if a Hindu took a Hindu's money, why are you upset?") are quoted and dismantled through reductio ad absurdum (\texttt{"kal koi goli maar dega tumhen aur yeh kahenge ki Hindu ka khoon Hindu ne bahaya"}, \texttt{[03:12]}).
  \end{itemize}
  \item \textbf{Does the presenter distinguish criticism of the trust or its financial practices from criticism of the temple or religion itself?}
  \begin{itemize}[leftmargin=1.4em, itemsep=2pt, topsep=2pt]
    \item \textbf{Yes, explicitly and sharply.} He contrasts the ancient tradition of Hinduism (\texttt{"hajaron lakhon saal purana dharma"}, \texttt{[03:38]}) and Lord Ram's moral character against opportunistic political contractors and online abusers (\texttt{"jo Mandir ke bahar chappal churate the unko andar ka theka de diya tumne... Ram ka chanda nahin Ram ka charitra chori ho chuka hai"}, \texttt{[06:16--06:26]}).
  \end{itemize}
  \item \textbf{Who is portrayed as protecting donors' interests, and who as exploiting religious sentiment?}
  \begin{itemize}[leftmargin=1.4em, itemsep=2pt, topsep=2pt]
    \item \textbf{Exploiting religious sentiment:} Trust insiders, VIP-darshan profiteers (\texttt{[02:29--02:38]}), BJP-RSS-linked contractors who peddled the \texttt{"Hindu khatre mein hai"} [Hindus are in danger] tonic (\texttt{[03:20--03:28]}), and abusive online "dharma rakhvala" trolls (\texttt{[04:30--05:22]}).
    \item \textbf{Protecting donors' interests:} No political party is idealized as a savior (later at \texttt{[26:44--27:55]} he equally mocks Kejriwal/AAP as \texttt{"sasti BJP"} for temple-appeasement politics); instead, Sarthak calls upon ordinary Hindu citizens themselves to wake up (\texttt{"Hindu jago"}, \texttt{[06:26]}) and reclaim ethical accountability.
  \end{itemize}
\end{enumerate}

\subsection*{Apply the five shared features (Event 1)}
\begin{longtable}{@{}>{\raggedright\arraybackslash}p{0.19\textwidth} >{\raggedright\arraybackslash}p{0.38\textwidth} >{\raggedright\arraybackslash}p{0.37\textwidth}@{}}
\toprule
\textbf{Feature} & \textbf{Exact wording + timestamp} & \textbf{What it does in context} \\
\midrule
\endhead
\textbf{1. Naming} & \texttt{"champat raay RSS VHP ka aadmi hai aur tinnoo draivar bali ka bakra."} \texttt{[00:46]} \newline  \texttt{"laik vhat ij dis kangresiphikeshan of BJP politiks."} \texttt{[00:29]} & Labels Champat Rai by his ideological patronage network (\texttt{RSS VHP ka aadmi}) while naming arrested driver Tinnu Yadav a \texttt{"bali ka bakra"} (scapegoat), challenging the official narrative that only low-level staff were culpable. \texttt{"Kangresiphikeshan"} equates BJP-era institutional graft with older Congress-era corruption scandals. \\
\addlinespace[4pt]
\textbf{2. Presupposition} & \texttt{"par tum mujhe batao tum sarapraizd kyon ho? ... jab itne salon se dharma ka dhandha ban gaya tha toh usmen ghapla toh hona hi tha."} \texttt{[02:12--02:21]} & The \textit{why}-question (\texttt{tum sarapraizd kyon ho?}) and temporal clause (\texttt{jab itne salon se dharma ka dhandha ban gaya tha}) presuppose as an established fact that religion had already been turned into a commercial racket (\texttt{dhandha}) long before the theft surfaced. \\
\addlinespace[4pt]
\textbf{3. Agency} & \texttt{"dharma chalane ka bachane ka theka toh tum hi ne diya tha in logon ko... aap hi ne nilam kiya aur ab aap boli ka hisab bhi mang rahe ho."} \texttt{[02:46--03:38]} & Shifts agency directly onto the viewer/voter (\texttt{tum hi ne}, \texttt{aap hi ne}) using active ergative constructions (\texttt{ne diya tha}, \texttt{ne nilam kiya}), refusing to let the public pose purely as passive victims of corrupt trustees. \\
\addlinespace[4pt]
\textbf{4. Antagonisation} & \texttt{"jo Mandir ke bahar chappal churate the unko andar ka theka de diya tumne... Ram ka chanda nahin Ram ka charitra chori ho chuka hai."} \texttt{[06:16--06:26]} & Constructs a sharp moral opposition between authentic ethical devotion (\texttt{Ram ka charitra}) and lumpenized, abusive self-appointed guardians of Hindutva (\texttt{chappal churate the}, online abusers chanting \texttt{"jay shri Ram"} while issuing rape threats at \texttt{[04:47--05:22]}). \\
\addlinespace[4pt]
\textbf{5. Intensification / softening} & \textbf{Irony/Hook reversal:} \texttt{"Ram Mandir mein ghotala nahin huaa hai... kyunki asli ghotala hi kahin aur huaa hai."} \texttt{[00:00--00:10]} \newline  \textbf{Intensification:} \texttt{"sher ko pala tha ab uski dait se dikkat... Hindu khatre mein tumhen ghutti pilae ja rahe the."} \texttt{[03:20--03:28]} & Opens with mock-denial (ironic softening) as a rhetorical hook, then intensifies critique through visceral metaphors---raising a predator (\texttt{sher ko pala}), force-feeding fear (\texttt{ghutti pilae}), and auctioning religion (\texttt{nilam kiya}, \texttt{[03:28]}). \\
\bottomrule
\end{longtable}

\subsection*{Short summary (Event 1)}
\begin{itemize}[leftmargin=1.5em, itemsep=3pt, topsep=4pt]
  \item \textbf{What interpretation of the event does this video encourage?} That the Ayodhya donation embezzlement is the inevitable outcome (\texttt{"prakriti ka niyam"}, \texttt{[02:21]}) of mixing electoral politics, unchecked VIP commercialization, and blind majoritarian loyalty (\texttt{[02:12--03:38]}), where junior staff like driver Tinnu Yadav take the fall (\texttt{[00:46]}) while deeper moral corruption---weaponizing Lord Ram's name for power and online abuse (\texttt{[04:30--06:26]})---goes unpunished.
  \item \textbf{Whose perspective is foregrounded, challenged, or missing?} \textit{Foregrounded:} Secular-liberal / self-reflective Hindu citizen critique and victims of online Hindutva harassment (\texttt{[04:38--05:48]}). \textit{Challenged:} Trust leadership's "rogue driver" excuse (\texttt{[01:31--01:40]}) and online bhakt apologia (\texttt{[02:54--03:12]}). \textit{Missing:} Detailed institutional explanation from SBI or SIT officials and the perspective of sincere grassroots temple volunteers.
  \item \textbf{What evidence supports your interpretation? What complicates it?} \textit{Supports:} He cites SIT arrests of 8 staffers (\texttt{[00:37]}), specific asset seizures (\texttt{Rs.\,20 lakh} from Avinash Shukla, Lavkush Mishra's \texttt{Rs.\,12,000} salary vs luxury assets at \texttt{[01:14--01:23]}), and verbatim screenshots/quotes of abusive online Hindutva accounts (\texttt{[04:47--05:22]}). \textit{Complicates:} Over half of the 6.5-minute segment (\texttt{[03:47--06:15]}) pivots away from the Ayodhya Trust's accounting mechanics to an unrelated viral online harassment incident to make a broader cultural point, blurring financial evidence with cultural commentary.
\end{itemize}
\vspace{4pt}\noindent\rule{\textwidth}{0.4pt}\vspace{4pt}

\section*{Event2. Students' protest regarding NEET paper leak}
\textbf{Focus:} How does coverage construct the legitimacy of student protest and responsibility for the underlying problem?

\begin{itemize}[leftmargin=1.5em, itemsep=3pt, topsep=4pt]
  \item \textbf{Outlet:} Sarthak Goswami (\textit{The Sunday Show}, YouTube)
  \item \textbf{Date:} \texttt{2026-07-26}
  \item \textbf{Video title:} \textit{Dharmendra Pradhan Resigned, CJP won, Modi on Instagram \& more | Sunday Show} (\url{https://youtu.be/EM1GeZufO6E?si=PgvQ7GHqNjZJIZET})
  \item \textbf{Timestamp(s):} Relevant Event 2 segments \texttt{[00:00--08:55]} and \texttt{[10:20--23:38]} (total video duration \texttt{33:04}; Kiwi sponsor ads \texttt{[08:56--10:20]}, \texttt{[20:52--21:03]} and unrelated news segments \texttt{[23:39--33:04]} excluded per project guidelines).
  \item \textbf{Analyst:} Baani Solanki (\texttt{2025114016})
  \item \textbf{Language:} Hindi / Hinglish (Native Hindi audio, Romanized script; non-AI-dubbed)
\end{itemize}

\subsection*{Event-specific questions (Event 2)}
\begin{enumerate}[leftmargin=1.6em, itemsep=4pt, topsep=4pt]
  \item \textbf{Are participants called "students", "candidates", "protesters", "agitators", "mobs", or "political pawns"?}
  \begin{itemize}[leftmargin=1.4em, itemsep=2pt, topsep=2pt]
    \item Sarthak consistently legitimizes participants as \texttt{"students"} (\texttt{[00:08]}), \texttt{"aam janta"} [ordinary public] (\texttt{[03:39]}), \texttt{"jenjis"} / Gen-Z youth (\texttt{[10:40]}, \texttt{[12:04]}), \texttt{"lakhon bachche"} [lakhs of children/kids] (\texttt{[22:48]}), and \texttt{"lidarles moovment"} whose leader became \texttt{"khud pablik"} [the public itself] (\texttt{[08:08]}). He explicitly quotes and mocks mainstream TV anchors who label them \texttt{"aatankvadi"} [terrorists] (\texttt{[05:31]}, \texttt{[22:22]}), \texttt{"goons oph da CJP"} (\texttt{[05:38]}), \texttt{"desh virodhi gaing"} [anti-national gang] (\texttt{[05:47]}), and \texttt{"pakistan phanded"} (\texttt{[22:48]}). He also honestly notes peripheral \texttt{"ansoshal eliments"} [anti-social elements] hooting near Connaught Place (\texttt{[07:12--07:21]}), separating them from the core Jantar Mantar site.
  \end{itemize}
  \item \textbf{Are their stated demands explained, summarised, mocked, questioned, or left out?}
  \begin{itemize}[leftmargin=1.4em, itemsep=2pt, topsep=2pt]
    \item \textbf{Clearly explained, summarised, and endorsed.} He opens with Education Minister Dharmendra Pradhan's resignation (\texttt{[00:08]}) and itemizes the remaining three accepted demands at \texttt{[08:16--08:35]}: (i) \texttt{Rs.\,1 crore} compensation to families of students who died by suicide due to paper leaks, (ii) a formal Delhi Police apology for the 20 July crackdown, and (iii) withdrawal/non-filing of criminal cases against protesters, alongside deeper systemic anti-leak reforms (\texttt{[16:25--17:07]}).
  \end{itemize}
  \item \textbf{Does the presenter focus on the original grievance, disruption, political affiliations, or public inconvenience?}
  \begin{itemize}[leftmargin=1.4em, itemsep=2pt, topsep=2pt]
    \item Focuses squarely on the \textbf{original systemic grievance} (repeated paper leaks, student suicides, collapsing education budget from \texttt{4.5\%} to \texttt{2.5\%} of GDP, \texttt{[15:59--18:10]}) and \textbf{democratic civic solidarity} (\texttt{[07:21--08:16]}), while debunking claims of foreign funding (\texttt{[00:40--00:48]}).
  \end{itemize}
  \item \textbf{Who is represented as responsible: examination authorities, government, students, political parties, or others?}
  \begin{itemize}[leftmargin=1.4em, itemsep=2pt, topsep=2pt]
    \item Primary responsibility is assigned to the \textbf{Union \& State Governments} (for reactive \texttt{"ditrens baay phiyar"} instead of preventive exam reforms and for slashing education spending in favor of electoral \texttt{"phri bis"} [freebies], \texttt{[16:34--18:53]}), \textbf{Delhi Police \& unbadged plainclothes men} for the 20 July crackdown (\texttt{[03:39--04:17]}), \textbf{Central Universities} (DU, JNU, Jamia, IITs) for issuing circulars threatening students with legal action (\texttt{[15:00--15:20]}), and \textbf{mainstream TV anchors} for vilifying youth (\texttt{[05:31--05:56]}, \texttt{[21:56--23:30]}).
  \end{itemize}
  \item \textbf{Are questions framed around whether students have legitimate grievances, or why they are causing disruption?}
  \begin{itemize}[leftmargin=1.4em, itemsep=2pt, topsep=2pt]
    \item Framed entirely around \textbf{students' legitimate grievances and state accountability}: why it took a massive street movement for a single resignation (\texttt{[01:51]}, \texttt{[15:59--16:16]}), who ordered unbadged men to assault female protesters (\texttt{[03:57--04:17]}), and why education budgets are shrinking (\texttt{[17:07--18:44]}).
  \end{itemize}
  \item \textbf{Are claims of violence, disorder, or political manipulation attributed to sources and substantiated?}
  \begin{itemize}[leftmargin=1.4em, itemsep=2pt, topsep=2pt]
    \item Claims of foreign/political conspiracy are refuted by citing the \textbf{Ministry of External Affairs (MEA)} clarification that there was \texttt{"koi bhi pakistan, chaina ya phoren phanding nahin"} (\texttt{[00:40--00:48]}). Claims of police violence (\texttt{lathi charj}, \texttt{aansoo gais}, \texttt{pelet gans}, plainclothes harassment at \texttt{[03:39--04:17]}) are backed by ground eyewitness observation (\texttt{[07:12--08:16]}) and cellphone videos (\texttt{[04:53]}), while fairly acknowledging that \texttt{"wahan pe kaee sare police vale bhi hart hue the. ritalieshan mein un pe bhi hamla huaa tha"} (\texttt{[04:45--04:53]}). He also transparently notes CJP founder Abhijit Dipke's prior AAP social-media stint (\texttt{[02:19]}).
  \end{itemize}
  \item \textbf{How are police and government actions described compared with protesters' actions? Are similar actions evaluated differently?}
  \begin{itemize}[leftmargin=1.4em, itemsep=2pt, topsep=2pt]
    \item Police actions on 18 July and 20 July are condemned in stark active verbs: \texttt{"abadakt karke le jati hai. not diten abadakt"} [abducts, not detains Sonam Wangchuk] (\texttt{[03:02--03:12]}), \texttt{"broot phors yooz karke"} (\texttt{[04:27]}), \texttt{"pelet gans ki bauchhar"} (\texttt{[03:39--03:47]}). Conversely, protester clashes are contextualized as \texttt{"ritalieshan"} [retaliation] (\texttt{[04:45]}), and post-20-July protesters are celebrated for peaceful mutual aid---forming human chains, medical camps, book libraries, and Zomato/Swiggy food sharing (\texttt{[07:21--08:08]}). Interestingly, at \texttt{[11:54--12:31]}, he humanizes young Gen-Z police jawans on the ground as sympathetic peers caught under senior orders.
  \end{itemize}
  \item \textbf{Are students allowed to explain their demands in their own words, or does the presenter mainly speak for them?}
  \begin{itemize}[leftmargin=1.4em, itemsep=2pt, topsep=2pt]
    \item The presenter \textbf{mainly speaks for and narrates on behalf of the students} in his studio monologue format, supplemented by reading out student protest posters (\texttt{"pavar of pipal ij gretar den da pipal in pavar"}, \texttt{[12:39]}) and recounting his own visit to Jantar Mantar (\texttt{[07:12--08:16]}), while inserting brief clips of politicians and rival TV anchors (\texttt{[05:38--05:56]}, \texttt{[22:30]}).
  \end{itemize}
  \item \textbf{Does the outlet's framing differ between coverage of the NEET paper-leak issue before the protest and coverage once students begin protesting?}
  \begin{itemize}[leftmargin=1.4em, itemsep=2pt, topsep=2pt]
    \item Rather than turning hostile once students hit the streets (as mainstream TV channels did), Sarthak argues that the street protest \textbf{elevated and validated} his prior anti-corruption stance on education---declaring \texttt{"bhaad mein gaee nyootraliti. agar tum ejukeshan karapt kar doge toh desh barbad hai"} (\texttt{[18:17--18:26]}). \textit{(Note: Direct pre-protest episodes are outside this single transcript, though he references his continuous week-long coverage at \texttt{[32:36]}.)}
  \end{itemize}
\end{enumerate}

\subsection*{Apply the five shared features (Event 2)}
\begin{longtable}{@{}>{\raggedright\arraybackslash}p{0.19\textwidth} >{\raggedright\arraybackslash}p{0.38\textwidth} >{\raggedright\arraybackslash}p{0.37\textwidth}@{}}
\toprule
\textbf{Feature} & \textbf{Exact wording + timestamp} & \textbf{What it does in context} \\
\midrule
\endhead
\textbf{1. Naming} & \texttt{"subah-subah hangar straik par baithe sonam vangachuk ko dilli police aati hai sivil kloths mein... wahan se abadakt karke le jati hai. not diten abadakt"} \texttt{[03:02--03:12]} & Explicitly rejects the legalistic police euphemism \texttt{"diten"} (detain) and replaces it with \texttt{"abadakt"} (abduct), criminalizing the plainclothes removal of hunger-striker Sonam Wangchuk and framing 20 July as \texttt{"blaik mande"} [Black Monday] (\texttt{[03:12]}). \\
\addlinespace[4pt]
\textbf{2. Presupposition} & \texttt{"aakhiri mein eego todna hi pada... pipal aar not skerd enimor."} \texttt{[00:40--00:58]} & \texttt{"Eego todna hi pada"} presupposes that the government's 30-day delay was driven by personal arrogance (\texttt{ego}) rather than policy deliberation; \texttt{"not skerd enimor"} presupposes an earlier climate of state-induced public fear. \\
\addlinespace[4pt]
\textbf{3. Agency} & \texttt{"dilli police aur sikyoriti phorses ne aam janta par lathi charj kar diya. aansoo gais chhod die. pelet gans ki bauchhar kar di."} \texttt{[03:39--03:47]} vs \texttt{"ladakiyon ko seksual hais karna... kuch logon dvara kiya gaya jo na vardi mein the na nem taig the."} \texttt{[03:57--04:07]} & Places direct active agency (\texttt{ne... kar diya}) on Delhi Police for tear gas and pellet guns, while highlighting how removing uniforms and name-tags deliberately obscures institutional agency (\texttt{kuch logon dvara}) to evade accountability. \\
\addlinespace[4pt]
\textbf{4. Antagonisation} & \texttt{"pavar of pipal ij gretar den da pipal in pavar."} \texttt{[12:39]} \newline  \texttt{"hamara media jisko aapki aavaj uthani thi usi ne aapki aavaj ka daman kiya."} \texttt{[21:56]} & Positions organic, civic-minded youth and ordinary taxpayers (\texttt{pablik}, \texttt{lakhon bachche}) against an arrogant ruling establishment and complicit \texttt{"godi media"} anchors (\texttt{"doosra Goswami"}, \texttt{[22:22]}) who smear students as \texttt{"terarist"} and \texttt{"enti neshanal"}. \\
\addlinespace[4pt]
\textbf{5. Intensification / softening} & \textbf{Intensification:} \texttt{"pelet gans ki bauchhar kar di"} \texttt{[03:39]}, \texttt{"jan sailab aane laga"} [a human flood arrived] \texttt{[03:12]}, \texttt{"enkar ko dalali karne se nahin roka"} [stop anchors from pimping/brokering] \texttt{[23:13]} \newline  \textbf{Softening (protester clashes):} \texttt{"ritalieshan mein un pe bhi hamla huaa tha"} \texttt{[04:45]} & Heightens state repression and media complicity via flood/barrage metaphors and harsh moral vocabulary (\texttt{dalali}), while softening protester violence against 118 injured police personnel by framing it as reactive \texttt{"ritalieshan"}. \\
\bottomrule
\end{longtable}

\subsection*{Short summary (Event 2)}
\begin{itemize}[leftmargin=1.5em, itemsep=3pt, topsep=4pt]
  \item \textbf{What interpretation of the event does this video encourage?} That the NEET/CJP protest was a historic, organic, non-partisan youth movement (\texttt{[00:40]}, \texttt{[15:39]}) triggered by systemic educational rot and paper leaks (\texttt{[15:59]}), which overcame brutal police repression on 20 July and mainstream media defamation to force Union Education Minister Dharmendra Pradhan's resignation (\texttt{[00:08--00:48]}, \texttt{[08:16--08:52]}).
  \item \textbf{Whose perspective is foregrounded, challenged, or missing?} \textit{Foregrounded:} Protesting Gen-Z students, middle/lower-class parents burdened by private school fees (\texttt{[19:09--19:35]}), and ground TV reporters caught between angry crowds and toxic studio anchors (\texttt{[22:07--22:22]}). \textit{Challenged:} Government delay, university gag circulars (\texttt{[15:00--15:20]}), and Republic-style "terrorist/anti-national" smears (\texttt{[05:31--05:56]}, \texttt{[22:22--23:04]}). \textit{Missing:} Direct first-person interviews with NTA/Ministry officials or injured police officers explaining crowd-control failures.
  \item \textbf{What evidence supports your interpretation? What complicates it?} \textit{Supports:} MEA's confirmation of zero foreign funding (\texttt{[00:40--00:48]}), specific GDP education budget trends (\texttt{4.5\%} down to \texttt{2.5\%} at \texttt{[17:52--18:10]}; \texttt{86 lakh} students leaving government schools at \texttt{[19:19]}), and firsthand observation of Jantar Mantar's medical/library/food camps (\texttt{[07:21--08:08]}). \textit{Complicates:} He acknowledges CJP founder Abhijit Dipke's past AAP social-media role (\texttt{[02:19]}), presence of hooting \texttt{"ansoshal eliments"} near CP (\texttt{[07:12--07:21]}), and injuries to 118 police personnel (\texttt{[04:37--04:53]}), showing the event was more complex than pure unprovoked one-way state action.
\end{itemize}
\vspace{4pt}\noindent\rule{\textwidth}{0.4pt}\vspace{4pt}

\section*{Common form: Blank excerpt record}
\textit{Copy or print this page whenever you find a passage worth analysing. Keep the exact wording and enough context to interpret it.}


\subsection*{Excerpt Record 1 (Event 1 --- Ram Mandir Donation Controversy)}
\begin{itemize}[leftmargin=1.5em, itemsep=3pt, topsep=4pt]
  \item \textbf{Outlet:} Sarthak Goswami (\textit{Sunday Show}) | \textbf{Date:} \texttt{2026-07-05}
  \item \textbf{Video title:} \textit{Ram Mandir Scam, Gadkari's Ethanol, Kejriwal's Temple \& More | Sunday Show}
  \item \textbf{Timestamp(s):} \texttt{[00:37--00:56]} | \textbf{Analyst:} Baani Solanki | \textbf{Language:} Hindi / Hinglish (Roman script)
  \item \textbf{Exact transcript excerpt:}
    \begin{quote}\small\ttfamily [00:37] apni report dayar ki jiske baad Trust ke aath logon ko arest kiya gaya. is poore ghotale ka mastarmaind bataya ja raha hai tinnoo e Ram shankar yadav... [00:46] champat raay RSS VHP ka aadmi hai aur tinnoo draivar bali ka bakra. [00:56] tinnoo par aarop hai ki usne apne malik champat raay ki pahunch ka phayda uthakar apne sare phaimili membars ko kaunting tim ka part bana liya.\end{quote}
  \item \textbf{Feature(s):} $\boxtimes$ Naming $\boxtimes$ Presupposition $\boxtimes$ Agency $\square$ Antagonisation $\boxtimes$ Intensification / softening
  \item \textbf{Who is speaking? Is this the outlet's wording or a quotation?} Sarthak Goswami speaking in his own editorial voice, summarizing official SIT claims (\texttt{bataya ja raha hai}) before interjecting his own judgement (\texttt{bali ka bakra}).
  \item \textbf{What does the wording imply, foreground, soften, or leave unstated?} The passive construction \texttt{"mastarmaind bataya ja raha hai"} [is being called the mastermind] distances Sarthak from the official police theory, while the copular assertion \texttt{"champat raay RSS VHP ka aadmi hai aur tinnoo draivar bali ka bakra"} presupposes political protection for Champat Rai and frames the driver as a sacrificial fall guy.
  \item \textbf{What context or evidence supports or complicates this interpretation?} Supported by the stark power asymmetry between a Rs.\,12,000--20,000/month driver and the VHP Vice-President running the Trust (\texttt{[01:14--01:40]}). Complicated by CCTV and asset-recovery reports showing counting staff personally amassed lakhs and crores (\texttt{[01:14--01:31]}), meaning Tinnu was an active participant, not an innocent bystander.
  \item \textbf{How does another outlet describe the same event or actor?} \textbf{Republic TV} (\texttt{0H-bruoEkaY}, \texttt{[03:22--03:42]} \& \texttt{[07:05--07:14]}) describes the exact same 8 arrested men strictly as \texttt{"low-level functionaries in the Ram Mandir trust, some cashiers"} to insulate senior leadership and praise \texttt{"Yogi's zero tolerance"}, whereas \textbf{Aaj Tak} panelist Sangeet Ragi (\texttt{B4L1pz0o4dc}, \texttt{[12:25--12:56]}) defends Champat Rai as an ascetic \texttt{"do dhoti aur do kurte par baithne vala aadmi"}.
  \item \textbf{My interpretation, with reasons:} Sarthak uses colloquial idiom (\texttt{bali ka bakra}) and patron-client naming (\texttt{malik champat raay}) to pierce bureaucratic scapegoating, encouraging viewers to question how a personal driver could pack the counting room with relatives without senior complicity.
\end{itemize}

\subsection*{Excerpt Record 2 (Event 1 --- Ram Mandir Donation Controversy)}
\begin{itemize}[leftmargin=1.5em, itemsep=3pt, topsep=4pt]
  \item \textbf{Outlet:} Sarthak Goswami (\textit{Sunday Show}) | \textbf{Date:} \texttt{2026-07-05}
  \item \textbf{Video title:} \textit{Ram Mandir Scam, Gadkari's Ethanol, Kejriwal's Temple \& More | Sunday Show}
  \item \textbf{Timestamp(s):} \texttt{[02:54--03:28]} | \textbf{Analyst:} Baani Solanki | \textbf{Language:} Hindi / Hinglish (Roman script)
  \item \textbf{Exact transcript excerpt:}
    \begin{quote}\small\ttfamily [02:54] dharma chalane ka bachane ka theka toh tum hi ne diya tha in logon ko. ab jab chori pakdi gaee toh tum hi ko kya slaid [gaslight] kar rahe hain ye? Hindu ka paisa agar kisi Hindu ne le bhi [03:02] liya bhaee tum logon ko kyon mirchi lag rahi hai? kahta hai hinduon ka paisa agar Hindu ne chori kar liya toh galat kya hai? ... [03:12] dekh lo yahi hai inke Hindu rashtra ka trelar... [03:20] main yahi log Hindu khatre mein hai. Hindu khatre mein tumhen ghutti pilae ja rahe the. tum pie ja rahe the.\end{quote}
  \item \textbf{Feature(s):} $\boxtimes$ Naming $\boxtimes$ Presupposition $\boxtimes$ Agency $\boxtimes$ Antagonisation $\boxtimes$ Intensification / softening
  \item \textbf{Who is speaking? Is this the outlet's wording or a quotation?} Sarthak Goswami's commentary embedding a paraphrased quotation (\texttt{kahta hai...}) from pro-Hindutva social media defenders.
  \item \textbf{What does the wording imply, foreground, soften, or leave unstated?} Foregrounds ideological betrayal: \texttt{"dharma... ka theka"} [contract/franchise over religion] implies religion was privatized to political brokers; \texttt{"Hindu rashtra ka trelar"} intensifies a temple theft into a dystopian preview of majoritarian impunity where co-religionist crime is excused.
  \item \textbf{What context or evidence supports or complicates this interpretation?} Supported by real viral rationalizations circulating online (also cited independently by Dhruv Rathee in \texttt{oeH0fOSJuX8} at \texttt{[14:34--14:42]}: \textit{"If a Hindu takes the money of a Hindu, why are you people feeling guilty"}). Complicated by the fact that many Hindu organizations and saints (e.g., Nirmohi Akhada cited on The Wire \texttt{7rUL8fGgRJA} at \texttt{[29:20]}) fiercely condemned the theft rather than excusing it.
  \item \textbf{How does another outlet describe the same event or actor?} \textbf{Republic TV} (\texttt{0H-bruoEkaY}, \texttt{[01:23--01:52]}) inverts the accusation entirely, claiming critics of the scam are \texttt{"people who say let's divide Hindus... people who did not want the Ram Mandir itself"}, whereas \textbf{Dhruv Rathee} (\texttt{oeH0fOSJuX8}, \texttt{[14:24--14:42]}) labels such apologists \texttt{"blind devotees who look at Narendra Modi above Lord Ram"}.
  \item \textbf{My interpretation, with reasons:} By quoting the apologist's own communal logic and extending it to murder (\texttt{"Hindu ka khoon Hindu ne bahaya"}, \texttt{[03:12]}), Sarthak exposes how majoritarian identity politics erodes basic legal and moral accountability.
\end{itemize}

\subsection*{Excerpt Record 3 (Event 1 --- Callback on Temple Appeasement)}
\begin{itemize}[leftmargin=1.5em, itemsep=3pt, topsep=4pt]
  \item \textbf{Outlet:} Sarthak Goswami (\textit{Sunday Show}) | \textbf{Date:} \texttt{2026-07-05}
  \item \textbf{Video title:} \textit{Ram Mandir Scam, Gadkari's Ethanol, Kejriwal's Temple \& More | Sunday Show}
  \item \textbf{Timestamp(s):} \texttt{[27:37--28:13]} | \textbf{Analyst:} Baani Solanki | \textbf{Language:} Hindi / Hinglish (Roman script)
  \item \textbf{Exact transcript excerpt:}
    \begin{quote}\small\ttfamily [27:37] disaid kar lo. dharma ki rajniti rokni hai ya karni hai. kyunki wo aap BJP se betar nahin kar paoge. aap sasti BJP banke rah jaoge... [27:55] BJP ke khilaph jate ho BJP ban jate hain. aise hi wo CJP vala hai. wo abhi risentli karne laga ki ham dilli police ne chhatrapati shivaji maharaj aur bhagat sinh ki kitab kaise phenki maphi mango... [28:04] tum wahan pe ejukeshan ke lie baithe ho rashtravad ke lie. BJP ka gem bajane aate hain, BJP ka hi gem khelne lagte hain.\end{quote}
  \item \textbf{Feature(s):} $\boxtimes$ Naming $\boxtimes$ Presupposition $\square$ Agency $\boxtimes$ Antagonisation $\boxtimes$ Intensification / softening
  \item \textbf{Who is speaking? Is this the outlet's wording or a quotation?} Sarthak Goswami's direct editorial commentary critiquing Arvind Kejriwal/AAP and briefly warning CJP protesters.
  \item \textbf{What does the wording imply, foreground, soften, or leave unstated?} Coining \texttt{"sasti BJP"} [cut-price/cheap BJP] presupposes BJP owns the copyright on religious nationalism (\texttt{BJP ka gem}) and implies that secular/issue-based challengers lose credibility whenever they adopt religious or hyper-nationalist symbolism.
  \item \textbf{What context or evidence supports or complicates this interpretation?} Supported by Punjab AAP's Rs.\,80-crore temple redevelopment and taxpayer-funded pilgrimage announcements (\texttt{[26:14--26:35]}). Complicates any simplistic reading of Sarthak as an uncritical opposition or CJP cheerleader---even 3 weeks before CJP's victory, he critiques CJP for drifting from \texttt{"ejukeshan"} into \texttt{"rashtravad"} symbolism (\texttt{[28:04]}).
  \item \textbf{How does another outlet describe the same event or actor?} \textbf{Republic TV} (\texttt{0H-bruoEkaY}, \texttt{[12:06--12:36]}) also attacks Kejriwal's temple politics as \texttt{"opportunistic"}, though from a pro-BJP angle demanding he \texttt{"atone for one's sins"} before Lord Ram.
  \item \textbf{My interpretation, with reasons:} This passage demonstrates ideological consistency: Sarthak critiques religious instrumentalization across party lines, strengthening his credibility when he defends student protesters in Event 2.
\end{itemize}

\subsection*{Excerpt Record 4 (Event 2 --- NEET Paper Leak Protest)}
\begin{itemize}[leftmargin=1.5em, itemsep=3pt, topsep=4pt]
  \item \textbf{Outlet:} Sarthak Goswami (\textit{Sunday Show}) | \textbf{Date:} \texttt{2026-07-26}
  \item \textbf{Video title:} \textit{Dharmendra Pradhan Resigned, CJP won, Modi on Instagram \& more | Sunday Show}
  \item \textbf{Timestamp(s):} \texttt{[03:02--04:17]} | \textbf{Analyst:} Baani Solanki | \textbf{Language:} Hindi / Hinglish (Roman script)
  \item \textbf{Exact transcript excerpt:}
    \begin{quote}\small\ttfamily [03:02] hangar straik par baithe sonam vangachuk ko dilli police aati hai sivil kloths mein aur unko saphed kapde mein lapet ke embulens mein daal ke wahan se abadakt karke le jati hai. not [03:12] diten abadakt... end den kem da blaik mande 20th julaee 2026... [03:39] dilli police aur sikyoriti phorses ne aam janta par lathi charj kar diya. aansoo gais chhod die. pelet gans ki bauchhar kar di... [03:57] ladakiyon ko seksual hais karna, bachchiyon ko marna yeh sab kuch logon dvara kiya gaya jo na vardi mein the na nem taig the. police ke the toh police aisa kyon karti hai? aur agar police ke nahin the toh wo wahan kya kar rahe the police ke saath?\end{quote}
  \item \textbf{Feature(s):} $\boxtimes$ Naming $\square$ Presupposition $\boxtimes$ Agency $\boxtimes$ Antagonisation $\boxtimes$ Intensification / softening
  \item \textbf{Who is speaking? Is this the outlet's wording or a quotation?} Sarthak Goswami's own narrative reconstruction of the 18 July and 20 July events in Delhi.
  \item \textbf{What does the wording imply, foreground, soften, or leave unstated?} Self-repairing \texttt{"not diten abadakt"} explicitly strips state action of lawful policing legitimacy. The dilemma posed at \texttt{[04:07]} (\texttt{police ke the toh... agar police ke nahin the toh...}) traps the administration in an inescapable accountability bind over unbadged men assaulting \texttt{"ladakiyon"} and \texttt{"bachchiyon"} [young girls].
  \item \textbf{What context or evidence supports or complicates this interpretation?} Corroborated across multiple independent testimonies---Dhruv Rathee (\texttt{l-6Pm-112io}, \texttt{[05:05--05:14]} \& \texttt{[11:37--12:25]}), Mohak Mangal's injured student interviewees (\texttt{v-RikSqb0Vw}, \texttt{[01:38]} \& \texttt{[04:38]}), and Naseeruddin Shah on The Wire (\texttt{pYaTVvHeGoM}, \texttt{[07:34--07:42]}). Complicated by Delhi Police's counter-claim (on Aaj Tak \texttt{JXLIs0WtJAM}, \texttt{[19:43--23:14]}) that officers faced heavy stone-pelting at barricades, which Sarthak himself notes at \texttt{[04:45--04:53]}.
  \item \textbf{How does another outlet describe the same event or actor?} \textbf{Republic TV} (\texttt{HTvb6E\_1Qdc}, \texttt{[01:37--01:47]} \& \texttt{[03:36--04:19]}) completely omits police pellet guns, plainclothes men, and Wangchuk's removal, instead portraying \texttt{"our police, our RAF"} solely as victims attacked by \texttt{"rabid, bloodthirsty mobs"} and \texttt{"goons of Abhijit Dipke"}.
  \item \textbf{My interpretation, with reasons:} The metalinguistic correction (\texttt{"not diten abadakt"}) is textbook critical discourse framing: Sarthak foregrounds how state vocabulary ("detention", "crowd control") sanitizes irregular force.
\end{itemize}

\subsection*{Excerpt Record 5 (Event 2 --- NEET Paper Leak Protest)}
\begin{itemize}[leftmargin=1.5em, itemsep=3pt, topsep=4pt]
  \item \textbf{Outlet:} Sarthak Goswami (\textit{Sunday Show}) | \textbf{Date:} \texttt{2026-07-26}
  \item \textbf{Video title:} \textit{Dharmendra Pradhan Resigned, CJP won, Modi on Instagram \& more | Sunday Show}
  \item \textbf{Timestamp(s):} \texttt{[07:21--08:16]} | \textbf{Analyst:} Baani Solanki | \textbf{Language:} Hindi / Hinglish (Roman script)
  \item \textbf{Exact transcript excerpt:}
    \begin{quote}\small\ttfamily [07:21] par jo apne aap mein protest sait thi. bahut hi achchhi kamyuniti spes ban gaee thi. ek hyooman chen logon ne banaee huee hai taki rasta alag rakha [07:30] ja sake. wahan pe koi log postar bana raha hai, koi khana bant raha hai... paas mein ek medikal kaimp laga huaa hai... [07:44] ek laibreri bani huee hai toh apne log kitaben donet kar rahe hain... [07:58] Zomato Swiggi vale khade hue hain khana le gae. dait ij hau sarvaiv kiya poora moovment. toh ek taim pe jo lidarles moovment tha na uska lidar khud pablik hi ban [08:08] gaee... jab koi nahin tha kisi ke lie toh sab the sabke lie.\end{quote}
  \item \textbf{Feature(s):} $\boxtimes$ Naming $\square$ Presupposition $\boxtimes$ Agency $\square$ Antagonisation $\boxtimes$ Intensification / softening
  \item \textbf{Who is speaking? Is this the outlet's wording or a quotation?} Sarthak Goswami speaking as a firsthand eyewitness who visited Jantar Mantar (\texttt{[08:08]}).
  \item \textbf{What does the wording imply, foreground, soften, or leave unstated?} Renames a contested protest zone a \texttt{"bahut hi achchhi kamyuniti spes"} [a wonderful community space] and distributes constructive civic agency across anonymous citizens (\texttt{koi postar bana raha hai}, \texttt{koi khana bant raha hai}), countering "violent mob" stereotypes.
  \item \textbf{What context or evidence supports or complicates this interpretation?} Supported by ground visuals of medical/food aid and peaceful sit-ins echoed by Mohak Mangal's ground interviewees (\texttt{v-RikSqb0Vw}, \texttt{[05:57]}, \texttt{[10:36--10:47]}). Complicated by Sarthak's own admission seconds earlier (\texttt{[07:12--07:21]}) that adjacent Connaught Place lanes had \texttt{"ansoshal eliments"} arriving in vans and hooting.
  \item \textbf{How does another outlet describe the same event or actor?} \textbf{Republic TV} (\texttt{HTvb6E\_1Qdc}, \texttt{[04:46--05:26]}) characterizes the exact same Jantar Mantar site as a scene of \texttt{"drunken violence"}, \texttt{"sexual perversion"}, and anti-national \texttt{"Nepal model"} anarchy where \texttt{"the students had evaporated"}.
  \item \textbf{My interpretation, with reasons:} By separating the core protest camp from peripheral troublemakers (\texttt{[07:12--07:21]}) and detailing mundane acts of care (libraries, medicine lists, Swiggy orders), Sarthak constructs democratic legitimacy from below.
\end{itemize}

\subsection*{Excerpt Record 6 (Event 2 --- Media Critique \& Anti-National Labelling)}
\begin{itemize}[leftmargin=1.5em, itemsep=3pt, topsep=4pt]
  \item \textbf{Outlet:} Sarthak Goswami (\textit{Sunday Show}) | \textbf{Date:} \texttt{2026-07-26}
  \item \textbf{Video title:} \textit{Dharmendra Pradhan Resigned, CJP won, Modi on Instagram \& more | Sunday Show}
  \item \textbf{Timestamp(s):} \texttt{[22:22--23:04]} | \textbf{Analyst:} Baani Solanki | \textbf{Language:} Hindi / Hinglish (Roman script)
  \item \textbf{Exact transcript excerpt:}
    \begin{quote}\small\ttfamily [22:22] aaj bhi wo doosra Goswami laga hi huaa hai. yeh enti neshanal hain. ye terarist hain... [22:30] da kangrechuleshans CJP nau ophishiyli terarist sar kaming aaut. itne ghatiya, itne daphar, itne dam aargument men strim ho gae na... [22:48] lakhon lakhon bachche sadkon pe dikh rahe hain. sab entineshanal hain. sab pakistan phanded hain. ab stoodiyo mein jane se pahle akhbar padh le. pakistan khud \$10 biliyan dolar ka [22:56] lon mang raha hai amerika se. wo kya phand karega?\end{quote}
  \item \textbf{Feature(s):} $\boxtimes$ Naming $\boxtimes$ Presupposition $\square$ Agency $\boxtimes$ Antagonisation $\boxtimes$ Intensification / softening
  \item \textbf{Who is speaking? Is this the outlet's wording or a quotation?} Sarthak Goswami's wording parodying and directly responding to Republic TV anchor Arnab Goswami (\texttt{"doosra Goswami"}).
  \item \textbf{What does the wording imply, foreground, soften, or leave unstated?} Foregrounds how state-aligned TV channels weaponize national-security labels (\texttt{enti neshanal}, \texttt{terarist}, \texttt{pakistan phanded}) to delegitimize domestic youth grievances, and punctures the conspiracy via geopolitical satire (Pakistan begging for a \$10B bailout).
  \item \textbf{What context or evidence supports or complicates this interpretation?} Directly supported by MEA's official denial of foreign funding (\texttt{[00:40--00:48]}) and by Republic TV's actual broadcast (\texttt{HTvb6E\_1Qdc}, \texttt{[00:23--00:33]} calling Dipke \texttt{"Deep State Dipke"} and \texttt{[04:19--04:46]} contrasting protesters with soldiers fighting terrorists in Anantnag).
  \item \textbf{How does another outlet describe the same event or actor?} \textbf{Aaj Tak} (\texttt{JXLIs0WtJAM}, \texttt{[10:04--11:10]}) avoids calling students terrorists, instead airing Maharashtra CM Devendra Fadnavis's softer delegitimization that outsiders \texttt{"apna-apna ejenda chala rahe hain"} without even knowing what NEET is.
  \item \textbf{My interpretation, with reasons:} This meta-discursive critique directly exemplifies the Language \& Society theme: Sarthak exposes \textit{Naming} ("terrorist" vs \texttt{"lakhon bachche"}) as a deliberate ideological tool of media delegitimization.
\end{itemize}
\vspace{4pt}\noindent\rule{\textwidth}{0.4pt}\vspace{4pt}

\section*{Common form: End-of-event summary}

\subsection*{End-of-Event Summary --- Event 1 (Ram Mandir Donation Controversy)}
\begin{itemize}[leftmargin=1.5em, itemsep=3pt, topsep=4pt]
  \item \textbf{The main way this outlet framed the event was:} Not merely as an Rs.\,80-lakh/crore accounting theft by temple staff, but as a deeper moral and political \texttt{"asli ghotala"} [real scam] resulting from the commercialization of Hindu faith (\texttt{"dharma ka dhandha"}) and citizens outsourcing religious authority to unaccountable political-ideological contractors (\texttt{[00:00--00:56]}, \texttt{[02:12--06:26]}).
  \item \textbf{The strongest linguistic evidence was (quote + timestamp):} \texttt{"champat raay RSS VHP ka aadmi hai aur tinnoo draivar bali ka bakra... Hindu jago Ram ka chanda nahin Ram ka charitra chori ho chuka hai."} (\texttt{[00:46]} \& \texttt{[06:26]}).
  \item \textbf{The shared feature(s) most clearly illustrated were:} \textbf{Naming} (\texttt{bali ka bakra}, \texttt{dharma ka dhandha}, \texttt{charitra chori}), \textbf{Presupposition} (\texttt{tum sarapraizd kyon ho?} at \texttt{[02:12]}), and \textbf{Agency} (holding voters/devotees active agents who handed over the \texttt{"theka"} of religion at \texttt{[02:46--03:38]}).
  \item \textbf{An alternative interpretation of the same evidence is:} That the embezzlement was an operational lapse in cash-counting protocols by colluding counting staff and bank supervisors---not an indictment of all temple devotees or of the Ram Mandir movement itself---and that the UP government's SIT arrest of 8 suspects (\texttt{[00:37]}) represented lawful institutional action rather than a cover-up.
  \item \textbf{What I cannot conclude from this transcript alone:} Whether senior trustees like Champat Rai personally profited from or directed Tinnu Yadav's theft, what the final SIT forensic audit concluded, or the exact rupee total stolen across cash and jewellery.
\end{itemize}

\subsection*{End-of-Event Summary --- Event 2 (Students' Protest Regarding NEET Paper Leak)}
\begin{itemize}[leftmargin=1.5em, itemsep=3pt, topsep=4pt]
  \item \textbf{The main way this outlet framed the event was:} As a victorious, organic, citizen-led democratic awakening (\texttt{"lidar khud pablik hi ban gaee"}, \texttt{[08:08]}) rooted in genuine educational despair (paper leaks, suicides, falling GDP share for education) that defeated police brutality (\texttt{20 July Black Monday}) and \texttt{"godi media"} smear campaigns to compel ministerial accountability (\texttt{[00:08--01:26]}, \texttt{[15:39--18:44]}).
  \item \textbf{The strongest linguistic evidence was (quote + timestamp):} \texttt{"sonam vangachuk ko dilli police aati hai sivil kloths mein... wahan se abadakt karke le jati hai. not diten abadakt... jab koi nahin tha kisi ke lie toh sab the sabke lie."} (\texttt{[03:02--03:12]} \& \texttt{[08:08--08:16]}).
  \item \textbf{The shared feature(s) most clearly illustrated were:} \textbf{Naming} (\texttt{not diten abadakt}, \texttt{blaik mande}, \texttt{kamyuniti spes} vs TV anchors' \texttt{terarist}/\texttt{enti neshanal}), \textbf{Agency} (foregrounding police lathi/pellet-gun force and unbadged plainclothes men at \texttt{[03:39--04:17]}), and \textbf{Antagonisation} (\texttt{pavar of pipal} vs \texttt{pipal in pavar} and \texttt{"dalali"} anchors at \texttt{[12:39]}, \texttt{[23:13]}).
  \item \textbf{An alternative interpretation of the same evidence is:} That once an unauthorized crowd breached Section 163 barricades to march toward Parliament on 20 July and injured 118 police personnel (\texttt{[03:30--03:39]}, \texttt{[04:37--04:53]}), security forces faced an acute law-and-order breakdown, and the government's eventual acceptance of demands reflected negotiated democratic responsiveness rather than mere bruised \texttt{"eego"} (\texttt{[00:40]}).
  \item \textbf{What I cannot conclude from this transcript alone:} The exact identity/command chain of the plainclothes individuals seen without name-tags (\texttt{[04:07]}), how many injuries on 20 July were caused by initial police action versus crowd stone-pelting, or whether post-resignation structural NTA reforms actually prevented future leaks.
\end{itemize}
\textit{Research reminder: interpret wording in context, distinguish allegations from established facts, and support claims with exact excerpts. Mixed or contradictory evidence is worth recording.}


\clearpage

\begin{center}
    {\large\bfseries LANGUAGE \& SOCIETY PROJECT}\\[3pt]
    {\LARGE\bfseries Transcript Analysis Worksheet}\\[3pt]
    {\normalsize Qualitative analysis of political news discourse}
\end{center}
\vspace{2pt}
\noindent\rule{\textwidth}{0.4pt}
\vspace{4pt}


\section*{Creator 2: Dhruv Rathee (Independent Investigative / Political Explainer YouTuber)}
\begin{itemize}[leftmargin=1.5em, itemsep=3pt, topsep=4pt]
  \item \textbf{Analyst:} Baani Solanki (\texttt{2025114016})
  \item \textbf{Channel / Outlet:} Dhruv Rathee (\texttt{@dhruvrathee}, YouTube)
  \item \textbf{Ideological / Editorial Positioning:} Anti-establishment, progressive-liberal investigative explainer; frames both controversies as systemic regime-level corruption and authoritarian state repression (\texttt{"Tana Shah government"}, \texttt{"dhritarashtra"}) directly implicating Prime Minister Narendra Modi, Home Minister Amit Shah, and the BJP-RSS ecosystem.
  \item \textbf{Transcript Files Referenced:}
  \begin{itemize}[leftmargin=1.4em, itemsep=2pt, topsep=2pt]
    \item Event 1: \texttt{transcripts/event1\_ram\_mandir/02\_dhruv\_rathee\_oeH0fOSJuX8.txt}
    \item Event 2: \texttt{transcripts/event2\_neet\_protest/02\_dhruv\_rathee\_l-6Pm-112io.txt}
    \item \textit{(Caption-track note: Dhruv Rathee's team manually uploaded English-translated Roman-script subtitles preserving key Hindi lexical items such as \texttt{Chanda Chori}, \texttt{Ramlal}, \texttt{dhritarashtra}, \texttt{Anshan}, and \texttt{Tana Shah} directly into YouTube's manual \texttt{hi} native slot; non-AI-dubbed original audio verified.)}
  \end{itemize}
\end{itemize}
\vspace{4pt}\noindent\rule{\textwidth}{0.4pt}\vspace{4pt}

\section*{Guidelines: Five features to look for in every event}
\begin{longtable}{@{}>{\raggedright\arraybackslash}p{0.17\textwidth} >{\raggedright\arraybackslash}p{0.20\textwidth} >{\raggedright\arraybackslash}p{0.24\textwidth} >{\raggedright\arraybackslash}p{0.32\textwidth}@{}}
\toprule
\textbf{Feature} & \textbf{What to look for} & \textbf{Illustrative example} & \textbf{Questions to answer} \\
\midrule
\endhead
\textbf{1. Naming} & What are people, groups, actions, or events called? & "Students", "protesters", "agitators"; "irregularity", "scam", "controversy". & What exact label is used? Who or what does it refer to? Who uses it? What alternative label could have been used? \\
\addlinespace[4pt]
\textbf{2. Presupposition} & What is treated as already true or taken for granted? & "Why are students creating chaos again?" may presuppose earlier disruption. & What does the wording assume? Is that assumption supported, disputed, or left unexamined? Does it appear in a question or words such as "again", "still", or "finally"? \\
\addlinespace[4pt]
\textbf{3. Agency} & Who is shown doing something? Who is acted upon? Who disappears? & "Police injured protesters" vs. "Protesters were injured." & Who is the actor? Who experiences the action? Is responsibility named, backgrounded, or omitted? \\
\addlinespace[4pt]
\textbf{4. Antagonisation} & Who is positioned against whom? & "Ordinary citizens" vs. "political troublemakers". & Which groups are contrasted? Is one presented as more legitimate, moral, patriotic, dangerous, or disloyal? Is an "us vs. them" division implied? \\
\addlinespace[4pt]
\textbf{5. Intensification / softening} & What is exaggerated, minimised, softened, or euphemised? & "Violent mob" vs. "peaceful demonstrators"; "incident" instead of a direct description. & Which words heighten emotion or seriousness? Which minimise it? Are metaphors, dramatic adjectives, hedges, or euphemisms used? Does the wording match the evidence? \\
\bottomrule
\end{longtable}
\textit{Reminder: A label or implication is not automatically false or manipulative. Examine the surrounding context and evidence.}

\vspace{4pt}\noindent\rule{\textwidth}{0.4pt}\vspace{4pt}

\section*{Event1. Ram Mandir donation controversy}
\textbf{Focus:} How is the controversy framed---as financial accountability, religious sentiment, political motivation, or something else? Be precise about the specific allegation or episode covered.

\begin{itemize}[leftmargin=1.5em, itemsep=3pt, topsep=4pt]
  \item \textbf{Outlet:} Dhruv Rathee (YouTube)
  \item \textbf{Date:} \texttt{2026-07-04}
  \item \textbf{Video title:} \textit{Ram Mandir Scam is WORSE Than You Think | Dhruv Rathee} (\url{https://youtu.be/oeH0fOSJuX8?si=tGAGQeFY_MRU32Az})
  \item \textbf{Timestamp(s):} Relevant Event 1 segments \texttt{[00:00--01:54]} and \texttt{[02:49--15:58]} (total video duration \texttt{16:25}; self-promotional AI Masterclass / Bootcamp segments at \texttt{[01:55--02:48]} and \texttt{[15:59--16:25]} excluded per project guidelines).
  \item \textbf{Analyst:} Baani Solanki (\texttt{2025114016})
  \item \textbf{Language:} Hindi / Hinglish (Spoken native Hindi; creator-uploaded Roman-script manual \texttt{hi} track, non-AI-dubbed)
\end{itemize}

\subsection*{Event-specific questions (Event 1)}
\begin{enumerate}[leftmargin=1.6em, itemsep=4pt, topsep=4pt]
  \item \textbf{What terms are used: "scam", "alleged irregularity", "corruption", "political controversy", or something else?}
  \begin{itemize}[leftmargin=1.4em, itemsep=2pt, topsep=2pt]
    \item Rathee uses unmitigated, high-severity crime labels throughout: \texttt{"Scam"} (in title and at \texttt{[01:28]}, \texttt{[08:40]}, \texttt{[11:07]}, \texttt{[13:14]}), \texttt{"Chanda Chori"} / \texttt{"Chanda robbery"} (\texttt{[00:00]}, \texttt{[05:55]}, \texttt{[09:36]}, \texttt{[10:32]}), \texttt{"open-air robbery in the Ram Temple for years"} (\texttt{[09:36]}), \texttt{"dacoits"} (\texttt{[00:00]}, \texttt{[10:32]}), \texttt{"big frauds"} (\texttt{[10:21]}, \texttt{[11:48]}), and \texttt{"looted"} (\texttt{[01:44]}, \texttt{[15:06]}).
  \end{itemize}
  \item \textbf{How are the temple, trust, donors, and people making allegations described?}
  \begin{itemize}[leftmargin=1.4em, itemsep=2pt, topsep=2pt]
    \item \textbf{Trust \& Leadership:} Explicitly stripped of institutional autonomy---\texttt{"this is not an independent Trust. The government is directly involved in this"} (\texttt{[04:46]})---noting that 9 of 15 members have direct Modi/BJP/RSS links, including PM's ex-Principal Secretary Nripendra Misra, CEC Gyanesh Kumar, Anil Mishra, and VHP Vice-President Champat Rai (\texttt{[04:18--05:02]}). Trustees are called \texttt{"temple-owners"} who blocked SBI's warnings (\texttt{[07:33]}), while arrested counting agents (\texttt{Tinnu Yadav}, \texttt{Subhash Srivastav}, \texttt{Lavkush Mishra}, \texttt{K.D. Tiwari}, \texttt{Anukalp Mishra}) are portrayed as low-salaried (\texttt{Rs.\,14,000--20,000/month}) insiders who \texttt{"became billionaires"} (\texttt{Rs.\,50 crore} fortune attributed to Tinnu Yadav at \texttt{[07:33--07:54]}).
    \item \textbf{Donors:} Depicted empathetically as sincere believers whose \texttt{"faith was played with in the name of Lord Ram"} (\texttt{[01:44]})---exemplified by Pune's Sindhi community donating \texttt{210} 1-kg silver bricks worth \texttt{Rs.\,2 crore} (\texttt{[00:17--01:20]})---except for \texttt{"blind devotees who look at Narendra Modi above Lord Ram"} (\texttt{[14:24]}).
    \item \textbf{People Making Allegations:} Whistleblower account-in-charge Mahipal Singh (\texttt{[06:55]}, \texttt{[08:03--08:31]}), SP chief Akhilesh Yadav (\texttt{[08:50]}), Shiv Sena's Sanjay Raut (\texttt{[09:18]}), Mumbai businessman Anil Vishwakarma (\texttt{[09:27]}), BJP politician Rajni Singh (\texttt{[13:14--13:24]}), and AAP's Arvind Kejriwal (\texttt{[13:42]}) are all treated as credible corroborating witnesses.
  \end{itemize}
  \item \textbf{Is the issue framed as financial accountability, an attack on Hindu faith, a political dispute, or a matter of public trust?}
  \begin{itemize}[leftmargin=1.4em, itemsep=2pt, topsep=2pt]
    \item Framed as \textbf{systemic, state-protected financial corruption and betrayal of public/devotee trust} spanning every stage---\texttt{"From the purchase of land, to the construction and donation of the Ram Mandir, there have been frauds"} (\texttt{[11:56]})---enabled by deliberate executive blindness (\texttt{[12:03]}, \texttt{[13:54]}).
  \end{itemize}
  \item \textbf{Are questions about donations presented as legitimate scrutiny or as an attempt to malign the temple or its supporters?}
  \begin{itemize}[leftmargin=1.4em, itemsep=2pt, topsep=2pt]
    \item Presented as \textbf{vital public-interest scrutiny}. When UP CM Yogi Adityanath warns critics to \texttt{"Stop playing with their faith"} (\texttt{[14:13]}), Rathee immediately counters with a rhetorical question: \texttt{"Who is playing with their faith? The scammer or the one who is questioning the scam?"} (\texttt{[14:24]}).
  \end{itemize}
  \item \textbf{Are allegations attributed to specific people and supported with evidence, or presented as established facts?}
  \begin{itemize}[leftmargin=1.4em, itemsep=2pt, topsep=2pt]
    \item Rathee weaves \textbf{named documentary/witness attributions} into an \textbf{established-fact narrative}: he cites whistleblower Mahipal Singh's testimony (\texttt{[06:55]}, \texttt{[08:03--08:40]}), SBI's warning 3 months prior (\texttt{[07:15]}), a November 2020 private audit report (\texttt{[10:12--10:21]}), \textit{Dainik Bhaskar} investigations into \texttt{40\%} construction commissions (Somesh Anand) and \texttt{17x} land-price inflation (\texttt{[10:49--11:24]}), June 2021 land records (\texttt{[11:38]}), and BJP leader Rajni Singh's PMO complaint (\texttt{[13:14--13:42]}). Yet unverified upper-bound figures (\texttt{"donations ranging from Rs 200 crore"}, \texttt{[01:28]}; Tinnu Yadav's \texttt{"fortune of around Rs.\,50 crore"}, \texttt{[07:54]}; Kejriwal's claim of \texttt{"12 IB reports"}, \texttt{[13:42]}) are absorbed smoothly into his indictment.
  \end{itemize}
  \item \textbf{How are denials and counterarguments represented: explained, dismissed, challenged, or omitted?}
  \begin{itemize}[leftmargin=1.4em, itemsep=2pt, topsep=2pt]
    \item \textbf{Quoted briefly only to be dismantled.} Champat Rai's post-June-7 denial (\texttt{"there was no problem in the donation audit"}, \texttt{[08:58]}) is refuted with missing gold/silver receipts (\texttt{[09:08--10:03]}); Champat Rai and Anil Mishra's resignations on 27 June are dismissed as insufficient (\texttt{"Is it enough to resign? If there's a scam here, should they also be arrested?"}, \texttt{[13:05--13:14]}); Yogi Adityanath's counter-attack is inverted (\texttt{[14:13--14:24]}).
  \end{itemize}
  \item \textbf{Does the presenter distinguish criticism of the trust or its financial practices from criticism of the temple or religion itself?}
  \begin{itemize}[leftmargin=1.4em, itemsep=2pt, topsep=2pt]
    \item \textbf{Yes, philosophically, though with a radical anti-donation conclusion.} Rathee invokes the \textit{Swades} poem (\texttt{"Ram is in unity, in progress... Whoever takes out Ravan from the heart, Ram is in his heart"}, \texttt{[15:23--15:40]}) to separate ethical devotion to Lord Ram from \texttt{"the one who shouts Jai Shri Ram the loudest"} while harbouring \texttt{"Ravan"} (\texttt{[15:49]}). However, unlike all other 5 creators, he advises viewers: \texttt{"Whenever you go to a temple, any temple, never donate money... give it to a poor person. Give it to an NGO"} (\texttt{[14:54--15:15]}).
  \end{itemize}
  \item \textbf{Who is portrayed as protecting donors' interests, and who as exploiting religious sentiment?}
  \begin{itemize}[leftmargin=1.4em, itemsep=2pt, topsep=2pt]
    \item \textbf{Exploiting religious sentiment:} Trust General Secretary Champat Rai, his associates/nephews (Tinnu Yadav, Chandan Rai, Somesh Anand), UP CM Yogi Adityanath, and PM Modi/PMO (\texttt{"When you ask for votes in the name of Lord Ram, they stand in the front line first. But when the Lord Ram is being robbed, they let it happen"}, \texttt{[12:03--12:11]}).
    \item \textbf{Protecting donors' interests:} Fired whistleblower Mahipal Singh (\texttt{[08:03--08:31]}), investigative journalists (\textit{Dainik Bhaskar}), dissenting saints like Mahant Kamal Nayan Das (\texttt{[12:30--12:38]}), and opposition/citizen whistleblowers raising pressure (\texttt{[08:50]}, \texttt{[12:11]}, \texttt{[13:42]}).
  \end{itemize}
\end{enumerate}

\subsection*{Apply the five shared features (Event 1)}
\begin{longtable}{@{}>{\raggedright\arraybackslash}p{0.19\textwidth} >{\raggedright\arraybackslash}p{0.38\textwidth} >{\raggedright\arraybackslash}p{0.37\textwidth}@{}}
\toprule
\textbf{Feature} & \textbf{Exact wording + timestamp} & \textbf{What it does in context} \\
\midrule
\endhead
\textbf{1. Naming} & \texttt{"this is not an independent Trust. The government is directly involved... The bank was saying to remove them, but the temple-owners were saying, let them be."} \texttt{[04:46]}, \texttt{[07:33]} \newline  \texttt{"Kejriwal called Modi a dhritarashtra."} \texttt{[13:54]} & Reclassifies fiduciary trustees into proprietary \texttt{"temple-owners"} and fuses the Trust with the Modi government; amplifies the Mahabharata metaphor \texttt{"dhritarashtra"} (the willfully blind king protecting corrupt kin) to frame executive silence as complicity. \\
\addlinespace[4pt]
\textbf{2. Presupposition} & \texttt{"Ram Mandir Scam is WORSE Than You Think"} \texttt{[Title]} \newline  \texttt{"Let's understand how in today's video, people's faith was played with in the name of Lord Ram, and the public was looted."} \texttt{[01:44]} & Both the comparative adjective (\texttt{WORSE Than You Think}) and the \textit{how}-complement (\texttt{understand how... people's faith was played with}) presuppose the existence of a massive scam as indisputable common ground before the main explanation even begins. \\
\addlinespace[4pt]
\textbf{3. Agency} & \texttt{"On the contrary, Mahipal Singh was removed from his job. Apart from that, CCTV footage of 7-8 months was also deleted."} \texttt{[08:22]} vs \texttt{"And the government was blind to this. Was it intentional or was it just like that?"} \texttt{[12:03]} & Uses strategic passives (\texttt{was removed}, \texttt{was deleted}) where insiders covered tracks, then immediately pivots to pointed rhetorical questions forcing the viewer to assign intentional agency up the chain to Champat Rai and the government (\texttt{[08:03]}, \texttt{[12:03]}). \\
\addlinespace[4pt]
\textbf{4. Antagonisation} & \texttt{"the one who shouts Jai Shri Ram the loudest, is not Ram in their hearts, but Ravan. The biggest criminals who rob in the name of Ram, and even more, those who are saving them."} \texttt{[15:49--15:58]} & Constructs a stark mythological-moral binary (\texttt{Ram} vs \texttt{Ravan}) casting vocal Hindutva politicians and Trust protectors not merely as negligent administrators, but as \texttt{"Ravan"} and \texttt{"the biggest criminals"}. \\
\addlinespace[4pt]
\textbf{5. Intensification / softening} & \textbf{Intensification:} \texttt{"Six months passed. Then a year passed. Two years. Five years... everything was disappeared"} \texttt{[01:03--01:28]}; \texttt{"open-air robbery in the Ram Temple for years"} \texttt{[09:36]}; \texttt{"have you ever heard of them happening in one project?"} \texttt{[11:56]} & Narrative pacing (counting elapsed years in short staccato sentences), hyperbole (\texttt{everything was disappeared}, \texttt{became billionaires}), and superlative rhetorical questions intensify moral outrage across the entire 6-year timeline. \\
\bottomrule
\end{longtable}

\subsection*{Short summary (Event 1)}
\begin{itemize}[leftmargin=1.5em, itemsep=3pt, topsep=4pt]
  \item \textbf{What interpretation of the event does this video encourage?} That the June 2026 cash-counting arrests are merely the tip of a six-year, multi-layered institutional plunder---encompassing \texttt{210 kg} missing silver bricks (\texttt{[00:26]}), unrecorded gold/silver jewellery (\texttt{[09:08--10:03]}), \texttt{40\%} construction kickbacks (\texttt{[10:49]}), and \texttt{17x} land-purchase inflation (\texttt{[11:15]})---shielded by BJP-RSS appointees and a willfully blind PMO (\texttt{[04:18--05:22]}, \texttt{[13:14--14:04]}).
  \item \textbf{Whose perspective is foregrounded, challenged, or missing?} \textit{Foregrounded:} Unreceipted donors (Pune Sindhi community), fired whistleblower Mahipal Singh, \textit{Dainik Bhaskar} investigative reporters, and opposition critics. \textit{Challenged:} Champat Rai's denials, the UP SIT's restriction of the FIR to 8 junior staff (\texttt{[12:48--13:05]}), and CM Yogi Adityanath's faith-based deflection (\texttt{[14:13--14:24]}). \textit{Missing:} Any direct interview or detailed rebuttal from Trust members, SBI officials, or SIT investigators explaining why senior trustees were not named in the initial FIR.
  \item \textbf{What evidence supports your interpretation? What complicates it?} \textit{Supports:} Specific donor testimonies (\texttt{210} silver bricks in Jan 2021 at \texttt{[00:17--01:20]}), November 2020 private audit warnings (\texttt{[10:12--10:21]}), Mahipal Singh's on-record account of bundle manipulation (\texttt{12--13} notes hidden in 10-note bundles at \texttt{[06:47--07:06]}) and deleted CCTV footage (\texttt{[08:22--08:40]}), and documented 2021 land-price jumps (\texttt{[11:15--11:38]}). \textit{Complicates:} Several headline numbers (\texttt{Rs.\,200 crore} total disappearance at \texttt{[01:28]}, \texttt{Rs.\,50 crore} for Tinnu Yadav at \texttt{[07:54]}, \texttt{12 IB reports} claimed by Kejriwal at \texttt{[13:42]}) rely on unverified media/political claims, and his sweeping conclusion (\texttt{"never donate money"} to \textit{any} temple, \texttt{[14:54]}) generalizes one Trust's failure to all religious institutions.
\end{itemize}
\vspace{4pt}\noindent\rule{\textwidth}{0.4pt}\vspace{4pt}

\section*{Event2. Students' protest regarding NEET paper leak}
\textbf{Focus:} How does coverage construct the legitimacy of student protest and responsibility for the underlying problem?

\begin{itemize}[leftmargin=1.5em, itemsep=3pt, topsep=4pt]
  \item \textbf{Outlet:} Dhruv Rathee (YouTube)
  \item \textbf{Date:} \texttt{2026-07-24} \textit{(published one day BEFORE Dharmendra Pradhan's resignation on 25 July)}
  \item \textbf{Video title:} \textit{How CJP defeated BJP at Jantar Mantar Protests? | Full Story of 20th July | Dhruv Rathee} (\url{https://youtu.be/l-6Pm-112io?si=TXnmzWqiJKF-qXev})
  \item \textbf{Timestamp(s):} Relevant Event 2 segments \texttt{[00:00--02:39]} and \texttt{[03:01--17:19]} (total video duration \texttt{17:50}; AI Masterclass promo \texttt{[02:40--03:00]} and end-screen promo \texttt{[17:20--17:50]} excluded per project guidelines).
  \item \textbf{Analyst:} Baani Solanki (\texttt{2025114016})
  \item \textbf{Language:} Hindi / Hinglish (Spoken native Hindi; creator-uploaded Roman-script manual \texttt{hi} track, non-AI-dubbed)
\end{itemize}

\subsection*{Event-specific questions (Event 2)}
\begin{enumerate}[leftmargin=1.6em, itemsep=4pt, topsep=4pt]
  \item \textbf{Are participants called "students", "candidates", "protesters", "agitators", "mobs", or "political pawns"?}
  \begin{itemize}[leftmargin=1.4em, itemsep=2pt, topsep=2pt]
    \item Rathee consistently names them \texttt{"protesters"} (\texttt{[00:23]}, \texttt{[01:21]}, \texttt{[08:20]}), \texttt{"the smart people"} (\texttt{[01:21]}), \texttt{"the youth of the country, the people of the country"} (\texttt{[02:55--03:05]}), \texttt{"thousands of students"} (\texttt{[06:33]}, \texttt{[15:41]}), \texttt{"normal people"} (\texttt{[09:26]}), and \texttt{"children"} / \texttt{"its own children"} (\texttt{[16:21--16:44]}), elevating their ethos to \texttt{"the revolutionary spirit, which was seen in the freedom fighters during the independence"} (\texttt{[15:06--15:14]}).
  \end{itemize}
  \item \textbf{Are their stated demands explained, summarised, mocked, questioned, or left out?}
  \begin{itemize}[leftmargin=1.4em, itemsep=2pt, topsep=2pt]
    \item \textbf{Summarised empathetically as minimal and self-evidently just:} \texttt{"the only thing they demanded was the resignation of Education Minister Dharmendra Pradhan"} (\texttt{[03:13]}) and \texttt{"There was only one demand from the children. That the exams should not be stolen. That a better Education Minister should be appointed"} (\texttt{[16:29--16:37]}).
  \end{itemize}
  \item \textbf{Does the presenter focus on the original grievance, disruption, political affiliations, or public inconvenience?}
  \begin{itemize}[leftmargin=1.4em, itemsep=2pt, topsep=2pt]
    \item Focuses overwhelmingly on a \textbf{minute-by-minute forensic reconstruction of pre-meditated state repression on 17--20 July} (\texttt{[00:23--02:19]}, \texttt{[04:27--14:54]}) and nationwide student solidarity (\texttt{[15:14--16:10]}), completely rejecting disruption/inconvenience frames.
  \end{itemize}
  \item \textbf{Who is represented as responsible: examination authorities, government, students, political parties, or others?}
  \begin{itemize}[leftmargin=1.4em, itemsep=2pt, topsep=2pt]
    \item Responsibility is placed squarely on \textbf{PM Narendra Modi} (\texttt{"scared Tana Shah"}, \texttt{[10:28]}; \texttt{"Shame on you, Narendra Modi"}, \texttt{[16:44]}), \textbf{Home Minister Amit Shah} (for abruptly replacing Delhi Police Commissioner Satish Golcha with IB Special Director Anurag Kumar on 17 July at \texttt{[04:27--04:54]}), \textbf{Health Minister J.P. Nadda} (accused of staging a 2-hour distraction meeting to decapitate protest leadership at \texttt{[14:20--14:47]}), and \textbf{Delhi Police / RAF / civil-clothes goons} (\texttt{[05:05]}, \texttt{[11:37--13:25]}).
  \end{itemize}
  \item \textbf{Are questions framed around whether students have legitimate grievances, or why they are causing disruption?}
  \begin{itemize}[leftmargin=1.4em, itemsep=2pt, topsep=2pt]
    \item Framed entirely around \textbf{state conspiracy and abuse of peaceful citizens}: why broken vehicles and brick-laden trucks appeared inside police cordons at \texttt{04:40 AM} (\texttt{[00:23--01:29]}), why metal detectors were removed (\texttt{[01:29--02:01]}), why Commissioner Satish Golcha was sacked (\texttt{[04:27--05:37]}), and why police hid their nameplates and faces (\texttt{[11:37--12:25]}).
  \end{itemize}
  \item \textbf{Are claims of violence, disorder, or political manipulation attributed to sources and substantiated?}
  \begin{itemize}[leftmargin=1.4em, itemsep=2pt, topsep=2pt]
    \item Rathee attributes \textbf{all pre-planned disorder and violence to the state}, substantiating his timeline with timestamped ground videos from independent YouTubers (\texttt{Shyam Meera Singh} at \texttt{[09:03--09:15]}, \texttt{[12:37]}; \texttt{Sarthak Goswami} at \texttt{[10:46]}, \texttt{[12:48]}; \texttt{Sandeesh Bhatia} at \texttt{[12:04]}, \texttt{[13:07]}; \texttt{Being Human / Ramit Verma} bleeding from the head at \texttt{[13:16--13:25]}), DCP New Delhi's \texttt{7:00 PM} Section 163 tweet (\texttt{[05:46--05:55]}), and Sonam Wangchuk's \texttt{11:05 AM} hospital letter (\texttt{[07:52--08:10]}). Conversely, claims of pre-planted brick trucks (\texttt{[00:41--01:29]}) and Amit Shah's motive for transferring Golcha (\texttt{"Perhaps he was upset that the protest was going on peacefully"}, \texttt{[04:36--04:45]}) rely on circumstantial inference.
  \end{itemize}
  \item \textbf{How are police and government actions described compared with protesters' actions? Are similar actions evaluated differently?}
  \begin{itemize}[leftmargin=1.4em, itemsep=2pt, topsep=2pt]
    \item Starkly polarized: Protesters prior to 20 July are \texttt{"100\% non-violent and peaceful"} (\texttt{[00:53]}, \texttt{[03:50]}), and breaking metro/street barricades on 20 July is framed as excusable democratic fervour (\texttt{"When the patience was out of hand, people jumped out of the barriers... People were so excited"}, \texttt{[07:17]}). Police actions, by contrast, are described as a \texttt{"war against the people of its own country"} (\texttt{[02:19]}, \texttt{[16:21]}), \texttt{"kidnapped Sonam Vangchuk"} (\texttt{[05:14]}), \texttt{"Jallianwala Bagh scene. General Dyer"} (\texttt{[10:08]}), and \texttt{"behaving like terrorists"} (\texttt{[11:54]}). Notably, Rathee \textbf{completely omits} the injury of 110--118 police personnel reported by Aaj Tak (\texttt{JXLIs0WtJAM}, \texttt{[18:09]}, \texttt{[19:43]}) and acknowledged by Sarthak Goswami (\texttt{EM1GeZufO6E}, \texttt{[04:37--04:53]}).
  \end{itemize}
  \item \textbf{Are students allowed to explain their demands in their own words, or does the presenter mainly speak for them?}
  \begin{itemize}[leftmargin=1.4em, itemsep=2pt, topsep=2pt]
    \item \textbf{Hybrid narrator-led documentary:} Rathee's scripted voiceover drives the chronological investigation, punctuated by short, emotionally intense ground-video clips of protesters, doctors (\texttt{"We are doctors! We came to help you... They beat us up"}, \texttt{[11:26]}), citizens confronting unbadged police (\texttt{"Where is your tag? Where is your nameplate?"}, \texttt{[11:37--12:25]}), and assaulted YouTubers (\texttt{[12:58--13:25]}).
  \end{itemize}
  \item \textbf{Does the outlet's framing differ between coverage of the NEET paper-leak issue before the protest and coverage once students begin protesting?}
  \begin{itemize}[leftmargin=1.4em, itemsep=2pt, topsep=2pt]
    \item Rathee links the protest directly to his long-standing thesis of creeping authoritarianism (\texttt{"I had exposed the government's taunting of Shahi [tanashahi] two years ago. Today, you can see the result of that"}, \texttt{[17:26]}): once students hit the streets, the issue escalates in his framing from educational incompetence to a full-blown colonial freedom struggle (\texttt{"East India Company"} vs \texttt{"West India Company"} at \texttt{[16:10]}).
  \end{itemize}
\end{enumerate}

\subsection*{Apply the five shared features (Event 2)}
\begin{longtable}{@{}>{\raggedright\arraybackslash}p{0.19\textwidth} >{\raggedright\arraybackslash}p{0.38\textwidth} >{\raggedright\arraybackslash}p{0.37\textwidth}@{}}
\toprule
\textbf{Feature} & \textbf{Exact wording + timestamp} & \textbf{What it does in context} \\
\midrule
\endhead
\textbf{1. Naming} & \texttt{"A Tana Shah government was ready to strangle the people"} \texttt{[02:09]}; \texttt{"kidnapped Sonam Vangchuk in an ambulance"} \texttt{[05:14]}; \texttt{"Narendra Modi, like a scared Tana Shah, leaves the Parliament building and runs away"} \texttt{[10:28]}; \texttt{"West India Company"} \texttt{[16:10]} & Replaces constitutional terms ("Union Government", "preventive detention", "PM") with delegitimizing authoritarian/colonial epithets (\texttt{Tana Shah} [dictator], \texttt{kidnapped}, \texttt{West India Company}---alluding to Gujarat-based corporate/political rule). \\
\addlinespace[4pt]
\textbf{2. Presupposition} & \texttt{"How exactly the youth of the country, the people of the country, were abused and how it was planned, To understand this better, we need to rewind"} \texttt{[02:55--03:05]} & Presupposes as proven fact both that the crackdown was a pre-meditated conspiracy (\texttt{how it was planned}) and that the victims represent the entire nation (\texttt{the youth of the country, the people of the country}). \\
\addlinespace[4pt]
\textbf{3. Agency} & \texttt{"The government started a war against the people of its own country."} \texttt{[02:19]} \newline  \texttt{"the policemen who were mercilessly beating the students, they had hidden their faces with a handkerchief."} \texttt{[11:45--11:54]} vs \texttt{"the barricades had begun to break"} \texttt{[07:17]} & Assigns maximal, intentional active agency to the government/police (\texttt{started a war}, \texttt{mercilessly beating}, \texttt{hidden their faces}), while using an agentless inchoative construction (\texttt{the barricades had begun to break}, \texttt{[07:17]}) that softens protester force against police barriers. \\
\addlinespace[4pt]
\textbf{4. Antagonisation} & \texttt{"Assuming that it was the Jallianwala Bagh scene. General Dyer was giving orders to fire... in a time, Indians were against the East India Company. And today, these Indians are against the West India Company."} \texttt{[10:08]}, \texttt{[16:10]} & Casts an elected Indian government and its security forces as alien colonial oppressors (\texttt{General Dyer}, \texttt{East/West India Company}) and masked \texttt{"terrorists"} (\texttt{[11:54]}) waging war on \texttt{"freedom fighters"} (\texttt{[15:14]}) and \texttt{"its own children"} (\texttt{[16:21]}). \\
\addlinespace[4pt]
\textbf{5. Intensification / softening} & \textbf{Intensification:} \texttt{"100\% non-violent"} \texttt{[00:53]}, \texttt{"very sinister plan"} \texttt{[02:01]}, \texttt{"strangle the people"} \texttt{[02:09]}, \texttt{"terrifying video"} \texttt{[10:08]}, \texttt{"The more the Tana Shahs in history thought themselves to be great, the more terrifying their end would be"} \texttt{[17:00]} & Uses absolute quantifiers (\texttt{100\%}), visceral bodily violence metaphors (\texttt{strangle}), and historical warnings of dictators' downfalls (\texttt{terrifying their end}) to maximize moral indignation, while conjectural hedges (\texttt{Perhaps}, \texttt{[04:36]}, \texttt{[05:37]}) briefly soften unproven motive claims regarding Satish Golcha's transfer. \\
\bottomrule
\end{longtable}

\subsection*{Short summary (Event 2)}
\begin{itemize}[leftmargin=1.5em, itemsep=3pt, topsep=4pt]
  \item \textbf{What interpretation of the event does this video encourage?} That a \texttt{"100\% non-violent"} student sit-in (\texttt{[00:53]}) grew so popular during Sonam Wangchuk's hunger strike (\texttt{[03:40--04:18]}) that a \texttt{"Tana Shah"} Modi-Shah regime (\texttt{[02:09]}, \texttt{[04:27]}) replaced the Delhi Police Commissioner, \texttt{"kidnapped"} Wangchuk (\texttt{[05:14]}), planted stone-filled vehicles to stage a pretext (\texttt{[00:23--01:29]}), shut down 5 metro stations and 150 internet sites (\texttt{[06:23--06:53]}), and unleashed masked/unbadged police and civil-dress goons---only for resilient students across India to reclaim Jantar Mantar that same evening (\texttt{[15:23--16:10]}).
  \item \textbf{Whose perspective is foregrounded, challenged, or missing?} \textit{Foregrounded:} Protesting students, CJP organisers, independent YouTubers assaulted on the ground (\texttt{Shyam Meera Singh}, \texttt{Sarthak Goswami}, \texttt{Sandeesh Bhatia}, \texttt{Ramit Verma}, \texttt{[12:48--13:35]}), and assaulted volunteer doctors (\texttt{[11:26]}). \textit{Challenged:} PM Modi's \texttt{"mother of democracy"} rhetoric (juxtaposed ironically at \texttt{[00:00]} and \texttt{[02:27]}), J.P. Nadda's dialogue offer (\texttt{[14:20--14:47]}), and Delhi Police crowd-control claims. \textit{Missing:} Any mention of the \texttt{110--118} injured police officers (including 5 IPS officers with fractures), opposition party turf wars, or Abhijit Dipke's political background.
  \item \textbf{What evidence supports your interpretation? What complicates it?} \textit{Supports:} Timestamped pre-dawn visuals (\texttt{04:40 AM} broken cars/stone trucks inside cordoned zones at \texttt{[00:23--01:29]}), official Section 163 / DMRC 5-station closure / internet shutdown records (\texttt{[05:55--06:53]}), Commissioner Golcha's abrupt pre-retirement replacement on 17 July (\texttt{[04:27--04:54]}), and clear multi-angle videos of unbadged/masked police and civil-clothed men slapping reporters and beating doctors (\texttt{[11:05--13:25]}). \textit{Complicates:} Total omission of police casualties (\texttt{110+} injured) creates a strictly one-sided casualty ledger, and equating tear-gas/riot positions outside Parliament with \texttt{"Jallianwala Bagh"} and \texttt{"General Dyer"} (\texttt{[10:08]}) relies on hyper-charged historical analogy rather than lethal live-fire evidence.
\end{itemize}
\vspace{4pt}\noindent\rule{\textwidth}{0.4pt}\vspace{4pt}

\section*{Common form: Blank excerpt record}
\textit{Copy or print this page whenever you find a passage worth analysing. Keep the exact wording and enough context to interpret it.}


\subsection*{Excerpt Record 1 (Event 1 --- Ram Mandir Donation Controversy)}
\begin{itemize}[leftmargin=1.5em, itemsep=3pt, topsep=4pt]
  \item \textbf{Outlet:} Dhruv Rathee | \textbf{Date:} \texttt{2026-07-04}
  \item \textbf{Video title:} \textit{Ram Mandir Scam is WORSE Than You Think | Dhruv Rathee}
  \item \textbf{Timestamp(s):} \texttt{[04:18--05:02]} | \textbf{Analyst:} Baani Solanki | \textbf{Language:} Hindi / Hinglish (Roman script)
  \item \textbf{Exact transcript excerpt:}
    \begin{quote}\small\ttfamily [04:18] This trust has a total of 15 members, 6 of whom are religious figures. And most of the other members, have a direct connection with Narendra Modi and BJP RSS. [04:28] For example, Narendra Modi's Principal Secretary, Nipendra Mishra, is the head of this temple construction committee. Similarly, the entire structure of this trust, [04:37] was prepared by the current Chief Election Commissioner, Gyanesh Kumar... [04:46] I'm telling you all this so that you understand that this is not an independent Trust The government is directly involved in this And the people associated with the government are running this temple\end{quote}
  \item \textbf{Feature(s):} $\boxtimes$ Naming $\boxtimes$ Presupposition $\boxtimes$ Agency $\square$ Antagonisation $\square$ Intensification / softening
  \item \textbf{Who is speaking? Is this the outlet's wording or a quotation?} Dhruv Rathee's own scripted analytical narration.
  \item \textbf{What does the wording imply, foreground, soften, or leave unstated?} Explicitly collapses the legal boundary between the Shri Ram Janmabhoomi Teerth Kshetra Trust and the Union Government (\texttt{"this is not an independent Trust"}), ensuring that any financial malpractice inside the temple is attributed upward as state/BJP-RSS culpability.
  \item \textbf{What context or evidence supports or complicates this interpretation?} Supported by the statutory formation of the Trust in February 2020 (\texttt{[04:08]}) and the documented career backgrounds of Nripendra Misra, Gyanesh Kumar, and Champat Rai---a point independently corroborated by Nilanjan Mukhopadhyay on \textbf{The Wire} (\texttt{7rUL8fGgRJA}, \texttt{[11:00--12:27]}). Complicated by \textbf{Aaj Tak} anchor Anjana Om Kashyap's distinction (\texttt{B4L1pz0o4dc}, \texttt{[10:57--11:24]}) that Nripendra Misra chaired only the \textit{Construction Committee} rather than day-to-day \textit{donation counting management}.
  \item \textbf{How does another outlet describe the same event or actor?} \textbf{Republic TV} (\texttt{0H-bruoEkaY}, \texttt{[04:08]}) asserts the exact opposite---\texttt{"there is no religious angle, there is no political angle, there is no government angle in this controversy"}---treating the Trust and the BJP government purely as swift law-enforcers (\texttt{[03:42--03:51]}).
  \item \textbf{My interpretation, with reasons:} Rathee uses institutional network mapping (\texttt{[04:18--05:02]}) early in the video so viewers cannot accept the later FIR (\texttt{[12:48]}) as an isolated failure of 8 low-level counting clerks.
\end{itemize}

\subsection*{Excerpt Record 2 (Event 1 --- Ram Mandir Donation Controversy)}
\begin{itemize}[leftmargin=1.5em, itemsep=3pt, topsep=4pt]
  \item \textbf{Outlet:} Dhruv Rathee | \textbf{Date:} \texttt{2026-07-04}
  \item \textbf{Video title:} \textit{Ram Mandir Scam is WORSE Than You Think | Dhruv Rathee}
  \item \textbf{Timestamp(s):} \texttt{[07:15--08:31]} | \textbf{Analyst:} Baani Solanki | \textbf{Language:} Hindi / Hinglish (Roman script)
  \item \textbf{Exact transcript excerpt:}
    \begin{quote}\small\ttfamily [07:15] 3 months before the arrest, the State Bank of India itself had warned that the counting staff had doubts. That they should be removed... [07:33] The bank was saying to remove them, but the temple-owners were saying, let them be. Why? ... [08:03] According to Mahipal Singh he told Champat Rai about the donation theft in 2021... [08:22] But do you know what happened? On the contrary, Mahipal Singh was removed from his job. Apart from that, CCTV footage of 7-8 months was also deleted.\end{quote}
  \item \textbf{Feature(s):} $\boxtimes$ Naming $\boxtimes$ Presupposition $\boxtimes$ Agency $\boxtimes$ Antagonisation $\boxtimes$ Intensification / softening
  \item \textbf{Who is speaking? Is this the outlet's wording or a quotation?} Rathee's narration combined with reported/embedded testimony from former Trust accounts in-charge Mahipal Singh (\texttt{[08:03--08:40]}).
  \item \textbf{What does the wording imply, foreground, soften, or leave unstated?} Naming Trust administrators \texttt{"the temple-owners"} (\texttt{[07:33]}) implies feudal, unaccountable proprietorship over public donations. Juxtaposing \texttt{"the counting staff"} staying employed against whistleblower Mahipal Singh being \texttt{"removed from his job"} and \texttt{"CCTV footage of 7-8 months... deleted"} implies active cover-up rather than mere administrative oversight.
  \item \textbf{What context or evidence supports or complicates this interpretation?} Supported by Mahipal Singh's on-camera interview clip (\texttt{[08:12--08:40]}) and reports of SBI's prior red flags. Complicated by Nripendra Misra's interview on \textbf{Aaj Tak} (\texttt{B4L1pz0o4dc}, \texttt{[05:09--05:19]}), which blamed \textit{SBI itself} for implementing counting supervision \texttt{"in the most non-serious manner"}.
  \item \textbf{How does another outlet describe the same event or actor?} \textbf{Aaj Tak} panelist Prof. Sangeet Ragi (\texttt{B4L1pz0o4dc}, \texttt{[10:36--10:48]}, \texttt{[12:25--12:56]}) shifts blame onto SBI and portrays Champat Rai as a selfless ascetic (\texttt{"do dhoti aur do kurte par baithne vala aadmi"}) whose integrity cannot be questioned.
  \item \textbf{My interpretation, with reasons:} This excerpt is Rathee's strongest evidentiary pivot from "petty theft by clerks" to "institutional complicity," because firing a whistleblower and deleting 7--8 months of CCTV cannot occur without managerial authority.
\end{itemize}

\subsection*{Excerpt Record 3 (Event 1 --- Ram Mandir Donation Controversy)}
\begin{itemize}[leftmargin=1.5em, itemsep=3pt, topsep=4pt]
  \item \textbf{Outlet:} Dhruv Rathee | \textbf{Date:} \texttt{2026-07-04}
  \item \textbf{Video title:} \textit{Ram Mandir Scam is WORSE Than You Think | Dhruv Rathee}
  \item \textbf{Timestamp(s):} \texttt{[14:13--15:58]} | \textbf{Analyst:} Baani Solanki | \textbf{Language:} Hindi / Hinglish (Roman script)
  \item \textbf{Exact transcript excerpt:}
    \begin{quote}\small\ttfamily [14:13] The CM of UP, Adityanath, is threatening the people who are questioning him. Don't take the fire of the Ram devotees. Stop playing with their faith. [14:24] Stop it. Who is playing with their faith? The scammer or the one who is questioning the scam? There are some blind devotees who look at Narendra Modi above Lord Ram... [14:54] Whenever you go to a temple, any temple, never donate money... [15:49] the one who shouts Jai Shri Ram the loudest, is not Ram in their hearts, but Ravan. The biggest criminals who rob in the name of Ram, and even more, those who are saving them.\end{quote}
  \item \textbf{Feature(s):} $\boxtimes$ Naming $\boxtimes$ Presupposition $\boxtimes$ Agency $\boxtimes$ Antagonisation $\boxtimes$ Intensification / softening
  \item \textbf{Who is speaking? Is this the outlet's wording or a quotation?} Paraphrased warning from UP CM Yogi Adityanath (\texttt{[14:13]}) followed by Dhruv Rathee's direct first-person rebuttal and moral peroration.
  \item \textbf{What does the wording imply, foreground, soften, or leave unstated?} Reverses the speech act of Yogi Adityanath's warning by categorizing it as \texttt{"threatening"} (\texttt{[14:13]}), labels partisan apologists \texttt{"blind devotees"} (\texttt{[14:24]}), and culminates in a mythological inversion (\texttt{Ram} vs \texttt{Ravan}, \texttt{[15:49]}) where state protectors of temple thieves are worse than the thieves themselves (\texttt{"and even more, those who are saving them"}).
  \item \textbf{What context or evidence supports or complicates this interpretation?} Supported by Yogi Adityanath's real Ayodhya speech warning against \texttt{"analagadh tippani"} hurting devotees' sentiments (heard verbatim on \textbf{Aaj Tak} \texttt{B4L1pz0o4dc} at \texttt{[03:02--03:13]}). Complicated by Rathee's sweeping imperative at \texttt{[14:54]} (\texttt{"Whenever you go to a temple, any temple, never donate money"}), which hands critics easy ammunition to accuse him of anti-temple bias.
  \item \textbf{How does another outlet describe the same event or actor?} \textbf{Sarthak Goswami} (\texttt{lb7gZnyZH0c}, \texttt{[06:16--06:26]}) similarly contrasts Lord Ram's character with abusive slogan-chanters (\texttt{"Ram ka charitra chori"}), whereas \textbf{Republic TV} (\texttt{0H-bruoEkaY}, \texttt{[01:12]}, \texttt{[03:22]}) hails Yogi Adityanath's exact same intervention as \texttt{"zero tolerance for any form of corruption"}.
  \item \textbf{My interpretation, with reasons:} Rathee reclaims Hindu ethical vocabulary (\textit{Swades} lyric on Ram vs Ravan) to avoid sounding anti-Hindu while delivering his sharpest anti-BJP and anti-institutional-donation critique.
\end{itemize}

\subsection*{Excerpt Record 4 (Event 2 --- NEET Paper Leak Protest)}
\begin{itemize}[leftmargin=1.5em, itemsep=3pt, topsep=4pt]
  \item \textbf{Outlet:} Dhruv Rathee | \textbf{Date:} \texttt{2026-07-24}
  \item \textbf{Video title:} \textit{How CJP defeated BJP at Jantar Mantar Protests? | Full Story of 20th July | Dhruv Rathee}
  \item \textbf{Timestamp(s):} \texttt{[00:23--02:19]} | \textbf{Analyst:} Baani Solanki | \textbf{Language:} Hindi / Hinglish (Roman script)
  \item \textbf{Exact transcript excerpt:}
    \begin{quote}\small\ttfamily [00:23] On 20th July, at 4.40 am, some people protesting at the Jantar Mantar come out... [00:41] right outside the Jantar Mantar, suddenly, some people saw trucks full of stones... [00:53] everything was 100\% non-violent and peaceful. [01:03] But now, near the nearby Parliament, a broken glass car was also seen standing... [01:29] It was planted to defame the protest by the government. That same night, the Delhi Police removed all the metal detectors and security checks from the Jantar Mantar... [02:01] A very sinister plan was being created behind this... [02:09] A Tana Shah government was ready to strangle the people to save its power. [02:19] The government started a war against the people of its own country.\end{quote}
  \item \textbf{Feature(s):} $\boxtimes$ Naming $\boxtimes$ Presupposition $\boxtimes$ Agency $\boxtimes$ Antagonisation $\boxtimes$ Intensification / softening
  \item \textbf{Who is speaking? Is this the outlet's wording or a quotation?} Dhruv Rathee's investigative narration incorporating an on-site protester's written note (\texttt{"It was planted to defame the protest"}, \texttt{[01:29]}) and ground video (\texttt{[01:41]}).
  \item \textbf{What does the wording imply, foreground, soften, or leave unstated?} Pre-emptively inoculates protesters against \textit{any} subsequent accusation of stone-pelting or vandalism on 20 July by framing stones and damaged cars as a state false-flag operation (\texttt{"sinister plan"}, \texttt{"planted to defame"}), escalated via war/dictatorship metaphors (\texttt{"Tana Shah government"}, \texttt{"started a war"}).
  \item \textbf{What context or evidence supports or complicates this interpretation?} Supported by timestamped \texttt{04:40 AM} ground footage inside heavily cordoned zones and protesters' real-time warnings before clashes began. Complicated by \textbf{Aaj Tak}'s interviews with injured Addl. DCPs Anand Mishra and Jagdev (\texttt{JXLIs0WtJAM}, \texttt{[21:12--23:14]}), who suffered fractures from stones thrown from within the crowd at barricades later that day.
  \item \textbf{How does another outlet describe the same event or actor?} \textbf{Republic TV} (\texttt{HTvb6E\_1Qdc}, \texttt{[01:37--01:47]}) attributes 100\% of street violence to \texttt{"rabid, bloodthirsty mobs"} attacking police and RAF, never mentioning pre-dawn stone trucks or removed metal detectors.
  \item \textbf{My interpretation, with reasons:} Opening the video with \texttt{04:40 AM} pre-dawn anomalies establishes a powerful forensic frame: before showing a single clash at \texttt{[07:43]}, Rathee primes the viewer to interpret all disorder as state sabotage.
\end{itemize}

\subsection*{Excerpt Record 5 (Event 2 --- NEET Paper Leak Protest)}
\begin{itemize}[leftmargin=1.5em, itemsep=3pt, topsep=4pt]
  \item \textbf{Outlet:} Dhruv Rathee | \textbf{Date:} \texttt{2026-07-24}
  \item \textbf{Video title:} \textit{How CJP defeated BJP at Jantar Mantar Protests? | Full Story of 20th July | Dhruv Rathee}
  \item \textbf{Timestamp(s):} \texttt{[10:08--10:28]} \& \texttt{[11:37--12:16]} | \textbf{Analyst:} Baani Solanki | \textbf{Language:} Hindi / Hinglish (Roman script)
  \item \textbf{Exact transcript excerpt:}
    \begin{quote}\small\ttfamily [10:08] Where security forces were taking firing positions. Assuming that it was the Jallianwala Bagh scene. General Dyer was giving orders to fire... [10:28] Meanwhile, Narendra Modi, like a scared Tana Shah, leaves the Parliament building and runs away... [11:37] The most disturbing thing was that the police officers... the nameplates of the vehicles were removed from their uniform. Where is your tag? Where is your nameplate? ... [11:54] they had hidden their faces with a handkerchief. They were behaving like terrorists. Apart from this, we also saw that some people came in ordinary clothes, [12:04] were roaming around with the policemen...\end{quote}
  \item \textbf{Feature(s):} $\boxtimes$ Naming $\square$ Presupposition $\boxtimes$ Agency $\boxtimes$ Antagonisation $\boxtimes$ Intensification / softening
  \item \textbf{Who is speaking? Is this the outlet's wording or a quotation?} Dhruv Rathee's voiceover alternating with direct on-the-ground audio of protesters questioning police (\texttt{"Where is your tag? Where is your nameplate?"}, \texttt{[11:37--11:45]}, \texttt{[12:16]}).
  \item \textbf{What does the wording imply, foreground, soften, or leave unstated?} Performs a radical \textbf{reversal of the "terrorist" label}: whereas pro-government channels call student protesters terrorists, Rathee applies \texttt{"behaving like terrorists"} (\texttt{[11:54]}) and \texttt{"General Dyer"} (\texttt{[10:08]}) to masked, unbadged security personnel, while mocking PM Modi's exit from Parliament as cowardly flight (\texttt{"like a scared Tana Shah... runs away"}, \texttt{[10:28]}).
  \item \textbf{What context or evidence supports or complicates this interpretation?} Unbadged policemen and civil-clothed assailants are corroborated across \textbf{Sarthak Goswami} (\texttt{EM1GeZufO6E}, \texttt{[03:57--04:07]}), \textbf{Mohak Mangal} (\texttt{v-RikSqb0Vw}, \texttt{[01:38]}, \texttt{[04:38]}), \textbf{Aaj Tak} (Akhilesh Yadav quote in \texttt{JXLIs0WtJAM}, \texttt{[06:50--07:08]}), and \textbf{The Wire} (\texttt{pYaTVvHeGoM}, \texttt{[07:34--07:42]}). Even the Jallianwala Bagh visual comparison is independently noted by Naseeruddin Shah on \textbf{The Wire} (\texttt{pYaTVvHeGoM}, \texttt{[10:59--11:08]}). However, security forces taking positions at Parliament's iron gates during a breach attempt (\texttt{[09:59--10:17]}) is standard high-security Parliamentary protocol, complicating the \texttt{"General Dyer"} analogy.
  \item \textbf{How does another outlet describe the same event or actor?} \textbf{Aaj Tak} (\texttt{JXLIs0WtJAM}, \texttt{[20:02--20:41]}) quotes DCP Sachin Sharma describing the exact same deployments neutrally as \texttt{"priventiv arenjment"} around \texttt{"vaital instoleshanans"} [vital installations].
  \item \textbf{My interpretation, with reasons:} Rathee combines undeniable procedural illegality (covered vehicle plates, missing name badges) with colonial/terrorist counter-naming to strip the state apparatus of moral authority.
\end{itemize}

\subsection*{Excerpt Record 6 (Event 2 --- NEET Paper Leak Protest)}
\begin{itemize}[leftmargin=1.5em, itemsep=3pt, topsep=4pt]
  \item \textbf{Outlet:} Dhruv Rathee | \textbf{Date:} \texttt{2026-07-24}
  \item \textbf{Video title:} \textit{How CJP defeated BJP at Jantar Mantar Protests? | Full Story of 20th July | Dhruv Rathee}
  \item \textbf{Timestamp(s):} \texttt{[16:10--17:00]} | \textbf{Analyst:} Baani Solanki | \textbf{Language:} Hindi / Hinglish (Roman script)
  \item \textbf{Exact transcript excerpt:}
    \begin{quote}\small\ttfamily [16:10] This will be an inspiration for many generations to come. That in a time, Indians were against the East India Company. And today, these Indians are against the West India Company. [16:21] Today, the world is witnessing how a country called the Internationalist [Mother of Democracy], conspired against its own children, and waged a war on its own people. [16:29] Do you know what was the strangest part? There was only one demand from the children. That the exams should not be stolen... [16:44] I just want to say this to the Prime Minister, Shame on you, Narendra Modi. The government that is beating up the children, it's time for them to leave... [17:00] The more the Tana Shahs in history thought themselves to be great, the more terrifying their end would be.\end{quote}
  \item \textbf{Feature(s):} $\boxtimes$ Naming $\boxtimes$ Presupposition $\boxtimes$ Agency $\boxtimes$ Antagonisation $\boxtimes$ Intensification / softening
  \item \textbf{Who is speaking? Is this the outlet's wording or a quotation?} Dhruv Rathee speaking directly to the camera and addressing Prime Minister Modi in the second person (\texttt{"Shame on you, Narendra Modi"}, \texttt{[16:44]}).
  \item \textbf{What does the wording imply, foreground, soften, or leave unstated?} Contrasts maximal state aggression (\texttt{"conspired against its own children, and waged a war"}) with minimal, innocent youth agency (\texttt{"only one demand from the children. That the exams should not be stolen"}), culminating in a direct political prophecy of regime fall (\texttt{"time for them to leave"}, \texttt{"terrifying their end"}).
  \item \textbf{What context or evidence supports or complicates this interpretation?} Remarkably, \textbf{the very next day (25 July 2026)} Education Minister Dharmendra Pradhan resigned and the government accepted CJP's core demands (confirmed in Sarthak \texttt{EM1GeZufO6E} \texttt{[00:08]} and Republic \texttt{HTvb6E\_1Qdc} \texttt{[00:00]}), showing immense public pressure---yet also complicating Rathee's \texttt{"Tana Shah"} thesis because an unyielding dictatorship rarely concedes ministerial resignation within 5 days of a march.
  \item \textbf{How does another outlet describe the same event or actor?} \textbf{Republic TV} (\texttt{HTvb6E\_1Qdc}, \texttt{[01:56--02:56]}) spins Pradhan's resignation the next day not as a defeat of the BJP, but as a masterstroke whereby \texttt{"the government outplayed"} Rahul Gandhi and ended the street crisis.
  \item \textbf{My interpretation, with reasons:} Rathee transforms a single-issue exam-reform protest into an existential plebiscite on Modi's rule (\texttt{"West India Company"}), mobilizing youth grievance through anti-colonial and familial (\texttt{"children"}) framing.
\end{itemize}
\vspace{4pt}\noindent\rule{\textwidth}{0.4pt}\vspace{4pt}

\section*{Common form: End-of-event summary}

\subsection*{End-of-Event Summary --- Event 1 (Ram Mandir Donation Controversy)}
\begin{itemize}[leftmargin=1.5em, itemsep=3pt, topsep=4pt]
  \item \textbf{The main way this outlet framed the event was:} As a comprehensive, six-year institutional plunder (\texttt{"From the purchase of land, to the construction and donation of the Ram Mandir, there have been frauds"}, \texttt{[11:56]}) run by a government-linked BJP-RSS Trust that ignored 2020 private audit and SBI warnings, fired whistleblower Mahipal Singh, deleted 7--8 months of CCTV, and scapegoated 8 low-level counting staff while shielding senior figures like Champat Rai under a \texttt{"dhritarashtra"} PMO (\texttt{[04:18--08:40]}, \texttt{[12:48--14:04]}).
  \item \textbf{The strongest linguistic evidence was (quote + timestamp):} \texttt{"this is not an independent Trust. The government is directly involved... the one who shouts Jai Shri Ram the loudest, is not Ram in their hearts, but Ravan. The biggest criminals who rob in the name of Ram, and even more, those who are saving them."} (\texttt{[04:46]} \& \texttt{[15:49--15:58]}).
  \item \textbf{The shared feature(s) most clearly illustrated were:} \textbf{Presupposition} (\texttt{Scam is WORSE Than You Think}; \texttt{"understand how... the public was looted"}, \texttt{[01:44]}), \textbf{Naming} (\texttt{temple-owners} \texttt{[07:33]}, \texttt{dhritarashtra} \texttt{[13:54]}, \texttt{Ravan} \texttt{[15:49]}), and \textbf{Antagonisation} (sincere devotees \& whistleblowers vs hypocritical \texttt{"Jai Shri Ram"} scammers and \texttt{"blind devotees"}, \texttt{[14:24--15:58]}).
  \item \textbf{An alternative interpretation of the same evidence is:} That the cash-counting theft (\texttt{\textasciitilde{}Rs.\,80 lakh} verified on CCTV in Republic \texttt{0H-bruoEkaY} \texttt{[07:26]}) resulted from lax operational protocols between SBI and counting contractors (as Nripendra Misra admitted on Aaj Tak \texttt{B4L1pz0o4dc} \texttt{[04:35--06:05]}), and that arresting 8 direct culprits within days alongside Champat Rai and Anil Mishra's resignations (\texttt{[12:56--13:05]}) demonstrated swift corrective accountability rather than PMO-directed collusion.
  \item \textbf{What I cannot conclude from this transcript alone:} Whether the unverified upper-end allegations (\texttt{Rs.\,200 crore} missing donations at \texttt{[01:28]}, \texttt{Rs.\,50 crore} amassed by Tinnu Yadav at \texttt{[07:54]}, or \texttt{12 IB reports} sent to the PM at \texttt{[13:42]}) are backed by forensic bank/SIT records, or whether senior trustees personally received kickbacks.
\end{itemize}

\subsection*{End-of-Event Summary --- Event 2 (Students' Protest Regarding NEET Paper Leak)}
\begin{itemize}[leftmargin=1.5em, itemsep=3pt, topsep=4pt]
  \item \textbf{The main way this outlet framed the event was:} As a pre-meditated authoritarian \texttt{"war against the people of its own country"} (\texttt{[02:19]}) in which a \texttt{"Tana Shah"} government (\texttt{[02:09]}) sacked the Delhi Police Commissioner, \texttt{"kidnapped"} hunger-striker Sonam Wangchuk (\texttt{[05:14]}), planted stone trucks (\texttt{[00:41--01:29]}), imposed an internet/metro siege (\texttt{[06:12--06:53]}), and unleashed masked/unbadged police and civil-dress goons on \texttt{"100\% non-violent"} \texttt{"children"}---who nevertheless triumphed like anti-colonial freedom fighters (\texttt{[15:06--16:37]}).
  \item \textbf{The strongest linguistic evidence was (quote + timestamp):} \texttt{"A Tana Shah government was ready to strangle the people to save its power. The government started a war against the people of its own country... Indians were against the East India Company. And today, these Indians are against the West India Company."} (\texttt{[02:09--02:19]} \& \texttt{[16:10]}).
  \item \textbf{The shared feature(s) most clearly illustrated were:} \textbf{Naming} (\texttt{Tana Shah}, \texttt{kidnapped}, \texttt{General Dyer}, \texttt{behaving like terrorists}, \texttt{West India Company}), \textbf{Agency} (attributing full pre-planned agency to Modi/Shah/Delhi Police while rendering broken barricades agentless at \texttt{[07:17]} and omitting all \texttt{110+} police injuries), and \textbf{Intensification} (\texttt{100\% non-violent}, \texttt{strangle}, \texttt{waged a war on its own children}).
  \item \textbf{An alternative interpretation of the same evidence is:} That closing nearby metro stations and erecting Section 163 barricades around Parliament during a Monsoon Session (\texttt{[05:55--06:53]}) were preventive security measures for a high-security constitutional zone, which escalated into violent two-way clashes once lakhs of protesters broke welded barricades (\texttt{[07:17--08:20]}) and pelted stones injuring over 110 police personnel (Aaj Tak \texttt{JXLIs0WtJAM}, \texttt{[19:43--23:14]}).
  \item \textbf{What I cannot conclude from this transcript alone:} Who actually parked the broken vehicles/brick trucks filmed at \texttt{04:40 AM} (\texttt{[00:23--01:12]}), why Satish Golcha was administratively transferred on 17 July (\texttt{[04:27--04:45]}), or how many police officers and protesters respectively initiated physical force at each flashpoint.
\end{itemize}
\textit{Research reminder: interpret wording in context, distinguish allegations from established facts, and support claims with exact excerpts. Mixed or contradictory evidence is worth recording.}


\clearpage

\begin{center}
    {\large\bfseries LANGUAGE \& SOCIETY PROJECT}\\[3pt]
    {\LARGE\bfseries Transcript Analysis Worksheet}\\[3pt]
    {\normalsize Qualitative analysis of political news discourse}
\end{center}
\vspace{2pt}
\noindent\rule{\textwidth}{0.4pt}
\vspace{4pt}


\section*{Creator 3: Mohak Mangal (Independent Ground-Reportage \& Documentary YouTuber)}
\begin{itemize}[leftmargin=1.5em, itemsep=3pt, topsep=4pt]
  \item \textbf{Analyst:} Baani Solanki (\texttt{2025114016})
  \item \textbf{Channel / Outlet:} Mohak Mangal (\texttt{@mohak\_mangal}, YouTube)
  \item \textbf{Ideological / Editorial Positioning:} Civic-constitutionalist, evidence-driven documentary journalism; foregrounds citizen testimonies, legal/police manuals, and historical irony over loud studio monologue, contrasting state/TV propaganda soundbites directly against lived ground realities.
  \item \textbf{Transcript Files Referenced:}
  \begin{itemize}[leftmargin=1.4em, itemsep=2pt, topsep=2pt]
    \item Event 1: \textbf{SKIPPED} (No video published by Mohak Mangal on Event 1; explicitly marked \textit{Skip} in \texttt{lns\_project\_instructions.pdf})
    \item Event 2: \texttt{transcripts/event2\_neet\_protest/03\_mohak\_mangal\_v-RikSqb0Vw.txt}
  \end{itemize}
\end{itemize}
\vspace{4pt}\noindent\rule{\textwidth}{0.4pt}\vspace{4pt}

\section*{Guidelines: Five features to look for in every event}
\begin{longtable}{@{}>{\raggedright\arraybackslash}p{0.17\textwidth} >{\raggedright\arraybackslash}p{0.20\textwidth} >{\raggedright\arraybackslash}p{0.24\textwidth} >{\raggedright\arraybackslash}p{0.32\textwidth}@{}}
\toprule
\textbf{Feature} & \textbf{What to look for} & \textbf{Illustrative example} & \textbf{Questions to answer} \\
\midrule
\endhead
\textbf{1. Naming} & What are people, groups, actions, or events called? & "Students", "protesters", "agitators"; "irregularity", "scam", "controversy". & What exact label is used? Who or what does it refer to? Who uses it? What alternative label could have been used? \\
\addlinespace[4pt]
\textbf{2. Presupposition} & What is treated as already true or taken for granted? & "Why are students creating chaos again?" may presuppose earlier disruption. & What does the wording assume? Is that assumption supported, disputed, or left unexamined? Does it appear in a question or words such as "again", "still", or "finally"? \\
\addlinespace[4pt]
\textbf{3. Agency} & Who is shown doing something? Who is acted upon? Who disappears? & "Police injured protesters" vs. "Protesters were injured." & Who is the actor? Who experiences the action? Is responsibility named, backgrounded, or omitted? \\
\addlinespace[4pt]
\textbf{4. Antagonisation} & Who is positioned against whom? & "Ordinary citizens" vs. "political troublemakers". & Which groups are contrasted? Is one presented as more legitimate, moral, patriotic, dangerous, or disloyal? Is an "us vs. them" division implied? \\
\addlinespace[4pt]
\textbf{5. Intensification / softening} & What is exaggerated, minimised, softened, or euphemised? & "Violent mob" vs. "peaceful demonstrators"; "incident" instead of a direct description. & Which words heighten emotion or seriousness? Which minimise it? Are metaphors, dramatic adjectives, hedges, or euphemisms used? Does the wording match the evidence? \\
\bottomrule
\end{longtable}
\textit{Reminder: A label or implication is not automatically false or manipulative. Examine the surrounding context and evidence.}

\vspace{4pt}\noindent\rule{\textwidth}{0.4pt}\vspace{4pt}

\section*{Event1. Ram Mandir donation controversy}
\textbf{Focus:} How is the controversy framed---as financial accountability, religious sentiment, political motivation, or something else? Be precise about the specific allegation or episode covered.

\begin{itemize}[leftmargin=1.5em, itemsep=3pt, topsep=4pt]
  \item \textbf{Outlet:} Mohak Mangal (YouTube)
  \item \textbf{Date:} N/A (\textit{Skipped per \texttt{lns\_project\_instructions.pdf}})
  \item \textbf{Video title:} \textit{No video relevant to scam; there is video regarding history, so skipping} (\texttt{SKIPPED})
  \item \textbf{Timestamp(s):} N/A (\textit{Skipped})
  \item \textbf{Analyst:} Baani Solanki (\texttt{2025114016})
  \item \textbf{Language:} N/A (\textit{Skipped})
  \item \textbf{Analytical Note for Cross-Outlet Comparison:} Even this \textit{silence / non-coverage} is sociolinguistically noteworthy when comparing Independent YouTube vs Mainstream TV: unlike studio debate shows (Aaj Tak, Republic TV) that cover daily partisan flashpoints, or weekly news-roundup/anti-BJP channels (Sarthak Goswami, Dhruv Rathee), Mohak Mangal selectively produces focused civic-policy and ground-reportage documentaries---intervening in Event 2 (\texttt{"We the people of India..."}) when constitutional rights (\texttt{Article 19}) and exam governance directly affected students on the streets of Delhi.
\end{itemize}
\vspace{4pt}\noindent\rule{\textwidth}{0.4pt}\vspace{4pt}

\section*{Event2. Students' protest regarding NEET paper leak}
\textbf{Focus:} How does coverage construct the legitimacy of student protest and responsibility for the underlying problem?

\begin{itemize}[leftmargin=1.5em, itemsep=3pt, topsep=4pt]
  \item \textbf{Outlet:} Mohak Mangal (YouTube)
  \item \textbf{Date:} \texttt{2026-07-22} \textit{(filmed on the ground on 20--21 July immediately after the "Black Monday" lathi-charge)}
  \item \textbf{Video title:} \textit{"We the people of India..."} (\url{https://youtu.be/v-RikSqb0Vw?si=T4KGrNSGutHG4YYO})
  \item \textbf{Timestamp(s):} Relevant Event 2 segment \texttt{[00:00--11:09]} (\texttt{100\%} of video; total duration \texttt{11:09}; zero sponsor segments or unrelated topics).
  \item \textbf{Analyst:} Baani Solanki (\texttt{2025114016})
  \item \textbf{Language:} Hindi / Hinglish \& English (Native bilingual ground reportage, Romanized script; non-AI-dubbed)
\end{itemize}

\subsection*{Event-specific questions (Event 2)}
\begin{enumerate}[leftmargin=1.6em, itemsep=4pt, topsep=4pt]
  \item \textbf{Are participants called "students", "candidates", "protesters", "agitators", "mobs", or "political pawns"?}
  \begin{itemize}[leftmargin=1.4em, itemsep=2pt, topsep=2pt]
    \item Mohak and his ground interviewees consistently identify participants as \texttt{"students"} (\texttt{[00:10]}, \texttt{[01:03]}, \texttt{[03:54]}), \texttt{"bachchon"} / \texttt{"10tha klas ka bachcha"} / \texttt{"16 se 22 saal ke ladkon"} [children / 10th-class child / 16-to-22-year-old boys] (\texttt{[01:57--02:06]}, \texttt{[04:30]}, \texttt{[04:46--05:04]}), \texttt{"ladakiyan"} / \texttt{"aantiyon"} / \texttt{"ledi"} [girls / aunties / lady] (\texttt{[01:57]}, \texttt{[02:51]}, \texttt{[04:56]}), \texttt{"net eksapiritent"} [NEET aspirant] (\texttt{[04:30]}), \texttt{"risentli grejueted"} youth (\texttt{[08:16]}), and \texttt{"sitijans"} / \texttt{"We the people of India"} (\texttt{[Title]}, \texttt{[09:10]}). Crucially, Mohak intercuts \textbf{uncommented audio clips of ruling-camp politicians, judges, and TV anchors} calling them \texttt{"tukde-tukde"}, \texttt{"dahashat gardon ka bi tim"} [B-team of terrorists] (\texttt{[00:44--00:53]}), \texttt{"goons [gsa] oph da CJP"} (\texttt{[01:03]}, \texttt{[10:56]}), \texttt{"yangastars laik kokroch"} [youngsters like cockroaches] (\texttt{[02:59]}), and \texttt{"not students... CJP sponsard gundaj"} (\texttt{[06:06]})---letting the visual/verbal contrast expose those smears without shouting back.
  \end{itemize}
  \item \textbf{Are their stated demands explained, summarised, mocked, questioned, or left out?}
  \begin{itemize}[leftmargin=1.4em, itemsep=2pt, topsep=2pt]
    \item \textbf{Articulated directly by the protesters, parents, and teachers themselves:} fair, leak-free exams after repeated 2024--2026 leaks (\texttt{[03:27--04:12]}, \texttt{[05:14]}), accountability for \texttt{98} leaked papers over 10 years and student suicides (\texttt{[05:14]}, \texttt{[06:34]}), flawed CET/NEET score normalization (\texttt{630/800} reduced to \texttt{530} at \texttt{[08:44--09:05]}), and ministerial listening/accountability from Dharmendra Pradhan (\texttt{"Dharmendra Pradhan mast lisan too as end aavar gavarnament ij obligeted too lisan too its sitijans"}, \texttt{[09:10--09:18]}).
  \end{itemize}
  \item \textbf{Does the presenter focus on the original grievance, disruption, political affiliations, or public inconvenience?}
  \begin{itemize}[leftmargin=1.4em, itemsep=2pt, topsep=2pt]
    \item Focuses deeply on the \textbf{human and constitutional cost of the original grievance}---4-year repeaters losing seats (\texttt{[03:27--04:12]}), mothers fearing student suicides (\texttt{[06:14--07:01]}), schoolteachers unable to teach the NCERT \texttt{"demokresi chaiptar"} (\texttt{[07:14--07:58]}), brain drain abroad (\texttt{[07:32--07:40]}), and future patient safety from \texttt{"doktar aaj phek"} [fake doctors buying leaked papers] (\texttt{[07:40--07:58]})---alongside \textbf{police breach of minimum-force guidelines} (\texttt{[02:06--02:59]}).
  \end{itemize}
  \item \textbf{Who is represented as responsible: examination authorities, government, students, political parties, or others?}
  \begin{itemize}[leftmargin=1.4em, itemsep=2pt, topsep=2pt]
    \item \textbf{Education Minister Dharmendra Pradhan and the Union Government} (for hypocrisy, denying leaks via \texttt{"deyar ij no evidens oph paper leak too last sevan iyars"} at \texttt{[05:04]}, monologue governance via \texttt{"man ki baat"} instead of listening to youth at \texttt{[05:14--05:47]}, and dividing citizens via \texttt{"Hindu muslim"} rhetoric at \texttt{[05:47--05:57]}), and \textbf{Delhi Police} (for removing nameplates and violating official lathi-charge guidelines by striking heads, eyes, and collarbones behind closed exit gates at \texttt{[01:38--02:59]}, \texttt{[04:38--05:04]}, \texttt{[09:56--10:06]}).
  \end{itemize}
  \item \textbf{Are questions framed around whether students have legitimate grievances, or why they are causing disruption?}
  \begin{itemize}[leftmargin=1.4em, itemsep=2pt, topsep=2pt]
    \item Framed entirely around \textbf{legitimate citizen grievances, state hypocrisy, and democratic accountability}: e.g., \texttt{"toh ye sarkar itna bachchon se darti kyon hai?"} [why is this government so afraid of children?] (\texttt{[04:30]}), \texttt{"main ej a tichar demokresi chaiptar padha rahe ho. so vhat aaee tich student?"} (\texttt{[07:14]}), and Supriya's pointed question \texttt{"agar wo gilt mein nahin hai toh kyon itne kaimre hain? kyon logon ko hataya ja raha hai?"} (\texttt{[10:17--10:27]}).
  \end{itemize}
  \item \textbf{Are claims of violence, disorder, or political manipulation attributed to sources and substantiated?}
  \begin{itemize}[leftmargin=1.4em, itemsep=2pt, topsep=2pt]
    \item \textbf{Police violence} is substantiated through \textbf{first-person injured eyewitnesses showing/describing wounds on camera} (leg blows from a \texttt{"thri star vala"} without a nameplate at \texttt{[01:38--01:57]}, a 20-year-old woman's torn shoulder at \texttt{[02:51]}, torn clothes of female protesters and an eye injury at \texttt{[04:56--05:04]}, beating behind closed exit doors at \texttt{[09:46--10:06]}) measured directly against verbatim \textbf{Indian Police Minimum-Force Guidelines} read out by Mohak at \texttt{[02:06--02:39]}. Meanwhile, state/media claims of \texttt{"tukde-tukde"}, \texttt{"dahashat gardon ka bi tim"} (\texttt{[00:44]}), \texttt{"brootal vaylens kairid aaut baay da [goons] oph da CJP"} (\texttt{[01:03]}, \texttt{[10:56]}), and \texttt{"CJP sponsard gundaj"} (\texttt{[06:06]}) are played verbatim as \textbf{ironic counter-clips} immediately refuted by the next speaker on screen (e.g., \texttt{[06:06]} \texttt{"diz aar not students... diz aar CJP sponsard gundaj"} cuts straight to \texttt{[06:14]} a mother of two 10th/4th-grade daughters explaining why she came to sit for all India's children!).
  \end{itemize}
  \item \textbf{How are police and government actions described compared with protesters' actions? Are similar actions evaluated differently?}
  \begin{itemize}[leftmargin=1.4em, itemsep=2pt, topsep=2pt]
    \item Mohak's opening masterstroke (\texttt{[00:00--01:03]}) compares \textbf{the exact same action across two eras}: in \texttt{1997}, when \texttt{1,500} students marched to the Bhubaneswar Secretariat over an Odisha paper leak and suffered brutal police lathi-charge fracturing young ABVP National Secretary \texttt{"mukkoo bhaee"} (Dharmendra Pradhan), he was a heroic student victim; in \texttt{2026}, \texttt{"dait da sem gaay jo 1997 mein paper leak ke khilaph protest kar rahe the. aaj unke khilaph protest ho rahi hai aur unki sarkar ab students par lathi charj kar rahi hai"} (\texttt{[00:53--01:03]}). Protesters in 2026 are consistently depicted as \texttt{"pisaphul"} / \texttt{"shantipoorn"} (\texttt{[00:10]}, \texttt{[04:38]}, \texttt{[05:57]}, \texttt{[10:36]}) following \texttt{"satya aur ahinsa ka marg"} (\texttt{[10:27--10:36]}), whereas police action is \texttt{"inahyooman"} (\texttt{[10:06]}) and unlawful (\texttt{[02:16--02:51]}).
  \end{itemize}
  \item \textbf{Are students allowed to explain their demands in their own words, or does the presenter mainly speak for them?}
  \begin{itemize}[leftmargin=1.4em, itemsep=2pt, topsep=2pt]
    \item \textbf{Students and ordinary citizens speak overwhelmingly in their own words!} Mohak himself speaks for \textbf{less than 2 minutes} of the entire \texttt{11:09} runtime (\texttt{[00:00--00:44]}, \texttt{[00:53--01:03]}, \texttt{[02:06--02:59]}, and brief interviewer prompts at \texttt{[01:12]}, \texttt{[04:21]}, \texttt{[08:44]}, \texttt{[09:10]}). Over \textbf{80\% of the video} consists of uninterrupted, raw voices of an injured male protester (\texttt{[01:12--02:06]}), a female 4-year NEET repeater (\texttt{[03:18--04:21]}), a Hyderabad hospital worker whose brother is a NEET aspirant (\texttt{[04:21--05:04]}), an anguished citizen critiquing \texttt{"man ki baat"} and communal polarization (\texttt{[05:14--06:05]}), a mother of two schoolgirls (\texttt{[06:14--07:08]}), a civics teacher (\texttt{[07:14--07:58]}), recent graduates (\texttt{[08:16--09:18]}), and Supriya from Konkan, Maharashtra (\texttt{[09:29--10:56]}).
  \end{itemize}
  \item \textbf{Does the outlet's framing differ between coverage of the NEET paper-leak issue before the protest and coverage once students begin protesting?}
  \begin{itemize}[leftmargin=1.4em, itemsep=2pt, topsep=2pt]
    \item Rather than shifting from sympathy (on paper leaks) to alarmism (once street protests begin), Mohak treats the street protest as the \textbf{authentic embodiment of the Constitutional Preamble (\texttt{"We the people of India..."})}, arguing through his interviewees that when 98 papers leak and institutional channels ignore youth for 25 days (\texttt{[05:14--05:38]}), peaceful street assembly is democracy's essential safeguard (\texttt{[04:46]}, \texttt{[09:10--09:18]}, \texttt{[10:27--10:56]}).
  \end{itemize}
\end{enumerate}

\subsection*{Apply the five shared features (Event 2)}
\begin{longtable}{@{}>{\raggedright\arraybackslash}p{0.19\textwidth} >{\raggedright\arraybackslash}p{0.38\textwidth} >{\raggedright\arraybackslash}p{0.37\textwidth}@{}}
\toprule
\textbf{Feature} & \textbf{Exact wording + timestamp} & \textbf{What it does in context} \\
\midrule
\endhead
\textbf{1. Naming} & Title: \texttt{"We the people of India..."} \newline  Clip \texttt{[06:06]}: \texttt{"diz aar not students... diz aar CJP end CJP sponsard gundaj"} $\rightarrow$ Cut to \texttt{[06:14]}: \texttt{"meri beti hai do betiyan hain... ek 10tha mein ek phorth mein hai... bachche sabke hain"} & Frames protesters under the supreme constitutional self-designation (\texttt{"We the people of India..."}), and uses montage juxtaposition so the state/TV label \texttt{"gundaj"} [goons] collides head-on with a soft-spoken mother of two schoolgirls. \\
\addlinespace[4pt]
\textbf{2. Presupposition} & \texttt{"toh ye sarkar itna bachchon se darti kyon hai?"} \texttt{[04:30]} \newline  \texttt{"agar wo preshar mein nahin hai. agar wo shem mein nahin hai. agar wo gilt mein nahin hai toh kyon itne kaimre hain?"} \texttt{[10:17--10:27]} & The \textit{why}-question (\texttt{darti kyon hai?}) presupposes that the government is afraid of young students; Supriya's conditional rhetorical questions (\texttt{kyon itne kaimre hain?}) presuppose state guilt and shame behind heavy surveillance. \\
\addlinespace[4pt]
\textbf{3. Agency} & \texttt{"lathi blos shud bi emd et sopht porashan of da bodi end kantekt vith da hed aur kolar bon shud bi avoyded... yahan hamne dekha ke ladkon ke toh pair toot hi rahe hain... ledi ka kandha phaad diya"} \texttt{[02:27--02:51]} & Contrasts the deontic passive of official police restraint rules (\texttt{shud bi avoyded}) against active ergonomic/bodily destruction by unnamed officers (\texttt{unhonne chaar dande mere pair pe mare}, \texttt{[01:47]}; \texttt{kandha phaad diya}, \texttt{[02:51]}; \texttt{usne nem plet nikal rakhi thi}, \texttt{[01:38]}). \\
\addlinespace[4pt]
\textbf{4. Antagonisation} & \texttt{"hamesha apne man ki baat karke aap chale jaoge... jab bachche bolna chahte hain aap unko nahin sun rahe... aap Hindu muslim Hindu muslim karke tab tak ladvaoge... yahan... na koi Hindu bol raha tha na koi musalman sab ek saath milke baithe hue the"} \texttt{[05:14--05:57]} & Pits a monologic, communally divisive political class (\texttt{apne man ki baat}, \texttt{Hindu muslim... aap log karvate ho}) against an inclusive, inter-faith civic solidarity of youth (\texttt{sab ek saath milke baithe hue the}). \\
\addlinespace[4pt]
\textbf{5. Intensification / softening} & \textbf{Understated documentary irony:} Opening 1997 story (\texttt{"unke dost unko mukkoo bhaee bulate hain. hu ij dis gaay? mistar Dharmendra Pradhan"}, \texttt{[00:35--00:44]}) and closing hard-cut from \texttt{"wo jhukegi, sarkar jhukegi"} \texttt{[10:56]} to TV anchor screaming \texttt{"brootal vaylens kairid aaut baay da [goons] oph da CJP"} \texttt{[10:56]} & Instead of anchor shouting, Mohak uses \textbf{calm, understated narrative restraint} paired with jarring audio montage---making the state/TV hyperbole (\texttt{tukde-tukde}, \texttt{dahashat gard}, \texttt{kokroch}) sound absurd and cruel against quiet, injured citizens. \\
\bottomrule
\end{longtable}

\subsection*{Short summary (Event 2)}
\begin{itemize}[leftmargin=1.5em, itemsep=3pt, topsep=4pt]
  \item \textbf{What interpretation of the event does this video encourage?} That the 20 July Jantar Mantar protest was a deeply legitimate, peaceful, pan-Indian constitutional gathering (\texttt{"We the people of India..."}) of wronged NEET aspirants, anxious parents, and educators, brutally assaulted in violation of police manuals by a government led by an Education Minister who had himself been lathi-charged protesting a paper leak in 1997 (\texttt{[00:00--01:03]}, \texttt{[02:06--02:59]}), while TV channels slandered ordinary mothers and teenagers as \texttt{"gundaj"} (\texttt{[06:06--06:34]}).
  \item \textbf{Whose perspective is foregrounded, challenged, or missing?} \textit{Foregrounded:} First-person voices of injured students, repeat NEET candidates, out-of-state family members (Hyderabad, Konkan), mothers, and schoolteachers (\texttt{[01:12--10:56]}). \textit{Challenged:} Dharmendra Pradhan's moral consistency (\texttt{[00:00--01:03]}, \texttt{[05:04]}), police claims of lawful minimum force (\texttt{[02:06--02:59]}), the \texttt{"cockroach"} slur against unemployed youth (\texttt{[02:59]}), and TV "goon/terrorist" labels (\texttt{[00:44]}, \texttt{[01:03]}, \texttt{[06:06]}, \texttt{[10:56]}). \textit{Missing:} Any interview with Delhi Police officers injured at barricades or detailed discussion of CJP's organizational leadership/political visitors.
  \item \textbf{What evidence supports your interpretation? What complicates it?} \textit{Supports:} Archival 1997 photograph/biography of Dharmendra Pradhan (\texttt{[00:00--00:44]}), verbatim Indian Police crowd-dispersal guidelines (\texttt{[02:06--02:39]}), concrete NEET/CET grievance details (\texttt{2024} \& \texttt{2026} leaks at \texttt{[03:37]}, \texttt{630} normalized to \texttt{530} at \texttt{[08:54]}), and raw on-site interviews with visibly shaken/injured citizens (\texttt{[01:12--10:56]}). \textit{Complicates:} By curating exclusively peaceful, articulate protesters and omitting scenes of barricade-breaking or stone-pelting that injured police officers (documented in Aaj Tak \texttt{JXLIs0WtJAM} \texttt{[18:09--23:14]}), the documentary constructs a pure moral contrast that leaves out the messy street dynamics at Parliament Street barricades.
\end{itemize}
\vspace{4pt}\noindent\rule{\textwidth}{0.4pt}\vspace{4pt}

\section*{Common form: Blank excerpt record}
\textit{Copy or print this page whenever you find a passage worth analysing. Keep the exact wording and enough context to interpret it.}


\subsection*{Excerpt Record 1 (Event 2 --- Historical Irony Opening on Dharmendra Pradhan)}
\begin{itemize}[leftmargin=1.5em, itemsep=3pt, topsep=4pt]
  \item \textbf{Outlet:} Mohak Mangal | \textbf{Date:} \texttt{2026-07-22}
  \item \textbf{Video title:} \textit{"We the people of India..."}
  \item \textbf{Timestamp(s):} \texttt{[00:00--01:03]} | \textbf{Analyst:} Baani Solanki | \textbf{Language:} Hindi / Hinglish \& English (Roman script)
  \item \textbf{Exact transcript excerpt:}
    \begin{quote}\small\ttfamily [00:00] aaee vant too start dis video baay teling a short stori. yeh kahani hai 1997 ki. is aadmi ki. 1997 mein udisa mein ek paper leak [00:10] huaa tha. 1500 students protest karne gae... protest pisaphul thi [00:18] par police ne lathi charj kar diya. is aadmi jiski mainne aapko photo dikhaee unko bahut pita gaya. phraikchar huaa... [00:35] hi vaas da neshanal sekretri oph ebivipi. unke dost unko mukkoo bhaee bulate hain. hu ij dis gaay? mistar Dharmendra Pradhan. India ke [00:44] ejukeshan ministar. [Clip: "jo log desh ko tukde-tukde karna chahte hain... ye toh dahashat gardon ka bi tim hai."] [00:53] aap aayaran hi dekho. dait da sem gaay jo 1997 mein paper leak ke khilaph protest kar rahe the. aaj unke khilaph protest ho rahi hai aur [01:03] unki sarkar ab students par lathi charj kar rahi hai.\end{quote}
  \item \textbf{Feature(s):} $\boxtimes$ Naming $\boxtimes$ Presupposition $\boxtimes$ Agency $\boxtimes$ Antagonisation $\boxtimes$ Intensification / softening
  \item \textbf{Who is speaking? Is this the outlet's wording or a quotation?} Mohak Mangal's scripted opening narration intercut at \texttt{[00:44--00:53]} with a verbatim ruling-party political soundbite smearing Jantar Mantar protesters.
  \item \textbf{What does the wording imply, foreground, soften, or leave unstated?} Deliberately withholds the protagonist's identity (\texttt{"is aadmi ki"}, \texttt{"mukkoo bhaee"}) for 40 seconds so even pro-BJP viewers naturally sympathize with a peaceful student beaten by police during a paper-leak march---before revealing \texttt{"mistar Dharmendra Pradhan"} (\texttt{[00:35--00:44]}). This narrative trap makes delegitimizing 2026 NEET protesters logically impossible without condemning Pradhan's own ABVP past.
  \item \textbf{What context or evidence supports or complicates this interpretation?} Supported by Dharmendra Pradhan's documented 1990s ABVP student-politics career in Odisha. Complicated only by the administrative distinction between being an opposition student leader and bearing executive cabinet responsibility for national security around Parliament.
  \item \textbf{How does another outlet describe the same event or actor?} \textbf{The Wire} (\texttt{pYaTVvHeGoM}, \texttt{[12:42--14:02]}) attacks Dharmendra Pradhan via Naseeruddin Shah calling him \texttt{"this idiot"} guilty of Orwellian \texttt{"classic double speak"}, whereas \textbf{Mohak Mangal} avoids personal insults entirely, letting biographical irony (\texttt{"aap aayaran hi dekho"}, \texttt{[00:53]}) do far more persuasive work.
  \item \textbf{My interpretation, with reasons:} This is arguably the most rhetorically effective opening across all 11 videos: by delaying Naming (\texttt{is aadmi} $\rightarrow$ \texttt{Dharmendra Pradhan}), Mohak disarms partisan defensiveness before introducing a single 2026 clip.
\end{itemize}

\subsection*{Excerpt Record 2 (Event 2 --- Police Minimum-Force Guidelines vs Ground Injuries)}
\begin{itemize}[leftmargin=1.5em, itemsep=3pt, topsep=4pt]
  \item \textbf{Outlet:} Mohak Mangal | \textbf{Date:} \texttt{2026-07-22}
  \item \textbf{Video title:} \textit{"We the people of India..."}
  \item \textbf{Timestamp(s):} \texttt{[01:38--02:59]} | \textbf{Analyst:} Baani Solanki | \textbf{Language:} Hindi / Hinglish \& English (Roman script)
  \item \textbf{Exact transcript excerpt:}
    \begin{quote}\small\ttfamily [01:38] [Injured youth:] sar ek toh thri star vala tha sar usne nem plet nikal rakhi thi... main niche gir gaya toh unhonne chaar dande mere pair pe mare... ek ladki aaee mere [01:57] samne ki nahin mat maro gir gaya wo usko yahan pe mara unhonne aantiyon ko nahin chhoda ek 10tha klas ka bachcha tha usko mara hai... [02:06] [Mohak Mangal:] indiyan sarkar ki na kuch gaidlains hain. aaee jast vant too rid dem aaut too yoo. da eksatent of phors [02:16] yoojd mast bi sabjekt too da prinsipal of minimam yooz of phors... [02:27] lathi blos shud bi emd et sopht porashan of da bodi end kantekt vith da hed aur kolar [02:39] bon shud bi avoyded ez phaar es praiktikal. yahan hamne dekha ke ladkon ke toh pair toot hi rahe hain. [02:51] us aaee dont no kya aij thi uski 20 saal ki umra hogi uska kandha phaad diya batao ledi ka kandha phaad diya buri tarike se nange pair [02:59] behvas bhaag rahe hain log kya tarika hai\end{quote}
  \item \textbf{Feature(s):} $\boxtimes$ Naming $\square$ Presupposition $\boxtimes$ Agency $\square$ Antagonisation $\boxtimes$ Intensification / softening
  \item \textbf{Who is speaking? Is this the outlet's wording or a quotation?} First-person testimony of an injured protester (\texttt{[01:38--02:06]}), followed by Mohak Mangal reading verbatim from official Government of India policing guidelines (\texttt{[02:06--02:39]}) and adding his own observation (\texttt{[02:39--02:59]}).
  \item \textbf{What does the wording imply, foreground, soften, or leave unstated?} Instead of using hyperbolic labels like \texttt{"General Dyer"} or \texttt{"terrorists"} (as Dhruv Rathee does in \texttt{l-6Pm-112io} \texttt{[10:08]}, \texttt{[11:54]}), Mohak evaluates the state \textbf{by its own rulebook} (\texttt{"indiyan sarkar ki na kuch gaidlains hain"}). Contrasting bureaucratic legal English (\texttt{"kantekt vith da hed aur kolar bon shud bi avoyded"}) with raw Hindi bodily trauma (\texttt{"ledi ka kandha phaad diya"}, \texttt{"nange pair behvas bhaag rahe hain"}) foregrounds institutional lawlessness.
  \item \textbf{What context or evidence supports or complicates this interpretation?} Corroborated by hospital/surgical reports of head stitches, fractures, and neck pellets cited on \textbf{Aaj Tak} (\texttt{JXLIs0WtJAM}, \texttt{[06:30]}) and \textbf{The Wire} (\texttt{pYaTVvHeGoM}, \texttt{[10:07--10:38]}).
  \item \textbf{How does another outlet describe the same event or actor?} \textbf{Aaj Tak} (\texttt{JXLIs0WtJAM}, \texttt{[19:09]}, \texttt{[21:49--23:50]}) frames police action as \texttt{"kanoon vyavastha ko durust karna ye police ka kaam hai"} and interviews injured Addl. DCPs explaining they were merely persuading crowds at an \texttt{"illigal"} march.
  \item \textbf{My interpretation, with reasons:} Reading statutory guidelines aloud (\texttt{"aaee jast vant too rid dem aaut too yoo"}) softens the presenter's tone into calm legal objectivity, which paradoxically intensifies the indictment of police brutality.
\end{itemize}

\subsection*{Excerpt Record 3 (Event 2 --- Juxtaposing "CJP Goons" Clip Against Mother \& Teacher)}
\begin{itemize}[leftmargin=1.5em, itemsep=3pt, topsep=4pt]
  \item \textbf{Outlet:} Mohak Mangal | \textbf{Date:} \texttt{2026-07-22}
  \item \textbf{Video title:} \textit{"We the people of India..."}
  \item \textbf{Timestamp(s):} \texttt{[06:06--07:40]} | \textbf{Analyst:} Baani Solanki | \textbf{Language:} Hindi / Hinglish \& English (Roman script)
  \item \textbf{Exact transcript excerpt:}
    \begin{quote}\small\ttfamily [06:06] [Clip:] diz aar not students. de aar not NEET espirents. diz aar CJP end CJP sponsard gundaj. [06:14] [Mother on ground:] vaise toh meri beti hai do betiyan hain. wo padh rahi hain. ek 10tha mein ek phorth mein hai... main [06:24] unke lie aaee hoon. ki jo mudda hai wo sabka hai. kisi ek ka nahin hai. bachche sabke hain. ejukeshan sabki hai [06:34] aur jo paper leak ho rahe hain usse jo bachchon ne susaid kiya hai wo hona nahin chahie... [07:14] [Schoolteacher:] main ej a tichar demokresi chaiptar padha rahe ho. so vhat aaee tich student? toh wo rilet hi nahin kar pa rahe hain. bikoz de phil ki demokresi hai hi nahin... [07:32] sam of maay student aalso thot ki India mein nahin rahna de vant too stadi in [07:40] abrod. so maay kveshchan iz iz ki agar sabhi aise jane lagenge toh phir vhaee India?\end{quote}
  \item \textbf{Feature(s):} $\boxtimes$ Naming $\boxtimes$ Presupposition $\boxtimes$ Agency $\boxtimes$ Antagonisation $\square$ Intensification / softening
  \item \textbf{Who is speaking? Is this the outlet's wording or a quotation?} Editorial montage cutting directly from an external English soundbite (\texttt{[06:06]}) to two ordinary female citizens interviewed at Jantar Mantar---a mother (\texttt{[06:14--07:08]}) and a civics teacher (\texttt{[07:14--07:58]}).
  \item \textbf{What does the wording imply, foreground, soften, or leave unstated?} Directly answers the delegitimizing claim \texttt{"diz aar not students"} (which presupposes that non-NEET-aspirants at the protest must be \texttt{"sponsard gundaj"}): Mohak shows that non-candidates at Jantar Mantar are actually \textit{mothers of 10th-graders and schoolteachers} exercising collective civic responsibility (\texttt{"bachche sabke hain. ejukeshan sabki hai"}).
  \item \textbf{What context or evidence supports or complicates this interpretation?} Directly refutes \textbf{Republic TV} (\texttt{HTvb6E\_1Qdc}, \texttt{[04:58--05:26]}: \texttt{"these were not students... the students had evaporated"}) and Maharashtra CM Devendra Fadnavis on \textbf{Aaj Tak} (\texttt{JXLIs0WtJAM}, \texttt{[10:13--10:23]}: claiming attendees who aren't NEET examinees are pushing \texttt{"apna-apna ejenda"}).
  \item \textbf{How does another outlet describe the same event or actor?} While \textbf{Republic TV} treats the presence of non-examinees as proof of \texttt{"Deep State"} / \texttt{"goon"} infiltration (\texttt{HTvb6E\_1Qdc}, \texttt{[00:23]}, \texttt{[04:19]}), \textbf{Mohak Mangal} and \textbf{Sarthak Goswami} (\texttt{EM1GeZufO6E}, \texttt{[07:58--08:16]}) frame broad citizen participation as the highest proof of civic solidarity.
  \item \textbf{My interpretation, with reasons:} This edit demonstrates how \textit{montage structure itself} performs discourse analysis---dismantling elite political Naming (\texttt{"CJP sponsard gundaj"}) solely through ordinary citizens' self-presentation.
\end{itemize}

\subsection*{Excerpt Record 4 (Event 2 --- Konkan Citizen Supriya \& Closing Cut)}
\begin{itemize}[leftmargin=1.5em, itemsep=3pt, topsep=4pt]
  \item \textbf{Outlet:} Mohak Mangal | \textbf{Date:} \texttt{2026-07-22}
  \item \textbf{Video title:} \textit{"We the people of India..."}
  \item \textbf{Timestamp(s):} \texttt{[09:46--10:56]} | \textbf{Analyst:} Baani Solanki | \textbf{Language:} Hindi / Hinglish \& English (Roman script)
  \item \textbf{Exact transcript excerpt:}
    \begin{quote}\small\ttfamily [09:46] [Supriya from Konkan:] jab lathi charj ho gaya hamne bola ki nahin ham udhar hi khade rahenge agar marna hai toh maro kyunki agar satya ka shabd bolte hain toh [09:56] bahut takat honi chahie usko nibhane ke lie aur ek baat ki agar aap jis darvaje se unko bhagana chahte hain wo darvaja kam se kam khula toh rahna chahie aise band karke marna toh bahut [10:06] hi inahyooman hai... [10:27] ant mein jit usi ki hai jisne satya aur ahinsa ka marg [10:36] pakda hai... [10:56] wo jhukegi, sarkar jhukegi. [Hard cut to Arnab Goswami clip:] deyar haiz bin brootal vaylens kairid aaut baay da gsa [goons] oph da CJP.\end{quote}
  \item \textbf{Feature(s):} $\boxtimes$ Naming $\boxtimes$ Presupposition $\boxtimes$ Agency $\boxtimes$ Antagonisation $\boxtimes$ Intensification / softening
  \item \textbf{Who is speaking? Is this the outlet's wording or a quotation?} Citizen protester Supriya from Konkan (\texttt{[09:29--10:56]}), followed by Mohak's final 3-second audio cut to a mainstream English TV anchor (\texttt{[10:56]}).
  \item \textbf{What does the wording imply, foreground, soften, or leave unstated?} Supriya invokes Gandhian moral vocabulary (\texttt{"satya aur ahinsa ka marg"}, \texttt{[10:27--10:36]}) and exposes the tactical cruelty of police closing exit doors while beating crowds (\texttt{"aise band karke marna toh bahut hi inahyooman hai"}, \texttt{[09:56--10:06]}). Ending the video on the abrupt smash-cut from her calm prophecy (\texttt{"sarkar jhukegi"} --- fulfilled 3 days later when Pradhan resigned!) to \texttt{"brootal vaylens kairid aaut baay da [goons] oph da CJP"} indicts mainstream media's disconnect from reality without Mohak uttering a single closing word.
  \item \textbf{What context or evidence supports or complicates this interpretation?} Trapping protesters by shutting metro stations/gates while firing tear gas and lathis is corroborated by \textbf{Dhruv Rathee} (\texttt{l-6Pm-112io}, \texttt{[06:43--09:38]}).
  \item \textbf{How does another outlet describe the same event or actor?} The very voice sampled at \texttt{[01:03]} and \texttt{[10:56]}---\textbf{Republic TV's Arnab Goswami} (\texttt{HTvb6E\_1Qdc}, \texttt{[04:19--05:26]})---portrays the exact same Jantar Mantar crowd as \texttt{"drunken violence"}, \texttt{"sexual perversion"}, and \texttt{"lynch mob of Abhijit Dipke"}.
  \item \textbf{My interpretation, with reasons:} Ending on an adversary's quote (\texttt{[10:56]}) after 10 minutes of humanizing ground testimony turns the TV anchor's own words into self-incriminating evidence of media propaganda.
\end{itemize}
\vspace{4pt}\noindent\rule{\textwidth}{0.4pt}\vspace{4pt}

\section*{Common form: End-of-event summary}

\subsection*{End-of-Event Summary --- Event 1 (Ram Mandir Donation Controversy)}
\begin{itemize}[leftmargin=1.5em, itemsep=3pt, topsep=4pt]
  \item \textbf{Status:} \textbf{Skipped per project instructions} (\texttt{lns\_project\_instructions.pdf} specifies \textit{Skip} for Mohak Mangal under Event 1, as he did not publish a video on the Ram Mandir donation controversy).
\end{itemize}

\subsection*{End-of-Event Summary --- Event 2 (Students' Protest Regarding NEET Paper Leak)}
\begin{itemize}[leftmargin=1.5em, itemsep=3pt, topsep=4pt]
  \item \textbf{The main way this outlet framed the event was:} As a dignified, multi-generational constitutional struggle of \texttt{"We the people of India..."}---students, mothers, workers, and teachers peacefully demanding exam integrity and democratic accountability---who were subjected to unlawful, guideline-violating police lathi-charges by an Education Minister (Dharmendra Pradhan) ironically guilty of forgetting his own 1997 paper-leak protest injuries, and smeared as \texttt{"kokroch"} and \texttt{"gundaj"} by an out-of-touch establishment and TV media (\texttt{[00:00--02:59]}, \texttt{[06:06--10:56]}).
  \item \textbf{The strongest linguistic evidence was (quote + timestamp):} \texttt{"dait da sem gaay jo 1997 mein paper leak ke khilaph protest kar rahe the. aaj unke khilaph protest ho rahi hai aur unki sarkar ab students par lathi charj kar rahi hai."} (\texttt{[00:53--01:03]}), paired with the smash-cut from \texttt{"diz aar CJP sponsard gundaj"} (\texttt{[06:06]}) to \texttt{"bachche sabke hain. ejukeshan sabki hai"} (\texttt{[06:24]}).
  \item \textbf{The shared feature(s) most clearly illustrated were:} \textbf{Naming} (\texttt{"We the people of India..."} \& \texttt{"bachche sabke hain"} vs \texttt{"tukde-tukde"}, \texttt{"kokroch"}, \texttt{"CJP sponsard gundaj"}), \textbf{Agency} (measuring unbadged officers' head/shoulder/leg blows against official Government of India minimum-force guidelines at \texttt{[01:38--02:59]}), and \textbf{Intensification / Softening via Documentary Restraint} (letting quiet citizen voices and biographical irony expose hyperbolic state/media smears).
  \item \textbf{An alternative interpretation of the same evidence is:} That while individual NEET aspirants, parents, and citizens interviewed on 20--21 July were sincere and peaceful, a massive unauthorized march toward Parliament during an active session inevitably drew in unruly/political elements who clashed with police at barricades (injuring \texttt{110+} officers per Aaj Tak \texttt{JXLIs0WtJAM}, \texttt{[19:43]}), which Mohak's selective ground curation leaves off-screen.
  \item \textbf{What I cannot conclude from this transcript alone:} Whether any protesters at other barricade points pelted stones or injured police officers, how CJP's core leadership coordinated negotiations with Union Ministers, or what happened on 25 July when Dharmendra Pradhan resigned (since this video was published earlier on \texttt{2026-07-22}).
\end{itemize}
\textit{Research reminder: interpret wording in context, distinguish allegations from established facts, and support claims with exact excerpts. Mixed or contradictory evidence is worth recording.}


\clearpage

\begin{center}
    {\large\bfseries LANGUAGE \& SOCIETY PROJECT}\\[3pt]
    {\LARGE\bfseries Transcript Analysis Worksheet}\\[3pt]
    {\normalsize Qualitative analysis of political news discourse}
\end{center}
\vspace{2pt}
\noindent\rule{\textwidth}{0.4pt}
\vspace{4pt}


\section*{Creator 4: Aaj Tak (\textit{Halla Bol} \& \textit{AT2 Live} --- Mainstream Hindi Television News)}
\begin{itemize}[leftmargin=1.5em, itemsep=3pt, topsep=4pt]
  \item \textbf{Analyst:} Baani Solanki (\texttt{2025114016})
  \item \textbf{Channel / Outlet:} Aaj Tak (\texttt{@aajtak}, TV Today Network)
  \item \textbf{Ideological / Editorial Positioning:} Mainstream commercial Hindi broadcast news; combines high-drama devotional/emotional vocabulary (\texttt{"toophan"}, \texttt{"paap"}, \texttt{"dangal"}) with formal procedural/law-and-order deference toward state institutions (SIT, CM Yogi, PM Modi, Delhi Police), while staging multi-party partisan shouting matches and electoral horse-race framing.
  \item \textbf{Transcript Files Referenced:}
  \begin{itemize}[leftmargin=1.4em, itemsep=2pt, topsep=2pt]
    \item Event 1: \texttt{transcripts/event1\_ram\_mandir/04\_aaj\_tak\_B4L1pz0o4dc.txt}
    \item Event 2: \texttt{transcripts/event2\_neet\_protest/04\_aaj\_tak\_JXLIs0WtJAM.txt}
  \end{itemize}
\end{itemize}
\vspace{4pt}\noindent\rule{\textwidth}{0.4pt}\vspace{4pt}

\section*{Guidelines: Five features to look for in every event}
\begin{longtable}{@{}>{\raggedright\arraybackslash}p{0.17\textwidth} >{\raggedright\arraybackslash}p{0.20\textwidth} >{\raggedright\arraybackslash}p{0.24\textwidth} >{\raggedright\arraybackslash}p{0.32\textwidth}@{}}
\toprule
\textbf{Feature} & \textbf{What to look for} & \textbf{Illustrative example} & \textbf{Questions to answer} \\
\midrule
\endhead
\textbf{1. Naming} & What are people, groups, actions, or events called? & "Students", "protesters", "agitators"; "irregularity", "scam", "controversy". & What exact label is used? Who or what does it refer to? Who uses it? What alternative label could have been used? \\
\addlinespace[4pt]
\textbf{2. Presupposition} & What is treated as already true or taken for granted? & "Why are students creating chaos again?" may presuppose earlier disruption. & What does the wording assume? Is that assumption supported, disputed, or left unexamined? Does it appear in a question or words such as "again", "still", or "finally"? \\
\addlinespace[4pt]
\textbf{3. Agency} & Who is shown doing something? Who is acted upon? Who disappears? & "Police injured protesters" vs. "Protesters were injured." & Who is the actor? Who experiences the action? Is responsibility named, backgrounded, or omitted? \\
\addlinespace[4pt]
\textbf{4. Antagonisation} & Who is positioned against whom? & "Ordinary citizens" vs. "political troublemakers". & Which groups are contrasted? Is one presented as more legitimate, moral, patriotic, dangerous, or disloyal? Is an "us vs. them" division implied? \\
\addlinespace[4pt]
\textbf{5. Intensification / softening} & What is exaggerated, minimised, softened, or euphemised? & "Violent mob" vs. "peaceful demonstrators"; "incident" instead of a direct description. & Which words heighten emotion or seriousness? Which minimise it? Are metaphors, dramatic adjectives, hedges, or euphemisms used? Does the wording match the evidence? \\
\bottomrule
\end{longtable}
\textit{Reminder: A label or implication is not automatically false or manipulative. Examine the surrounding context and evidence.}

\vspace{4pt}\noindent\rule{\textwidth}{0.4pt}\vspace{4pt}

\section*{Event1. Ram Mandir donation controversy}
\textbf{Focus:} How is the controversy framed---as financial accountability, religious sentiment, political motivation, or something else? Be precise about the specific allegation or episode covered.

\begin{itemize}[leftmargin=1.5em, itemsep=3pt, topsep=4pt]
  \item \textbf{Outlet:} Aaj Tak (\textit{Halla Bol With Anjana Om Kashyap})
  \item \textbf{Date:} \texttt{2026-06-19} \textit{(Day 5 of the UP SIT probe --- one week BEFORE the 26 June arrests and Champat Rai's resignation)}
  \item \textbf{Video title:} \textit{Halla Bol With Anjana Om Kashyap Full Show: Ram Mandir ko kisne 'loota'? | Ram Mandir Donation | BJP} (\url{https://youtu.be/B4L1pz0o4dc?si=BRvYR4hn3KalxM_t})
  \item \textbf{Timestamp(s):} Relevant Event 1 segment \texttt{[00:00--38:26]} (\texttt{100\%} of broadcast dedicated to the Ram Mandir donation controversy; total duration \texttt{38:30}).
  \item \textbf{Analyst:} Baani Solanki (\texttt{2025114016})
  \item \textbf{Language:} Hindi / Hinglish (Native Hindi broadcast audio, Romanized script; non-AI-dubbed)
\end{itemize}

\subsection*{Event-specific questions (Event 1)}
\begin{enumerate}[leftmargin=1.6em, itemsep=4pt, topsep=4pt]
  \item \textbf{What terms are used: "scam", "alleged irregularity", "corruption", "political controversy", or something else?}
  \begin{itemize}[leftmargin=1.4em, itemsep=2pt, topsep=2pt]
    \item \textbf{Carefully hedged in the anchor's voice, dramatized in the VTR package, and polarized across panelists:}
    \item \textbf{Title \& Anchor (Anjana Om Kashyap):} Uses scare-quotes in the title (\texttt{Ram Mandir ko kisne 'loota'?}) and attributes allegations carefully as \texttt{"Ram Mandir chadhava chori ka aarop"} [allegation of offering theft] (\texttt{[00:10]}), \texttt{"gambhir aarop"} [serious allegations] (\texttt{[00:48]}, \texttt{[09:12]}), and \texttt{"bahut bada toophan"} [a massive storm] (\texttt{[00:10]}, \texttt{[00:58]}).
    \item \textbf{Aaj Tak Bureau Report (\texttt{[02:02--09:02]}):} Escalates into devotional-moral terms---\texttt{"bhrashtachariyon ne dhava bola hai"} [corrupt people have raided the temple] (\texttt{[02:02]}), \texttt{"karodon arbon bhakton ki aastha se chhal"} [betrayal of crores of devotees' faith] (\texttt{[02:19]}), \texttt{"chadhava chori kand"} (\texttt{[02:19--02:30]}), and \texttt{"ek aisa paap hai jiska prayashchit shayad nahin ho sakta"} [a sin perhaps beyond atonement] (\texttt{[08:51]}).
    \item \textbf{Panelists:} SP/Congress/analysts call it \texttt{"bahut bada ghotala dakaiti chori"} (\texttt{[08:14]}) and \texttt{"jamin ghotala"} (\texttt{[34:58]}), while pro-Trust/VHP voices soften it to \texttt{"gadabdi"} [irregularity] (\texttt{[10:08]}), \texttt{"kuch avyavastha huee hai"} [some mismanagement occurred] (\texttt{[34:21]}), and \texttt{"badant intjami"} [poor administration] (\texttt{[22:38]}).
  \end{itemize}
  \item \textbf{How are the temple, trust, donors, and people making allegations described?}
  \begin{itemize}[leftmargin=1.4em, itemsep=2pt, topsep=2pt]
    \item \textbf{Temple:} Reverentially named \texttt{"prabhu shri Ram ke Mandir"} (\texttt{[00:00]}) and \texttt{"500 salon ke sangharsh ke baad yeh Mandir bana"} (\texttt{[01:16]}, \texttt{[02:02]}).
    \item \textbf{Trust \& Suspects:} Anchor Anjana notes \texttt{"Mandir prabandhan samiti se jude log hi sandeh ke ghere mein hain"} (\texttt{[00:20]}) and repeatedly cites Construction Chair Nripendra Misra's admission that oversight wasn't foolproof (\texttt{[01:06]}, \texttt{[04:24--06:05]}, \texttt{[16:44--17:08]}, \texttt{[33:19--33:47]}), while pointedly \textit{exonerating} Misra himself (\texttt{"nirman alag hai, prabandhan alag hai"}, \texttt{[10:57--11:24]}). Suspects (\texttt{Ramashankar Yadav urph Tinnoo Yadav}, \texttt{Lavkush Mishra}, \texttt{Anukalp Mishra}, \texttt{Rajesh Pathak}, \texttt{K.D. Tiwari}, \texttt{Manish Yadav}, \texttt{[03:39--04:06]}) are framed as counting/admin staff, while right-leaning panelist Prof. Sangeet Ragi shields Champat Rai as a lifelong RSS pracharak living in \texttt{"do dhoti aur do kurte"} (\texttt{[12:25--12:56]}). Notably, Aaj Tak notes at \texttt{[13:21--13:30]} (via SP's Anurag Bhadauria) that no official BJP spokesperson even showed up, leaving VHP spokesperson Pravesh Chaudhary (\texttt{[01:38]}) and sympathetic analysts (\texttt{Sangeet Ragi}, \texttt{Saurabh/Soor Maithil}) to defend the Trust.
    \item \textbf{Donors:} Revered as \texttt{"karodon arbon bhakton"} (\texttt{[02:19]}) who made \texttt{"samarpan"} [sacred offering, not mere chanda] (\texttt{[17:17--17:25]}).
    \item \textbf{People Making Allegations:} Introduced through an electoral-partisan lens (\texttt{"vipaksh bhi is mudde ko bada bana raha hai aur BJP ko is mauke par gher raha hai"}, \texttt{[00:20--00:29]}), citing SP ex-minister Pawan Pandey, Akhilesh Yadav's June 7 X post (\texttt{[06:13--06:21]}), and Faizabad MP Awadhesh Prasad (\texttt{[01:06--01:16]}, \texttt{[06:32--07:46]}), whom the Aaj Tak reporter pointedly quizzes on why Akhilesh Yadav hasn't visited Ayodhya Ram Mandir (\texttt{[07:10--07:46]}, echoed by Anjana at \texttt{[16:10--16:20]}, \texttt{[26:52--27:10]}).
  \end{itemize}
  \item \textbf{Is the issue framed as financial accountability, an attack on Hindu faith, a political dispute, or a matter of public trust?}
  \begin{itemize}[leftmargin=1.4em, itemsep=2pt, topsep=2pt]
    \item \textbf{Hybrid devotional-trust + partisan slugfest frame:} Aaj Tak frames the theft itself as a grave wound to devotees' faith and public trust (\texttt{" kya Yogi aadityanath yeh bharosa Ram bhakton ko lauta paenge"}, \texttt{[00:38--00:48]}; \texttt{"paap"}, \texttt{[08:51]}) requiring managerial accountability (\texttt{[16:44--17:17]}, \texttt{[33:19--33:47]}), while simultaneously allowing over half the show (\texttt{[16:10--23:18]}, \texttt{[26:16--29:43]}, \texttt{[35:46--38:00]}) to devolve into a BJP/VHP vs SP/Congress litmus test over 1990 kar-sevak firing and Pran Pratishtha boycotts.
  \end{itemize}
  \item \textbf{Are questions about donations presented as legitimate scrutiny or as an attempt to malign the temple or its supporters?}
  \begin{itemize}[leftmargin=1.4em, itemsep=2pt, topsep=2pt]
    \item \textbf{Split between Anchor vs Right-Wing Panelists/Audience:} Anchor Anjana insists scrutiny of unchecked VIP/trustee access without frisking is legitimate (\texttt{"akauntebal sabko hona padega... sare niyam toh shraddhaluon ke lie hain"}, \texttt{[33:19--33:47]}), whereas CM Yogi (\texttt{[00:29--00:38]}, \texttt{[02:30--03:13]}), VHP's Pravesh Chaudhary (\texttt{[18:50--19:19]}, \texttt{[34:13]}: \texttt{"nyas par prashn khada karna hi invailid hai"}), Soor Maithil (\texttt{[21:13--21:39]}), and studio audience youth (\texttt{[26:33]}, \texttt{[28:08]}: \texttt{"aapko hak hi nahin hai... Ram ji ke bare mein bolne ke lie"}) attack opposition questioners as \texttt{"ramadrohi"} [traitors to Ram] (\texttt{[18:50]}, \texttt{[21:30]}).
  \end{itemize}
  \item \textbf{Are allegations attributed to specific people and supported with evidence, or presented as established facts?}
  \begin{itemize}[leftmargin=1.4em, itemsep=2pt, topsep=2pt]
    \item Carefully attributed at the editorial level: Aaj Tak cites \textbf{exclusive CCTV footage} of counting (\texttt{[03:30]}), \textbf{Rs.\,10 lakh cash seized} from Lavkush Mishra's house (\texttt{[03:49]}), and \textbf{Nripendra Misra's on-record interview} with Anjana admitting lack of frisking, SBI's \texttt{"most non siriyas mainar"}, and broken double-lock protocols (\texttt{[01:06]}, \texttt{[04:24--06:05]}), while repeatedly reminding viewers that \texttt{"abhi toh SIT ki poori report aani baki hai"} (\texttt{[00:58]}, \texttt{[38:09]}).
  \end{itemize}
  \item \textbf{How are denials and counterarguments represented: explained, dismissed, challenged, or omitted?}
  \begin{itemize}[leftmargin=1.4em, itemsep=2pt, topsep=2pt]
    \item \textbf{Given massive airtime, and occasionally challenged on procedure:} CM Yogi's \texttt{"doodh ka doodh aur pani ka pani"} SIT assurance is aired prominently twice (\texttt{[00:29--00:38]}, \texttt{[02:30--03:30]}); VHP's \texttt{"500 varsh intjar kiya hai. kya ham 15 din aur intjar nahin kar sakte?"} (\texttt{[18:00--18:38]}) and Sangeet Ragi's separation of senior trustees from 40 counting clerks (\texttt{[09:59--10:36]}) are heard at length. However, Anjana and analyst Satish Prakash do challenge VHP on \textbf{why no FIR was registered immediately} (\texttt{"Anjana ji poochhie FIR kyon nahin karaee?"}, \texttt{[22:09--22:28]}; \texttt{[23:18--23:29]}) and why trustees walked in above police frisking rules (\texttt{[16:52--17:17]}, \texttt{[33:19--33:47]}).
  \end{itemize}
  \item \textbf{Does the presenter distinguish criticism of the trust or its financial practices from criticism of the temple or religion itself?}
  \begin{itemize}[leftmargin=1.4em, itemsep=2pt, topsep=2pt]
    \item \textbf{Yes---devotion to the temple is the show's foundational premise.} Anjana speaks from \textit{within} a devout Hindu consensus (\texttt{"prabhu shri Ram ke Mandir"}, \texttt{[00:00]}; \texttt{"500 salon ke sangharsh"}, \texttt{[01:16]}), criticizing \textit{management} (\texttt{prabandhan}, lack of frisking at \texttt{[33:19--33:47]}) precisely because \texttt{"koi bhagvan se oopar kaise ho gaya"} [how did anyone become above God?] (\texttt{[17:08--17:17]}).
  \end{itemize}
  \item \textbf{Who is portrayed as protecting donors' interests, and who as exploiting religious sentiment?}
  \begin{itemize}[leftmargin=1.4em, itemsep=2pt, topsep=2pt]
    \item Every camp claims the mantle of protector in Aaj Tak's multi-voice arena: CM Yogi (\texttt{[02:53]}) and VHP (\texttt{[17:50]}) claim the Trust/State initiated the SIT; SP (\texttt{[14:57--15:25]}) and Congress (\texttt{[24:55--25:22]}) claim opposition whistleblowing exposed the \texttt{"kuchakra"}; meanwhile Anjana positions \textbf{Nripendra Misra} (who handed over keys cleanly after construction, \texttt{[10:57--11:15]}) and \textbf{Aaj Tak's own questioning} (\texttt{[01:06]}, \texttt{[33:19]}) as neutral guardians of devotees' trust, while testing opposition leaders for political opportunism (\texttt{[16:10--16:20]}, \texttt{[26:52--27:10]}).
  \end{itemize}
\end{enumerate}

\subsection*{Apply the five shared features (Event 1)}
\begin{longtable}{@{}>{\raggedright\arraybackslash}p{0.19\textwidth} >{\raggedright\arraybackslash}p{0.38\textwidth} >{\raggedright\arraybackslash}p{0.37\textwidth}@{}}
\toprule
\textbf{Feature} & \textbf{Exact wording + timestamp} & \textbf{What it does in context} \\
\midrule
\endhead
\textbf{1. Naming} & Title: \texttt{"Ram Mandir ko kisne 'loota'?"} \newline  VTR \texttt{[08:51]}: \texttt{"sirf ek bhrashtachar ka mamla bhar nahin balki ek aisa paap hai"} vs VHP \texttt{[17:17]}: \texttt{"chanda bhi nahin de sakte samarpan hai"} \& \texttt{[18:50]}, \texttt{[21:30]}: \texttt{"ramadrohi"} & Anchor scare-quotes \texttt{'loota'} to maintain journalistic distance, VTR sacralizes financial crime into religious \texttt{"paap"} [sin], and VHP/pro-BJP speakers rename donations \texttt{"samarpan"} [sacred surrender] while branding opposition critics \texttt{"ramadrohi"} [traitors to Ram]. \\
\addlinespace[4pt]
\textbf{2. Presupposition} & Anchor to SP \texttt{[16:20]}: \texttt{"par aisa kya hai ki akhilesh ji Ram Mandir nahin jate?"} \& \texttt{[26:52]}: \texttt{"pran pratishtha ke taim jise aapne BJP ka ivent kah ke apne aap ko doori bana li thi. aaj chande ke mudde pe kyon bhaee aapko... Ram bhakton ki aastha ki duhaee dene lag gae?"} & Presupposes that boycotting the 2024 Pran Pratishtha ceremony makes SP's concern for devotees' donations suspect and politically opportunistic, shifting the burden of proof onto the whistleblower party! \\
\addlinespace[4pt]
\textbf{3. Agency} & Nripendra Misra \texttt{[04:56--05:19]}: \texttt{"SBI ne isko most non siriyas mainar mein impliment kiya... suparavijan totli laiking tha"} \& Sangeet Ragi \texttt{[10:18]}: \texttt{"jo paise ki ginti mein shamil log the unhonne gadabdi ki hai"} & Deflects institutional agency away from top BJP/RSS/Trust leadership onto an external bank (\texttt{SBI}) and low-level counting clerks (\texttt{ginti mein shamil log}), insulating Champat Rai (\texttt{[12:25--12:56]}) and Nripendra Misra (\texttt{[10:57--11:24]}). \\
\addlinespace[4pt]
\textbf{4. Antagonisation} & CM Yogi \texttt{[00:29]}, \texttt{[02:30]}: \texttt{"Ram bhakton par, kaar sevkon par goli chalane... vale log aaj updesh dene chale hain"} \newline  Audience member \texttt{[28:08]}: \texttt{"aapko hak hi nahin hai... wo hamare bhagvan Ram hai ham apne aap dekh lenge"} & Invokes 1990 historical trauma (\texttt{kaar sevkon par goli chalane vale}) to draw an exclusionary boundary between authentic \texttt{"sanataniyon"} / \texttt{"Ram bhakton"} (\texttt{"hamare ghar ka mamla"}, \texttt{[26:33]}) and illegitimate \texttt{"ramadrohi"} / \texttt{"pappoo parti"} outsiders (\texttt{[25:49]}). \\
\addlinespace[4pt]
\textbf{5. Intensification / softening} & \textbf{Intensification (VTR):} \texttt{"bahut toophan khada kar chuka hai"} \texttt{[00:10]}, \texttt{"bhrashtachariyon ne dhava bola"} \texttt{[02:02]}, \texttt{"paap... prayashchit shayad nahin"} \texttt{[08:51]} \newline  \textbf{Softening (VHP):} \texttt{"500 varsh intjar kiya hai. kya ham 15 din aur intjar nahin kar sakte?"} \texttt{[18:00]}; \texttt{"kuch avyavastha huee hai"} \texttt{[34:21]} & Melodramatic opening metaphors (\texttt{toophan}, \texttt{dhava}, \texttt{paap}) hook prime-time viewers, while VHP softens systemic theft into minor \texttt{"avyavastha"} [disarray] and dwarfs a 15-day SIT wait against a \texttt{"500 varsh"} civilizational struggle. \\
\bottomrule
\end{longtable}

\subsection*{Short summary (Event 1)}
\begin{itemize}[leftmargin=1.5em, itemsep=3pt, topsep=4pt]
  \item \textbf{What interpretation of the event does this video encourage?} That a deeply distressing lapse in cash-counting supervision (\texttt{"nigrani ki kami"}, \texttt{[05:56]}) allowed insiders/clerks to steal sacred offerings (\texttt{chadhava}) at Ayodhya---demanding strict accountability (\texttt{[33:19--33:47]}) while giving viewers reason to await CM Yogi's 15-day SIT report (\texttt{[00:38]}, \texttt{[38:09]}) and to view SP/Congress outrage with skepticism given their past distance from the Ram Mandir (\texttt{[16:10--16:20]}, \texttt{[26:52--27:10]}).
  \item \textbf{Whose perspective is foregrounded, challenged, or missing?} \textit{Foregrounded:} CM Yogi Adityanath's SIT assurance, Nripendra Misra's technocratic diagnosis, VHP/RSS defenders, SP/Congress spokespersons, and vocal pro-BJP studio youth (\texttt{[26:16--29:53]}). \textit{Challenged:} Both the Trust's lack of frisking/FIR (\texttt{[16:52--17:17]}, \texttt{[22:09]}, \texttt{[33:19]}) AND the opposition's religious credentials (\texttt{[07:10--07:46]}, \texttt{[26:52]}, \texttt{[36:42]}). \textit{Missing:} Official BJP party spokesperson (\texttt{[13:21--13:30]}), SBI's institutional response to Nripendra Misra's blame (\texttt{[04:56]}), and unreceipted donors like the Sindhi silver-brick contributors.
  \item \textbf{What evidence supports your interpretation? What complicates it?} \textit{Supports:} Exclusive counting-room CCTV visuals (\texttt{[03:30]}), \texttt{Rs.\,10 lakh} cash recovery from Lavkush Mishra (\texttt{[03:49]}), and Nripendra Misra's candid on-camera admission of informal/missing frisking and double-lock failures (\texttt{[04:35--06:05]}). \textit{Complicates:} Whenever panelists raise structural precedents---like the 2021 Ayodhya land scam (\texttt{[34:58--35:08]}), Vaishno Devi Shrine Board reforms (\texttt{[23:56--24:14]}), or missing FIRs (\texttt{[23:18--23:29]})---the studio format drowns them out in personal insults (\texttt{"pharji prophesar"} \texttt{[14:15]}, \texttt{"pappoo parti"} \texttt{[25:49]}) and 1990 kar-sevak diversions (\texttt{"stop dis shauting maich"}, \texttt{[35:38]}).
\end{itemize}
\vspace{4pt}\noindent\rule{\textwidth}{0.4pt}\vspace{4pt}

\section*{Event2. Students' protest regarding NEET paper leak}
\textbf{Focus:} How does coverage construct the legitimacy of student protest and responsibility for the underlying problem?

\begin{itemize}[leftmargin=1.5em, itemsep=3pt, topsep=4pt]
  \item \textbf{Outlet:} Aaj Tak (\textit{AAJTAK 2 LIVE | AT2 LIVE} anchored by Himanshu Dixit, with ground reporter Sushant Mehra)
  \item \textbf{Date:} \texttt{2026-07-21} \textit{(the day AFTER the 20 July clashes)}
  \item \textbf{Video title:} \textit{AAJTAK 2 LIVE | NEET PROTEST| PM aavas ke bahar dharne par baithe Rahul Gandhi! | AT2 LIVE} (\url{https://www.youtube.com/live/JXLIs0WtJAM?si=3_7gAW9j4KDBWCUo})
  \item \textbf{Timestamp(s):} Relevant Event 2 segment \texttt{[00:04--23:50]} (\texttt{100\%} of live bulletin; total duration \texttt{24:10}).
  \item \textbf{Analyst:} Baani Solanki (\texttt{2025114016})
  \item \textbf{Language:} Hindi / Hinglish (Native Hindi broadcast audio, Romanized script; non-AI-dubbed)
\end{itemize}

\subsection*{Event-specific questions (Event 2)}
\begin{enumerate}[leftmargin=1.6em, itemsep=4pt, topsep=4pt]
  \item \textbf{Are participants called "students", "candidates", "protesters", "agitators", "mobs", or "political pawns"?}
  \begin{itemize}[leftmargin=1.4em, itemsep=2pt, topsep=2pt]
    \item \textbf{Shifts noticeably across the 24-minute broadcast from neutral to law-and-order delegitimization:}
    \item \textbf{Opening \& Opposition Soundbites (\texttt{[00:14--09:56]}):} Anchor uses neutral \texttt{"pradarshankariyon"} [protesters] (\texttt{[00:14]}) and \texttt{"chhatron"} [students] (\texttt{[14:50]}), while Rahul Gandhi, Akhilesh Yadav, and Dimple Yadav call them \texttt{"students"}, \texttt{"yuvaon"}, and \texttt{"bachche jo door-door se chalakar ke aae hain"} (\texttt{[04:37]}, \texttt{[06:40]}, \texttt{[13:34]}).
    \item \textbf{Mid-Bulletin Official Quote (\texttt{[10:04--11:10]}):} Airs Maharashtra CM Devendra Fadnavis claiming \texttt{"kuch log ismen sirf apna ejenda poora karne ke lie ghuse hue hain"} who don't even know what NEET is.
    \item \textbf{Ground \& Police Segment (\texttt{[15:43--23:50]}):} Reporter Sushant Mehra introduces \texttt{"asamajik tatva"} [anti-social elements] (\texttt{[15:52]}), and the anchor explicitly adopts \texttt{"upadraviyon"} [rioters/miscreants] (\texttt{"kaise upadraviyon ne police valon ko nishana banaya"}, \texttt{[21:20]}; \texttt{"upadravi unke chehre par hamla kar rahe the"}, \texttt{[22:38]})!
  \end{itemize}
  \item \textbf{Are their stated demands explained, summarised, mocked, questioned, or left out?}
  \begin{itemize}[leftmargin=1.4em, itemsep=2pt, topsep=2pt]
    \item \textbf{Briefly summarised two-thirds of the way into the bulletin (\texttt{[16:01--16:31]}):} Ground reporter Sushant Mehra lists CJP's three core demands presented to J.P. Nadda---(i) Education Minister's resignation, (ii) \texttt{Rs.\,1 crore} compensation to families of students who died by suicide after the NEET leak, and (iii) release of Sonam Wangchuk---while also covering Delhi High Court's order shifting Wangchuk to Medanta Hospital (\texttt{[02:02--02:54]}).
  \end{itemize}
  \item \textbf{Does the presenter focus on the original grievance, disruption, political affiliations, or public inconvenience?}
  \begin{itemize}[leftmargin=1.4em, itemsep=2pt, topsep=2pt]
    \item \textbf{Overwhelmingly focuses on (a) high-profile party politics / UP electoral horse-race (\texttt{[00:31--02:02]}, \texttt{[03:04--14:42]}) and (b) Parliamentary/street disruption \& police injuries (\texttt{[15:43--23:50]})}, rather than the NEET exam mechanics or student experiences. Even the video title foregrounds Rahul Gandhi (\texttt{PM aavas ke bahar dharne par baithe Rahul Gandhi!}), and an entire special package (\texttt{[11:18--14:42]}) evaluates the student protest purely as a \texttt{2027 UP election} contest (\texttt{"lakhanvi siyasat"}, \texttt{"dangal"}) comparing Dimple Yadav/SP against Rahul Gandhi/Congress, BSP, AAP, and AIMIM!
  \end{itemize}
  \item \textbf{Who is represented as responsible: examination authorities, government, students, political parties, or others?}
  \begin{itemize}[leftmargin=1.4em, itemsep=2pt, topsep=2pt]
    \item \textbf{Distributed across partisan claims, with editorial sympathy tilting toward police on street violence:} Opposition clips blame PM Modi, Home Minister, and Education Minister (\texttt{[04:47--05:05]}, \texttt{[06:19--07:08]}, \texttt{[08:00--09:48]}, \texttt{[13:34]}); PM Modi's NDA speech blames state-level \texttt{"dhandhe mein involva"} paper-leak criminals and accuses the opposition of \texttt{"rajniti"} (\texttt{[04:29--04:37]}, \texttt{[05:05--05:51]}); Parliamentary Minister blames opposition MPs for stalling the Supreme Court Judges Bill (\texttt{[07:08--07:50]}); and the final 8 minutes (\texttt{[15:43--23:50]}) blame \texttt{"asamajik tatva"} and \texttt{"upadraviyon"} breaking barricades at an \texttt{"illigal"} march for the 20 July violence.
  \end{itemize}
  \item \textbf{Are questions framed around whether students have legitimate grievances, or why they are causing disruption?}
  \begin{itemize}[leftmargin=1.4em, itemsep=2pt, topsep=2pt]
    \item \textbf{Framed primarily around political gamesmanship and law-and-order management:} Who is winning the opposition \texttt{"dangal"} over UP issues in Delhi (\texttt{[11:18--14:42]}), why Parliament was adjourned thrice (\texttt{11:00}, \texttt{12:00}, \texttt{2:00} at \texttt{[03:04--04:10]}), and how \texttt{"upadraviyon"} injured \texttt{110--118} police personnel during an \texttt{"illigal"} march (\texttt{[19:25--23:50]}).
  \end{itemize}
  \item \textbf{Are claims of violence, disorder, or political manipulation attributed to sources and substantiated?}
  \begin{itemize}[leftmargin=1.4em, itemsep=2pt, topsep=2pt]
    \item \textbf{Asymmetric substantiation:}
    \item Claims of \textbf{protester violence against police} are substantiated with official FIR charges (\texttt{5 FIRs} at \texttt{[19:25--19:43]}), casualty counts (\texttt{110--118} injured police including \texttt{5 IPS} vs \texttt{\textasciitilde{}70} protesters at \texttt{[18:09--18:20]}, \texttt{[19:43]}, \texttt{[20:51]}), and \textbf{three consecutive on-camera police interviews} (\texttt{DCP Sachin Sharma}, \texttt{Addl. DCP Anand Mishra} with a fractured hand, \texttt{Addl. DCP Jagdev} with a swollen/fractured hand at \texttt{[19:51--23:50]}).
    \item Claims of \textbf{police/plainclothes brutality against students} (stitched heads, broken limbs, Hitler-style \texttt{"praivet dres"} lathi gangs, torn girls' clothes) appear \textbf{only inside opposition politicians' soundbites} (Rahul Gandhi \texttt{[04:47]}, Akhilesh Yadav \texttt{[06:19--07:08]}, Kharge \texttt{[08:38--08:59]}, \texttt{[09:56]})---with \textbf{zero injured students interviewed directly by Aaj Tak}!
  \end{itemize}
  \item \textbf{How are police and government actions described compared with protesters' actions? Are similar actions evaluated differently?}
  \begin{itemize}[leftmargin=1.4em, itemsep=2pt, topsep=2pt]
    \item Anchor Himanshu Dixit explicitly endorses police action as civic duty: \texttt{"kanoon vyavastha ko durust karna ye police ka kaam hai"} (\texttt{[19:09]}) and praises injured officers returning to \texttt{"24 * 7 janta ki seva"} (\texttt{[23:30--23:50]}). Note the causal syntax in Sushant Mehra's ground report (\texttt{[15:43]}): \texttt{"us pradarshan ne ugra roop liya, pattharbaji huee aur uske baad tiyar gais chhode gae, police par hamla huaa"}---placing crowd escalation and stone-pelting \textit{first} and tear-gas \textit{after} in passive voice (\texttt{chhode gae}), the exact inverse of Dhruv Rathee (\texttt{l-6Pm-112io}, \texttt{[07:43--08:43]}) and Mohak Mangal (\texttt{v-RikSqb0Vw}, \texttt{[01:21--02:06]}).
  \end{itemize}
  \item \textbf{Are students allowed to explain their demands in their own words, or does the presenter mainly speak for them?}
  \begin{itemize}[leftmargin=1.4em, itemsep=2pt, topsep=2pt]
    \item \textbf{Zero students speak in their own words in the entire \texttt{24:10} broadcast.} Every single soundbite belongs to a politician (\texttt{Rahul Gandhi}, \texttt{PM Modi}, \texttt{Akhilesh Yadav}, \texttt{BJP Minister}, \texttt{Priyanka Gandhi}, \texttt{Mallikarjun Kharge}, \texttt{Devendra Fadnavis}, \texttt{Dimple Yadav}) or a senior IPS officer (\texttt{DCP Sachin Sharma}, \texttt{Addl. DCP Anand Mishra}, \texttt{Addl. DCP Jagdev}), mediated by Aaj Tak anchors and reporters.
  \end{itemize}
  \item \textbf{Does the outlet's framing differ between coverage of the NEET paper-leak issue before the protest and coverage once students begin protesting?}
  \begin{itemize}[leftmargin=1.4em, itemsep=2pt, topsep=2pt]
    \item \textbf{Yes, sharply.} While the underlying NEET paper leak is treated via PM Modi's quote as a non-partisan crime against youth (\texttt{"paap kaam"}, \texttt{[05:25]}), once students march on 20 July and politicians join on 21 July, Aaj Tak's coverage pivots entirely into \textbf{partisan spectacle} (\texttt{handcuff dharna}, \texttt{UP 2027 dangal}) and \textbf{law-and-order criminalization} (\texttt{5 FIRs}, \texttt{"upadraviyon"}, injured DCPs).
  \end{itemize}
\end{enumerate}

\subsection*{Apply the five shared features (Event 2)}
\begin{longtable}{@{}>{\raggedright\arraybackslash}p{0.19\textwidth} >{\raggedright\arraybackslash}p{0.38\textwidth} >{\raggedright\arraybackslash}p{0.37\textwidth}@{}}
\toprule
\textbf{Feature} & \textbf{Exact wording + timestamp} & \textbf{What it does in context} \\
\midrule
\endhead
\textbf{1. Naming} & Opening \texttt{[00:14]}: \texttt{"pradarshankariyon"} $\rightarrow$ Reporter \texttt{[15:52]}: \texttt{"asamajik tatva"} $\rightarrow$ Anchor \texttt{[21:20]}, \texttt{[22:38]}: \texttt{"kaise upadraviyon ne police valon ko nishana banaya"} & Tracks an editorial lexical slide across the bulletin from neutral \texttt{"pradarshankariyon"} [protesters] to \texttt{"asamajik tatva"} [anti-social elements] to unhedged \texttt{"upadraviyon"} [rioters] in the anchor's own voice, legitimizing the 5 police FIRs. \\
\addlinespace[4pt]
\textbf{2. Presupposition} & Anchor to Addl. DCP \texttt{[21:20]}: \texttt{"kaise upadraviyon ne police valon ko nishana banaya"} \& \texttt{[21:30]}: \texttt{"kyunki jo police hai wo toh rok rahi thi jo protestar the unko"} & The \textit{how}-clause (\texttt{kaise upadraviyon ne... nishana banaya}) and causal premise (\texttt{kyunki jo police hai wo toh rok rahi thi}) presuppose as established fact that police were merely restraining peacefully while rioters targeted them. \\
\addlinespace[4pt]
\textbf{3. Agency} & Reporter \texttt{[15:43]}: \texttt{"us pradarshan ne ugra roop liya, pattharbaji huee aur uske baad tiyar gais chhode gae, police par hamla huaa"} vs Anchor \texttt{[00:14]}: \texttt{"pradarshankariyon par jis tarah ki karvaee ki gaee"} & Police violence (\texttt{karvaee ki gaee}, \texttt{tiyar gais chhode gae}) is systematically passivized with the police agent deleted and sequenced \textit{after} \texttt{"pattharbaji huee"}, whereas \texttt{"asamajik tatva"} and \texttt{"upadraviyon"} receive active ergative verbs (\texttt{gadiyan todi}, \texttt{nishana banaya}). \\
\addlinespace[4pt]
\textbf{4. Antagonisation} & VTR \texttt{[12:01--12:10]}: \texttt{"satta paksh ke UP vale bhagidar... Mangal milan kar rahe hain toh satta ke lie lad raha vipaksh apne-apne hisab se dangal kar raha hai"} & Reframes a grassroots youth movement into a cynical, symmetric gladiatorial contest (\texttt{Mangal milan} vs \texttt{dangal}, \texttt{"siyasi rotiyan sek rahi hai"} \texttt{[06:08]}) among rival political parties eyeing UP elections. \\
\addlinespace[4pt]
\textbf{5. Intensification / softening} & \textbf{Softening (police force \& gov delay):} \texttt{"sootron ke havale se khabar hai ki sarkar... batchit ke darvaje khule rakhegi"} \texttt{[09:56]}, \texttt{[17:07]}; \texttt{"priventiv arenjment"} \texttt{[20:02]} \newline  \textbf{Intensification (police pain):} \texttt{"poora haath aapka laal pad gaya"} \texttt{[22:46]}, \texttt{"118 se zyada pulisakarmi... panch aaeepies adhikari... gambhir chot"} \texttt{[18:09]} & Softens 30 days of government silence via anonymous \texttt{"sources"} promising \texttt{"open doors for dialogue"}, while visually and emotionally intensifying police injuries (\texttt{poora haath aapka laal pad gaya}) over student injuries (\texttt{70 ke aaspas} mentioned in a single passing clause at \texttt{[18:20]}). \\
\bottomrule
\end{longtable}

\subsection*{Short summary (Event 2)}
\begin{itemize}[leftmargin=1.5em, itemsep=3pt, topsep=4pt]
  \item \textbf{What interpretation of the event does this video encourage?} That the NEET protest aftermath on 21 July is primarily a high-stakes \textbf{ruling-vs-opposition political showdown} (from Parliament adjournments and Rahul Gandhi's PM-residence dharna to SP vs Congress \texttt{UP 2027} positioning) alongside a serious \textbf{law-and-order breach} in which \texttt{"upadraviyon"} at an \texttt{"illigal"} march injured over \texttt{110--118} dedicated police personnel while the government keeps \texttt{"batchit ke darvaje khule"} (\texttt{[09:56]}, \texttt{[15:43--23:50]}).
  \item \textbf{Whose perspective is foregrounded, challenged, or missing?} \textit{Foregrounded:} Delhi Police officers (\texttt{DCP Sachin Sharma}, \texttt{Addl. DCP Anand Mishra}, \texttt{Addl. DCP Jagdev}), PM Modi, CM Fadnavis, and national/UP party leaders (\texttt{Rahul Gandhi}, \texttt{Akhilesh/Dimple Yadav}, \texttt{Kharge}). \textit{Challenged:} Congress's street commitment compared to SP (\texttt{[11:18--14:14]}) and protesters' breach of barricades (\texttt{[17:50--23:50]}). \textit{Missing:} \textbf{Protesting NEET students and CJP members themselves}---not a single student or injured youth is interviewed directly!
  \item \textbf{What evidence supports your interpretation? What complicates it?} \textit{Supports:} Verifiable Parliamentary adjournment timestamps (\texttt{11:00}, \texttt{12:00}, \texttt{2:00} at \texttt{[03:35--04:00]}), Delhi HC's Medanta transfer order for Sonam Wangchuk (\texttt{[02:02--02:54]}), \texttt{5 FIRs}, and firsthand interviews with visibly bandaged IPS officers (\texttt{[19:25--23:50]}). \textit{Complicates:} Opposition leaders within Aaj Tak's own broadcast cite stitched student heads, broken limbs, and Hitler-like \texttt{"praivet dres"} lathi squads (\texttt{[06:30--07:08]}), and Addl. DCP Anand Mishra notably evades the reporter's brief question (\texttt{"police ki taraph se huaa ya unki taraph se huaa?"}) with \texttt{"abhi toh investigeshan chal raha hai"} (\texttt{[21:39--21:49]}).
\end{itemize}
\vspace{4pt}\noindent\rule{\textwidth}{0.4pt}\vspace{4pt}

\section*{Common form: Blank excerpt record}
\textit{Copy or print this page whenever you find a passage worth analysing. Keep the exact wording and enough context to interpret it.}


\subsection*{Excerpt Record 1 (Event 1 --- Nripendra Misra Interview \& Anchor's Institutional Separation)}
\begin{itemize}[leftmargin=1.5em, itemsep=3pt, topsep=4pt]
  \item \textbf{Outlet:} Aaj Tak (\textit{Halla Bol}) | \textbf{Date:} \texttt{2026-06-19}
  \item \textbf{Video title:} \textit{Halla Bol With Anjana Om Kashyap Full Show: Ram Mandir ko kisne 'loota'? | Ram Mandir Donation | BJP}
  \item \textbf{Timestamp(s):} \texttt{[04:35--05:19]} \& \texttt{[10:57--11:24]} | \textbf{Analyst:} Baani Solanki | \textbf{Language:} Hindi / Hinglish (Roman script)
  \item \textbf{Exact transcript excerpt:}
    \begin{quote}\small\ttfamily [04:35] [Nripendra Misra:] aaj poora prabandhan anaupcharik tarike se hai. kis prakar se entri aur egjit par phrisking hogi us vyakti ki... [04:56] kadachit in kagjon ka koi anupalan nahin huaa hai. SBI ne isko most non siriyas mainar mein [05:09] impliment kiya... woh suparavijan totli laiking tha...\end{quote}
    \begin{quote}\small\ttfamily [10:57] [Anjana Om Kashyap:] ek toh kisi bhi prabandhan samiti mein nirpen mishra hai hi nahin... shri Ram janmabhoomi [11:07] nirman samiti ke adhyaksh hain. nirman ka kaam poora ho chuka hai. ek-ek jagah ki phail chabi sab kuch unhonne Trust ko diya hai. [11:15] prabandhan samiti nahin Trust hota hai... champat raay ji ki baat karen ya [11:24] anya logon ki baat karen prabandhan ki jimmedari un logon ki hoti hai.\end{quote}
  \item \textbf{Feature(s):} $\boxtimes$ Naming $\square$ Presupposition $\boxtimes$ Agency $\square$ Antagonisation $\boxtimes$ Intensification / softening
  \item \textbf{Who is speaking? Is this the outlet's wording or a quotation?} Quotation from Construction Committee Chair Nripendra Misra (\texttt{[04:35--05:19]}) followed by anchor Anjana Om Kashyap's forceful editorial correction (\texttt{[10:57--11:24]}).
  \item \textbf{What does the wording imply, foreground, soften, or leave unstated?} Nripendra Misra softens internal culpability into \texttt{"anaupcharik"} [informal] management while assigning active blame to \texttt{"SBI"} (\texttt{"most non siriyas mainar"}). Anjana then draws a firewall (\texttt{"nirman alag hai, prabandhan alag hai"}) that completely insulates PM Modi's former Principal Secretary Nripendra Misra (and thus Delhi/PMO) while localizing responsibility onto Trust managers like Champat Rai.
  \item \textbf{What context or evidence supports or complicates this interpretation?} Supported by the formal administrative division between the Construction Committee and the Trust Secretariat. Complicated by \textbf{The Wire} (\texttt{7rUL8fGgRJA}, \texttt{[11:48--14:55]}) and \textbf{Dhruv Rathee} (\texttt{oeH0fOSJuX8}, \texttt{[04:28]}), who argue that as PMO's deputed overseer, Misra scrutinized land/material acquisitions and necessarily kept top leadership informed.
  \item \textbf{How does another outlet describe the same event or actor?} \textbf{The Wire} (\texttt{7rUL8fGgRJA}, \texttt{[12:27--13:42]}) treats Nripendra Misra as the prime institutional link proving PMO awareness since 2021, whereas \textbf{Aaj Tak} presents him as an unimpeachable whistleblower-statesman who cleanly handed over \texttt{"ek-ek jagah ki phail chabi"} (\texttt{[11:07]}).
  \item \textbf{My interpretation, with reasons:} This excerpt shows subtle mainstream TV gatekeeping: Aaj Tak allows hard questions about Ayodhya's \textit{local counting management} while ring-fencing PMO-linked technocrats from political fallout.
\end{itemize}

\subsection*{Excerpt Record 2 (Event 1 --- Testing the Whistleblower's Devotion \& "Ramadrohi" Attacks)}
\begin{itemize}[leftmargin=1.5em, itemsep=3pt, topsep=4pt]
  \item \textbf{Outlet:} Aaj Tak (\textit{Halla Bol}) | \textbf{Date:} \texttt{2026-06-19}
  \item \textbf{Video title:} \textit{Halla Bol With Anjana Om Kashyap Full Show: Ram Mandir ko kisne 'loota'? | Ram Mandir Donation | BJP}
  \item \textbf{Timestamp(s):} \texttt{[18:00--19:00]} \& \texttt{[26:52--27:10]} | \textbf{Analyst:} Baani Solanki | \textbf{Language:} Hindi / Hinglish (Roman script)
  \item \textbf{Exact transcript excerpt:}
    \begin{quote}\small\ttfamily [18:00] [VHP Spokesperson Pravesh Chaudhary:] aap ye dekhie Anjana ji is Hindu samaj ne 500 varsh intjar kiya hai. kya ham 15 din aur intjar nahin kar [18:10] sakte? ... ye jo sootron ke havale se khabar aa rahi hai... sootron pe vishvas mat karie... [18:50] ye wo log hain jo prashn khada karte hain jinko Ram se kuch nahin hai. Ram drohi Ram ramadrohi...\end{quote}
    \begin{quote}\small\ttfamily [26:52] [Anjana Om Kashyap to SP's Anurag Bhadauria:] pran pratishtha ke taim jise aapne BJP ka ivent kah ke apne aap ko doori bana li thi. aaj chande ke [27:01] mudde pe kyon bhaee aapko Ram Mandir matalab Ram bhakton ki aastha ki duhaee dene lag gae.\end{quote}
  \item \textbf{Feature(s):} $\boxtimes$ Naming $\boxtimes$ Presupposition $\square$ Agency $\boxtimes$ Antagonisation $\boxtimes$ Intensification / softening
  \item \textbf{Who is speaking? Is this the outlet's wording or a quotation?} VHP spokesperson Pravesh Chaudhary (\texttt{[18:00--19:00]}) and anchor Anjana Om Kashyap's own question to SP spokesperson Anurag Bhadauria (\texttt{[26:52--27:10]}).
  \item \textbf{What does the wording imply, foreground, soften, or leave unstated?} VHP juxtaposes \texttt{"500 varsh"} against \texttt{"15 din"} to make immediate scrutiny look impatient, and labels critics \texttt{"ramadrohi"}. Crucially, even though Congress's Deepak Jha objects (\texttt{"sawal poochhne vale ko deshadrohi aur ramadrohi kyon bata rahe hain?"}, \texttt{[24:55]}), Anchor Anjana herself adopts the \textit{ad hominem} devotional test at \texttt{[26:52]}---implying only those who attended Pran Pratishtha have moral standing to question donation theft.
  \item \textbf{What context or evidence supports or complicates this interpretation?} Supported factually by SP/Congress leaders declining the January 2024 Pran Pratishtha invitation. Complicated by basic public-trust law: financial fraud in a statutory public trust requires investigation regardless of who raises the complaint, and even Ayodhya sants/kar-sevaks expressed anger in Aaj Tak's own VTR (\texttt{[08:23--08:42]}).
  \item \textbf{How does another outlet describe the same event or actor?} \textbf{Republic TV} (\texttt{0H-bruoEkaY}, \texttt{[02:46--03:14]}, \texttt{[04:18]}) pushes this exact same line to an extreme (\texttt{"don't have the asta to visit the Ram temple"}), whereas \textbf{Dhruv Rathee} (\texttt{oeH0fOSJuX8}, \texttt{[14:24]}) counters: \texttt{"Who is playing with their faith? The scammer or the one who is questioning the scam?"}.
  \item \textbf{My interpretation, with reasons:} Aaj Tak occupies a middle ground between Republic TV and independent YouTubers---acknowledging the theft hurts devotees, yet repeatedly validating the ruling camp's "who-are-you-to-ask?" credential test against the opposition.
\end{itemize}

\subsection*{Excerpt Record 3 (Event 1 --- Anchor Pressing VHP on Unchecked VIP Trustees \& Missing Frisking)}
\begin{itemize}[leftmargin=1.5em, itemsep=3pt, topsep=4pt]
  \item \textbf{Outlet:} Aaj Tak (\textit{Halla Bol}) | \textbf{Date:} \texttt{2026-06-19}
  \item \textbf{Video title:} \textit{Halla Bol With Anjana Om Kashyap Full Show: Ram Mandir ko kisne 'loota'? | Ram Mandir Donation | BJP}
  \item \textbf{Timestamp(s):} \texttt{[33:19--34:21]} | \textbf{Analyst:} Baani Solanki | \textbf{Language:} Hindi / Hinglish (Roman script)
  \item \textbf{Exact transcript excerpt:}
    \begin{quote}\small\ttfamily [33:19] [Anjana Om Kashyap:] akauntebal sabko hona padega. aur yeh sawal hai ki sab wahan par agar chadhava [33:28] chori ho raha tha kisi ko dikh hi nahin raha tha. woh log jo thaat baat se wahan ghoomte hain Mandir mein kyunki sare niyam toh shraddhaluon ke lie hain. yahan se jana hai. [33:38] iske aage nahin badhna hai. jinko ekses tha. 40 daan petiyon ka paisa ja raha tha. kalekshan ho raha tha. phrisking ho rahi hai nahin ho rahi hai. kisi ko dikh hi nahin raha [33:47] tha... [34:13] [VHP Spokesperson:] aapka nyas par prashn khada karna hi invailid hai... [Anjana:] jimmedari kiski thi? [34:21] thik hai agar Trust ki jimmedari nahin thi toh kiski jimmedari? [VHP:] bilakul Trust ki jimmedari thi kuch avyavastha huee hai.\end{quote}
  \item \textbf{Feature(s):} $\boxtimes$ Naming $\boxtimes$ Presupposition $\boxtimes$ Agency $\boxtimes$ Antagonisation $\boxtimes$ Intensification / softening
  \item \textbf{Who is speaking? Is this the outlet's wording or a quotation?} Anchor Anjana Om Kashyap challenging VHP spokesperson Pravesh Chaudhary (\texttt{[33:19--34:21]}).
  \item \textbf{What does the wording imply, foreground, soften, or leave unstated?} Anjana contrasts ordinary \texttt{"shraddhaluon"} [pilgrims], who face strict queues and barriers, against privileged insiders \texttt{"woh log jo thaat baat se wahan ghoomte hain"} [those who strut around in luxury/swagger]. When VHP tries to declare questioning the Trust \texttt{"invailid"} (\texttt{[34:13]}), her sharp question \texttt{"agar Trust ki jimmedari nahin thi toh kiski jimmedari?"} forces VHP to concede responsibility---though softened immediately to \texttt{"kuch avyavastha"} [some disarray] (\texttt{[34:21]}).
  \item \textbf{What context or evidence supports or complicates this interpretation?} Important complicating evidence against caricaturing all mainstream Hindi TV as 100\% uncritical: unlike Republic TV, Anjana directly forces a VHP spokesperson to retract \texttt{"nyas par prashn... invailid hai"}.
  \item \textbf{How does another outlet describe the same event or actor?} Echoes \textbf{Sarthak Goswami}'s critique of VIP temple culture (\texttt{lb7gZnyZH0c}, \texttt{[02:29--02:38]}), though Anjana stops short of naming BJP/RSS patronage networks behind those insiders.
  \item \textbf{My interpretation, with reasons:} Because ordinary Hindu viewers resent VIP temple privileges (\texttt{"sare niyam toh shraddhaluon ke lie hain"}), Aaj Tak safely channels populist devotee frustration against local Trust arrogance without implicating CM Yogi or PM Modi.
\end{itemize}

\subsection*{Excerpt Record 4 (Event 2 --- Reframing Student Protest as UP Electoral "Dangal")}
\begin{itemize}[leftmargin=1.5em, itemsep=3pt, topsep=4pt]
  \item \textbf{Outlet:} Aaj Tak (\textit{AT2 LIVE}) | \textbf{Date:} \texttt{2026-07-21}
  \item \textbf{Video title:} \textit{AAJTAK 2 LIVE | NEET PROTEST| PM aavas ke bahar dharne par baithe Rahul Gandhi! | AT2 LIVE}
  \item \textbf{Timestamp(s):} \texttt{[11:18--12:27]} \& \texttt{[14:04--14:42]} | \textbf{Analyst:} Baani Solanki | \textbf{Language:} Hindi / Hinglish (Roman script)
  \item \textbf{Exact transcript excerpt:}
    \begin{quote}\small\ttfamily [11:18] uttar pradesh ke muddon par bhi dilli mein rajniti ho rahi hai paper leak ko lekar chadhava chori tak ke mamlon ko lekar dilli [11:25] mein samajvadi parti jitni aakramak rajniti kar rahi hai utni Congress nahin kar pa rahi. jantaramantar se sansad bhavan tak chal rahi lakhanvi siyasat ki kuch chaliya dikhate hain... [12:01] ek taraph satta paksh ke UP vale bhagidar... Mangal milan kar rahe hain toh satta ke lie lad [12:10] raha vipaksh apne-apne hisab se dangal kar raha hai... [12:27] samajvadi parti apne hisse ki siyasat bakhoobi kar rahi hai... [14:33] ab dekhna hoga ki sansad satra ke baki 18 dinon tak dilli mein UP [14:42] ki chunavi rajniti ke muddon ki mashal kaun vipakshi dal kitni majbooti se thame rakhta hai.\end{quote}
  \item \textbf{Feature(s):} $\boxtimes$ Naming $\boxtimes$ Presupposition $\boxtimes$ Agency $\boxtimes$ Antagonisation $\square$ Intensification / softening
  \item \textbf{Who is speaking? Is this the outlet's wording or a quotation?} Anchor Himanshu Dixit (\texttt{[11:18--11:25]}) and Aaj Tak Bureau VTR narrator (\texttt{[12:01--14:42]}).
  \item \textbf{What does the wording imply, foreground, soften, or leave unstated?} Translates both Event 1 (\texttt{chadhava chori}) and Event 2 (\texttt{NEET paper leak protest}) into strategic moves (\texttt{"kuch chaliya"}, \texttt{"dangal"}, \texttt{"chunavi rajniti ke muddon ki mashal"}) in the upcoming Uttar Pradesh election horse-race. Students and CJP become mere stage props (\texttt{"kokroch janta parti ke manch par dimpal yadav najar aaee"}, \texttt{[13:17]}) for grading opposition politicians' street performance!
  \item \textbf{What context or evidence supports or complicates this interpretation?} Supported by the visible presence of SP (\texttt{Dimple Yadav}), ASP (\texttt{Chandrashekhar Azad}), Congress (\texttt{Rahul/Priyanka}), AAP (\texttt{Kejriwal}), and NCP (\texttt{Sharad Pawar}, \texttt{Supriya Sule}) leaders on 20--21 July (\texttt{[12:37--16:48]}). Complicated by ground accounts in \textbf{Sarthak Goswami} (\texttt{EM1GeZufO6E}, \texttt{[02:29]}, \texttt{[08:08]}) and \textbf{Mohak Mangal} (\texttt{v-RikSqb0Vw}, \texttt{[05:47--09:18]}) showing that CJP ran for 30 days primarily as a student/citizen movement wary of partisan hijacking.
  \item \textbf{How does another outlet describe the same event or actor?} \textbf{Mohak Mangal} (\texttt{v-RikSqb0Vw}) mentions \textit{zero} opposition politicians across 11 minutes, keeping 100\% focus on citizens, whereas \textbf{Republic TV} (\texttt{HTvb6E\_1Qdc}, \texttt{[00:00--02:56]}) similarly obsesses over Rahul Gandhi (\texttt{"Rahul The Biggest Loser"}) to erase genuine student agency.
  \item \textbf{My interpretation, with reasons:} "Horse-race politicization" is a classic mainstream TV discourse strategy: by treating a citizen protest as \texttt{"lakhanvi siyasat"} and \texttt{"chunavi rajniti"}, Aaj Tak subtly drains the protest of non-partisan moral urgency without needing to scream like Arnab Goswami.
\end{itemize}

\subsection*{Excerpt Record 5 (Event 2 --- Ground Report Causal Sequencing \& Passive Tear-Gas)}
\begin{itemize}[leftmargin=1.5em, itemsep=3pt, topsep=4pt]
  \item \textbf{Outlet:} Aaj Tak (\textit{AT2 LIVE}) | \textbf{Date:} \texttt{2026-07-21}
  \item \textbf{Video title:} \textit{AAJTAK 2 LIVE | NEET PROTEST| PM aavas ke bahar dharne par baithe Rahul Gandhi! | AT2 LIVE}
  \item \textbf{Timestamp(s):} \texttt{[15:43--16:01]} \& \texttt{[17:50--19:18]} | \textbf{Analyst:} Baani Solanki | \textbf{Language:} Hindi / Hinglish (Roman script)
  \item \textbf{Exact transcript excerpt:}
    \begin{quote}\small\ttfamily [15:43] [Reporter Sushant Mehra:] kal jis tarike se pradarshan huaa us pradarshan ne ugra roop liya ah pattharbaji huee aur uske baad tiyar gais chhode gae police par hamla huaa pradarshankariyon ke bich mein jis tarike se [15:52] jo asamajik tatva unhonne ne kis tarike se gadiyan todi, dukanon mein todphod kari...\end{quote}
    \begin{quote}\small\ttfamily [17:50] tamam bairikeding ko todakar ke pradarshankari jo hai wo aage ki taraph badhte hue najar aae the... [18:09] 118 se zyada pulisakarmi jo hai wo ghayal hue the. unmen takriban agar kaha jae toh panch aaeepies adhikari hai jinke gambhir chot bhi aaee thi. in sabko lekar ke ismen 70 ke aaspas [18:20] pradarshankari bhi jo hai wo ghayal hue the... [19:09] [Anchor:] kanoon vyavastha ko durust karna ye police ka kaam hai.\end{quote}
  \item \textbf{Feature(s):} $\boxtimes$ Naming $\square$ Presupposition $\boxtimes$ Agency $\square$ Antagonisation $\boxtimes$ Intensification / softening
  \item \textbf{Who is speaking? Is this the outlet's wording or a quotation?} Aaj Tak ground correspondent Sushant Mehra live from Jantar Mantar (\texttt{[15:43--18:30]}) and studio anchor Himanshu Dixit (\texttt{[19:09]}).
  \item \textbf{What does the wording imply, foreground, soften, or leave unstated?} Notice the textbook \textbf{transitivity and ordering asymmetry}: 1. \texttt{"us pradarshan ne ugra roop liya"} $\rightarrow$ active protest escalation; 2. \texttt{"pattharbaji huee"} $\rightarrow$ stone-pelting placed \textit{before} police action; 3. \texttt{"aur uske baad tiyar gais chhode gae"} $\rightarrow$ passive voice (\texttt{chhode gae} = were released) deleting \texttt{"Delhi Police"} as agent and framing tear gas purely as a subsequent reaction (\texttt{uske baad}); 4. Casualty foregrounding: \texttt{118+} injured police (\texttt{5 IPS} with \texttt{"gambhir chot"}) receive detailed elaboration before \texttt{\textasciitilde{}70} injured protesters are tacked on in a single agentless clause (\texttt{"ghayal hue the"}, \texttt{[18:20]}) with zero mention of pellet guns or unbadged civil-clothes men.
  \item \textbf{What context or evidence supports or complicates this interpretation?} Police fracture/injury statistics are backed by named DCP interviews (\texttt{[19:43--23:50]}). Complicated by independent ground videos in \textbf{Dhruv Rathee} (\texttt{l-6Pm-112io}, \texttt{[07:43--08:43]}, \texttt{[11:05]}) showing police striking marching students at Mandi House \textit{before} they reached Jantar Mantar and firing tear gas at calmly seated protesters, plus \textit{The Hindu}'s pellet-gun surgery report cited on \textbf{The Wire} (\texttt{pYaTVvHeGoM}, \texttt{[10:07--10:38]}).
  \item \textbf{How does another outlet describe the same event or actor?} \textbf{Sarthak Goswami} (\texttt{EM1GeZufO6E}, \texttt{[03:39--04:53]}) reverses the exact causal chain---foregrounding active police lathi-charge, tear gas, and \texttt{"pelet gans ki bauchhar"} first, and describing police injuries as subsequent crowd \texttt{"ritalieshan"}.
  \item \textbf{My interpretation, with reasons:} Even in seemingly matter-of-fact live reporting, syntactic agency (active vs agentless passive) and chronological ordering (\texttt{"aur uske baad"}) quietly align viewer sympathy with state policing.
\end{itemize}

\subsection*{Excerpt Record 6 (Event 2 --- Anchor Adopting "Upadraviyon" \& Sympathetic DCP Interviews)}
\begin{itemize}[leftmargin=1.5em, itemsep=3pt, topsep=4pt]
  \item \textbf{Outlet:} Aaj Tak (\textit{AT2 LIVE}) | \textbf{Date:} \texttt{2026-07-21}
  \item \textbf{Video title:} \textit{AAJTAK 2 LIVE | NEET PROTEST| PM aavas ke bahar dharne par baithe Rahul Gandhi! | AT2 LIVE}
  \item \textbf{Timestamp(s):} \texttt{[21:12--22:55]} \& \texttt{[23:30--23:50]} | \textbf{Analyst:} Baani Solanki | \textbf{Language:} Hindi / Hinglish (Roman script)
  \item \textbf{Exact transcript excerpt:}
    \begin{quote}\small\ttfamily [21:12] [Anchor:] toh kal dilli mein hue CJP ke pradarshan mein edishanal disipi aanand mishra ko bhi chot lagi. unke haath mein phraikchar huaa. aaj tak se batchit mein aanand mishra ne bataya ki kaise [21:20] upadraviyon ne police valon ko nishana banaya... [21:49] [Addl. DCP Anand Mishra:] us samay kraud kuch aisa kraud tha ki jo ki ham log barabar unko samjha rahe the... kaphi unko parsuet kar rahe the ki bhaee [21:58] aap log jo hai aage na badhe kyunki ye march illigal hai...\end{quote}
    \begin{quote}\small\ttfamily [22:38] [Anchor:] disipi sekand jagdev ne bataya ki upadravi unke chehre par hamla kar rahe the... [22:46] [Reporter:] sar kahan-kahan chot lagi? toh poora haath aapka laal pad gaya... [23:39] [Addl. DCP Jagdev:] police ka kaam 24 * 7 janta ki seva mein khade rahna hai.\end{quote}
  \item \textbf{Feature(s):} $\boxtimes$ Naming $\boxtimes$ Presupposition $\boxtimes$ Agency $\boxtimes$ Antagonisation $\boxtimes$ Intensification / softening
  \item \textbf{Who is speaking? Is this the outlet's wording or a quotation?} Studio anchor Himanshu Dixit (\texttt{[21:12--21:20]}, \texttt{[22:38]}), Aaj Tak correspondent (\texttt{[22:46]}), and injured Delhi Police Addl. DCPs Anand Mishra and Jagdev.
  \item \textbf{What does the wording imply, foreground, soften, or leave unstated?} The anchor replaces \texttt{"chhatra"} / \texttt{"pradarshankari"} with \texttt{"upadraviyon"} [rioters/hooligans] (\texttt{[21:20]}, \texttt{[22:38]}) as an unhedged fact. Police action is softened to patient persuasion (\texttt{"barabar unko samjha rahe the"}, \texttt{"parsuet kar rahe the"}, \texttt{"24 * 7 janta ki seva"}), while the reporter's empathetic close-up (\texttt{"poora haath aapka laal pad gaya"}, \texttt{[22:46]}) humanizes injured officers---a care never extended to any injured NEET student on the broadcast.
  \item \textbf{What context or evidence supports or complicates this interpretation?} Supported by the officers' real X-rayed fractures and swelling (\texttt{[22:55--23:23]}). Complicated by the complete absence of parallel hospital interviews with injured students at RML Hospital (even though Sushant Mehra mentions at \texttt{[15:25]} and \texttt{[18:40]} that detained/injured students were at RML Hospital right nearby!).
  \item \textbf{How does another outlet describe the same event or actor?} \textbf{Mohak Mangal} (\texttt{v-RikSqb0Vw}, \texttt{[01:12--02:06]}) asks the exact same empathetic question (\texttt{"kaise huaa?"}) exclusively to \textit{injured young protesters}, whereas \textbf{Aaj Tak} asks \texttt{"yeh chot aapko kaise lagi?"} (\texttt{[21:20]}, \texttt{[22:46]}) exclusively to \textit{senior IPS officers}.
  \item \textbf{My interpretation, with reasons:} Comparing Mohak Mangal and Aaj Tak side by side reveals how journalistic \textbf{source selection} (whose bandages get the microphone) shapes reality even before a single overtly partisan adjective is spoken.
\end{itemize}
\vspace{4pt}\noindent\rule{\textwidth}{0.4pt}\vspace{4pt}

\section*{Common form: End-of-event summary}

\subsection*{End-of-Event Summary --- Event 1 (Ram Mandir Donation Controversy)}
\begin{itemize}[leftmargin=1.5em, itemsep=3pt, topsep=4pt]
  \item \textbf{The main way this outlet framed the event was:} As a painful devotional scandal (\texttt{"karodon arbon bhakton ki aastha se chhal"}, \texttt{"paap"}, \texttt{[02:19]}, \texttt{[08:51]}) caused by lax local supervision and SBI/counting-room informality (per Nripendra Misra, \texttt{[04:35--06:05]}), which warrants management accountability (\texttt{[33:19--33:47]}) while trusting CM Yogi's 15-day SIT probe (\texttt{[00:38]}, \texttt{[02:53]}) and scrutinizing opposition critics for partisan opportunism (\texttt{[16:10--16:20]}, \texttt{[26:52--27:10]}).
  \item \textbf{The strongest linguistic evidence was (quote + timestamp):} Title \texttt{"Ram Mandir ko kisne 'loota'?"} (with scare-quotes), Nripendra Misra's \texttt{"SBI ne isko most non siriyas mainar mein impliment kiya"} (\texttt{[04:56--05:09]}), and Anchor Anjana's \texttt{"pran pratishtha ke taim... doori bana li thi. aaj chande ke mudde pe kyon bhaee aapko... Ram bhakton ki aastha ki duhaee dene lag gae"} (\texttt{[26:52--27:01]}).
  \item \textbf{The shared feature(s) most clearly illustrated were:} \textbf{Presupposition} (challenging SP/Congress standing to ask questions because they skipped Pran Pratishtha at \texttt{[16:20]}, \texttt{[26:52]}), \textbf{Agency} (separating Construction Chair Nripendra Misra from Trust management at \texttt{[10:57--11:24]} and blaming SBI/40 counting staff), and \textbf{Antagonisation} (allowing \texttt{"ramadrohi"} vs \texttt{"kaar sevkon par goli chalane vale"} binaries to dominate the studio floor).
  \item \textbf{An alternative interpretation of the same evidence is:} That long-standing lack of frisking, absence of an immediate FIR (\texttt{[22:09]}, \texttt{[23:18]}), and prior 2021 land-deal controversies (\texttt{[34:58--35:08]}) point not to mere SBI informality, but to systemic political impunity inside a BJP/RSS-nominated Trust requiring a Supreme Court-monitored investigation rather than a state-government SIT (\texttt{[06:32--07:02]}).
  \item \textbf{What I cannot conclude from this transcript alone:} Because this episode aired on \texttt{2026-06-19} (Day 5 of the SIT probe), the transcript alone cannot show who would eventually be arrested on 26 June, whether senior trustees would resign, or how much total cash/jewellery was stolen.
\end{itemize}

\subsection*{End-of-Event Summary --- Event 2 (Students' Protest Regarding NEET Paper Leak)}
\begin{itemize}[leftmargin=1.5em, itemsep=3pt, topsep=4pt]
  \item \textbf{The main way this outlet framed the event was:} Through a dual \textbf{electoral-spectacle + police-law-and-order frame}: foregrounding Parliament adjournments, Rahul Gandhi's handcuff dharna, and SP-vs-Congress \texttt{UP 2027} turf battles (\texttt{"lakhanvi siyasat"}, \texttt{"dangal"}, \texttt{[11:18--14:42]}), while portraying the 20 July street clashes through 5 police FIRs and sympathetic interviews with injured IPS officers attacked by \texttt{"upadraviyon"} (\texttt{[15:43--23:50]}).
  \item \textbf{The strongest linguistic evidence was (quote + timestamp):} \texttt{"us pradarshan ne ugra roop liya, pattharbaji huee aur uske baad tiyar gais chhode gae"} (\texttt{[15:43]}) and \texttt{"aanand mishra ne bataya ki kaise upadraviyon ne police valon ko nishana banaya... police ka kaam 24 * 7 janta ki seva mein khade rahna hai"} (\texttt{[21:20]} \& \texttt{[23:39]}).
  \item \textbf{The shared feature(s) most clearly illustrated were:} \textbf{Naming} (evolving across the bulletin from \texttt{"pradarshankariyon"} \texttt{[00:14]} to \texttt{"asamajik tatva"} \texttt{[15:52]} to \texttt{"upadraviyon"} \texttt{[21:20]}), \textbf{Agency} (passivizing tear-gas and police force as \texttt{"karvaee ki gaee"} / \texttt{"tiyar gais chhode gae"} while giving active verbs and 3 interviews to injured police), and \textbf{Softening} (citing \texttt{"sootron"} that government dialogue doors remain open at \texttt{[09:56]}, \texttt{[17:07]}).
  \item \textbf{An alternative interpretation of the same evidence is:} That the protest was driven not by UP electoral calculations or \texttt{"upadraviyon"}, but by desperate NEET students ignored for a month by the Union Government (\texttt{[08:09--09:48]}, \texttt{[13:34]}), and that unbadged police/plainclothes squads triggered the escalation through excessive force including pellet guns and lathi-charges (\texttt{[06:19--07:08]}).
  \item \textbf{What I cannot conclude from this transcript alone:} What protesting students themselves experienced or thought on 20--21 July (since zero students were interviewed), whether police used pellet guns or unbadged civil-dress personnel (mentioned only in Akhilesh Yadav's clip at \texttt{[06:50--07:08]}), or who threw the first stone at the barricades.
\end{itemize}
\textit{Research reminder: interpret wording in context, distinguish allegations from established facts, and support claims with exact excerpts. Mixed or contradictory evidence is worth recording.}


\clearpage

\begin{center}
    {\large\bfseries LANGUAGE \& SOCIETY PROJECT}\\[3pt]
    {\LARGE\bfseries Transcript Analysis Worksheet}\\[3pt]
    {\normalsize Qualitative analysis of political news discourse}
\end{center}
\vspace{2pt}
\noindent\rule{\textwidth}{0.4pt}
\vspace{4pt}


\section*{Creator 5: Republic TV (\textit{The Debate With Arnab Goswami} --- Nationalist / Pro-Establishment English TV News)}
\begin{itemize}[leftmargin=1.5em, itemsep=3pt, topsep=4pt]
  \item \textbf{Analyst:} Baani Solanki (\texttt{2025114016})
  \item \textbf{Channel / Outlet:} Republic TV (\texttt{@RepublicWorld}, Republic Media Network)
  \item \textbf{Ideological / Editorial Positioning:} Hyper-nationalist, pro-BJP/Hindutva English prime-time broadcast television; deploys prosecutorial anchor monologues ("putting the opposition in the dock"), moral panics, and national-security binaries (\texttt{"Nationalists of Bharat"} vs \texttt{"Deep State"} / \texttt{"rabid, bloodthirsty mobs"} / \texttt{"Taliban"} mindset) to insulate the ruling establishment and counter-attack critics.
  \item \textbf{Transcript Files Referenced:}
  \begin{itemize}[leftmargin=1.4em, itemsep=2pt, topsep=2pt]
    \item Event 1: \texttt{transcripts/event1\_ram\_mandir/05\_republic\_tv\_0H-bruoEkaY.txt}
    \item Event 2: \texttt{transcripts/event2\_neet\_protest/05\_republic\_tv\_HTvb6E\_1Qdc.txt}
  \end{itemize}
\end{itemize}
\vspace{4pt}\noindent\rule{\textwidth}{0.4pt}\vspace{4pt}

\section*{Guidelines: Five features to look for in every event}
\begin{longtable}{@{}>{\raggedright\arraybackslash}p{0.17\textwidth} >{\raggedright\arraybackslash}p{0.20\textwidth} >{\raggedright\arraybackslash}p{0.24\textwidth} >{\raggedright\arraybackslash}p{0.32\textwidth}@{}}
\toprule
\textbf{Feature} & \textbf{What to look for} & \textbf{Illustrative example} & \textbf{Questions to answer} \\
\midrule
\endhead
\textbf{1. Naming} & What are people, groups, actions, or events called? & "Students", "protesters", "agitators"; "irregularity", "scam", "controversy". & What exact label is used? Who or what does it refer to? Who uses it? What alternative label could have been used? \\
\addlinespace[4pt]
\textbf{2. Presupposition} & What is treated as already true or taken for granted? & "Why are students creating chaos again?" may presuppose earlier disruption. & What does the wording assume? Is that assumption supported, disputed, or left unexamined? Does it appear in a question or words such as "again", "still", or "finally"? \\
\addlinespace[4pt]
\textbf{3. Agency} & Who is shown doing something? Who is acted upon? Who disappears? & "Police injured protesters" vs. "Protesters were injured." & Who is the actor? Who experiences the action? Is responsibility named, backgrounded, or omitted? \\
\addlinespace[4pt]
\textbf{4. Antagonisation} & Who is positioned against whom? & "Ordinary citizens" vs. "political troublemakers". & Which groups are contrasted? Is one presented as more legitimate, moral, patriotic, dangerous, or disloyal? Is an "us vs. them" division implied? \\
\addlinespace[4pt]
\textbf{5. Intensification / softening} & What is exaggerated, minimised, softened, or euphemised? & "Violent mob" vs. "peaceful demonstrators"; "incident" instead of a direct description. & Which words heighten emotion or seriousness? Which minimise it? Are metaphors, dramatic adjectives, hedges, or euphemisms used? Does the wording match the evidence? \\
\bottomrule
\end{longtable}
\textit{Reminder: A label or implication is not automatically false or manipulative. Examine the surrounding context and evidence.}

\vspace{4pt}\noindent\rule{\textwidth}{0.4pt}\vspace{4pt}

\section*{Event1. Ram Mandir donation controversy}
\textbf{Focus:} How is the controversy framed---as financial accountability, religious sentiment, political motivation, or something else? Be precise about the specific allegation or episode covered.

\begin{itemize}[leftmargin=1.5em, itemsep=3pt, topsep=4pt]
  \item \textbf{Outlet:} Republic TV (\textit{The Debate With Arnab Goswami})
  \item \textbf{Date:} \texttt{2026-06-26} \textit{(the exact day 8 Trust/counting staff were arrested and Champat Rai \& Anil Mishra resigned)}
  \item \textbf{Video title:} \textit{Debate With Arnab: Yogi's Zero Tolerance For Ram Mandir Donation Money Scamsters} (\url{https://youtu.be/0H-bruoEkaY?si=T9EqDVN39KNILYKg})
  \item \textbf{Timestamp(s):} Relevant Event 1 segment \texttt{[00:54--54:06]} (total duration \texttt{54:30}; brief sponsor jingles \texttt{[00:08--00:53]} and \texttt{[54:12--54:30]} excluded per project guidelines).
  \item \textbf{Analyst:} Baani Solanki (\texttt{2025114016})
  \item \textbf{Language:} English (Native English broadcast audio, Roman script; non-AI-dubbed)
\end{itemize}

\subsection*{Event-specific questions (Event 1)}
\begin{enumerate}[leftmargin=1.6em, itemsep=4pt, topsep=4pt]
  \item \textbf{What terms are used: "scam", "alleged irregularity", "corruption", "political controversy", or something else?}
  \begin{itemize}[leftmargin=1.4em, itemsep=2pt, topsep=2pt]
    \item Arnab Goswami concedes the basic offence in bounded terms---\texttt{"theft from the coffers of the Bhavya Raander [Ram Mandir]"} (\texttt{[00:54--01:02]}), \texttt{"irregularities"} (\texttt{[03:42]}), \texttt{"corruption by low-level functionaries"} (\texttt{[03:22--03:33]}), and \texttt{"donation money scamsters"} (\texttt{[Title]})---while explicitly denying systemic scope: \texttt{"there is no religious angle, there is no political angle, there is no government angle in this controversy"} (\texttt{[04:00--04:08]}).
  \end{itemize}
  \item \textbf{How are the temple, trust, donors, and people making allegations described?}
  \begin{itemize}[leftmargin=1.4em, itemsep=2pt, topsep=2pt]
    \item \textbf{Temple \& State Leadership:} Revered as \texttt{"Bhavya Raander"} (\texttt{[01:02]}); CM Yogi Adityanath is hailed as \texttt{"Yogiad who's proved that there is zero tolerance for any form of corruption"} (\texttt{[03:22]}, \texttt{[09:00]}).
    \item \textbf{Trust \& Arrested Staff:} The 8 arrested men (\texttt{Avinash Shukla}, \texttt{Lavkush Mishra}, \texttt{Anukalp Mishra}, \texttt{Manish Kumar Yadav}, \texttt{Karunesh Pandey}, \texttt{Ram Shankar Mishra}, \texttt{Ram Shankar alias Tinnu Yadav}, \texttt{Subhash Srivastava}) are repeatedly ring-fenced as \texttt{"low-level functionaries in the Ramand trust, some cashiers"} (\texttt{[03:22--03:33]}, \texttt{[07:14]}, \texttt{[21:09]}, \texttt{[32:30--32:41]}). General Secretary Champat Rai (who resigned that day, \texttt{[35:17--35:25]}) is defended during the middle 30 minutes as a \texttt{"79year-old prominent leader of the Vishvah Hindu Parishad... a man who's worked tirelessly"} (\texttt{[20:23--21:28]}), though later at \texttt{[26:47--27:05]} and \texttt{[39:35--40:16]} Arnab briefly questions whether Champat Rai was \texttt{"not skilled administratively"} or given \texttt{"an honorable exit"}.
    \item \textbf{Donors:} Devotees who \texttt{"throng the Aayodhya Ram mandr"} (\texttt{[06:02]}) and \texttt{"contributed, fought for 500 years"} (\texttt{[35:52--36:02]}), exemplified in VTR by a devotee proclaiming \texttt{"I will never lose faith in Ram because of some selfish people"} (\texttt{[09:19--09:30]}).
    \item \textbf{People Making Allegations:} Demonized en masse in the opening monologue as \texttt{"People who say let's divide Hindus... People who did not want the Ram mandr itself... People who said Lord Ram did never exist"} (\texttt{[01:23--01:52]}), \texttt{"don't have the asta [faith] to visit the Ram temple"} (\texttt{[04:18]}), and heirs to Mulayam Singh Yadav's \texttt{"GENOCIDE OF CAR SEAKS"} (\texttt{[04:47--05:06]}).
  \end{itemize}
  \item \textbf{Is the issue framed as financial accountability, an attack on Hindu faith, a political dispute, or a matter of public trust?}
  \begin{itemize}[leftmargin=1.4em, itemsep=2pt, topsep=2pt]
    \item Framed overwhelmingly as \textbf{swift BJP/Yogi administrative triumph (\texttt{"Zero Tolerance"}) under opportunistic political attack by anti-Hindu opposition hypocrites}. Arnab explicitly announces his prosecutorial pivot at \texttt{[05:14--05:24]}: \texttt{"I believe action has been taken, but since they have politicized the matter I suggest the best response is to put them in the dock and ask them what their record is on Hindutva"}.
  \end{itemize}
  \item \textbf{Are questions about donations presented as legitimate scrutiny or as an attempt to malign the temple or its supporters?}
  \begin{itemize}[leftmargin=1.4em, itemsep=2pt, topsep=2pt]
    \item Opposition questions are framed as \textbf{hypocritical attempts to malign the temple and divide Hindus} (\texttt{"looking at this opportunity to target the Ramander itself"}, \texttt{[01:23]}; \texttt{"want to fish in troubled waters"}, \texttt{[23:01]}). Only late in the 54-minute show (\texttt{[39:35--40:26]}, \texttt{[53:10--53:19]}) does Arnab concede that \texttt{"phase by phase, week by week... details must come out"} because devotees were hurt.
  \end{itemize}
  \item \textbf{Are allegations attributed to specific people and supported with evidence, or presented as established facts?}
  \begin{itemize}[leftmargin=1.4em, itemsep=2pt, topsep=2pt]
    \item Official SIT findings---\texttt{8} named arrests, \texttt{"over 80 lakh rupees may have been stolen"}, and \texttt{"CCTV footage reportedly showing the accused taking cash at least 70 times"} (\texttt{[06:48--07:34]})---are presented with cautious journalistic hedging (\texttt{may have been}, \texttt{reportedly}). Conversely, whenever opposition panelists allege wider BJP/Trust complicity (e.g., Priyanka Kakkar citing Raju Manwani's \texttt{200 kg} silver bricks \& Champat Rai refusing PM summons at \texttt{[14:46--15:23]}; Major Ashish Chaturvedi alleging BJP leaders' involvement and \texttt{2020} private audit findings at \texttt{[27:34--28:14]}, \texttt{[31:25]}), Arnab aggressively demands courtroom-grade proof on the spot (\texttt{"NAME ONE... VALIDATE IT RIGHT NOW"}, \texttt{[31:43--32:10]}).
  \end{itemize}
  \item \textbf{How are denials and counterarguments represented: explained, dismissed, challenged, or omitted?}
  \begin{itemize}[leftmargin=1.4em, itemsep=2pt, topsep=2pt]
    \item \textbf{State/BJP defenses are amplified as the show's headline thesis}, with BJP's Nalin Kohli and Shantanu Gupta given extended turns to praise \texttt{"breakneck speed"} arrests (\texttt{[12:47]}, \texttt{[48:20--48:48]}, \texttt{[49:37--50:40]}). Meanwhile, \textbf{opposition counterarguments are derailed for over 20 minutes (\texttt{[16:50--35:00]} \& \texttt{[46:18--48:02]})} when AAP spokesperson Priyanka Kakkar makes a hyperbolically phrased demand (\texttt{"Why are you shielding Champat and the bigger chors? Why should they not be publicly executed?"}, \texttt{[16:50]})---which Arnab and Nalin Kohli seize upon as \texttt{"BREAKING NEWS"}, \texttt{"TALIBAN"}, \texttt{"lynch mob attitude"}, and \texttt{"PART OF THE TALIBAN SHOOTING SQUAD"} (\texttt{[17:13--22:35]}, \texttt{[33:17--33:37]}) to drown out her preceding factual questions about missing silver receipts and unmade public SIT reports!
  \end{itemize}
  \item \textbf{Does the presenter distinguish criticism of the trust or its financial practices from criticism of the temple or religion itself?}
  \begin{itemize}[leftmargin=1.4em, itemsep=2pt, topsep=2pt]
    \item \textbf{Rhetorically conflates opposition criticism of the Trust with hatred of the Temple and Hinduism} for most of the broadcast (\texttt{"looking at this opportunity to target the Ramander itself"}, \texttt{[01:23]}; Nalin Kohli equating Priyanka Kakkar's attack on Champat Rai with \texttt{"contempt for Lord Ram"} and \texttt{"EXECUTE THE RAM BHAKTS, EXECUTE THE SANATANIS"} at \texttt{[24:51--25:21]}), while reserving for \textit{himself} and pro-BJP lawyer Ishkaran Bhandari (\texttt{[35:44--36:56]}, \texttt{[39:35--40:16]}) the right to distinguish administrative failure from faith.
  \end{itemize}
  \item \textbf{Who is portrayed as protecting donors' interests, and who as exploiting religious sentiment?}
  \begin{itemize}[leftmargin=1.4em, itemsep=2pt, topsep=2pt]
    \item \textbf{Protecting donors' interests:} UP CM Yogi Adityanath (\texttt{[01:12]}, \texttt{[03:22]}, \texttt{[09:00]}), the Trust itself (claimed to have written the letter requesting action, \texttt{[03:42--03:51]}), and the three-member IAS/IPS/Finance SIT (\texttt{[48:29--48:39]}).
    \item \textbf{Exploiting religious sentiment:} Arvind Kejriwal (mocked repeatedly for his \texttt{"nani told me"} Babri Masjid remark, \texttt{[02:01--02:28]}, \texttt{[10:34--11:10]}, \texttt{[47:24--47:34]}), Akhilesh Yadav, and Rahul Gandhi (\texttt{[02:36--05:14]}).
  \end{itemize}
\end{enumerate}

\subsection*{Apply the five shared features (Event 1)}
\begin{longtable}{@{}>{\raggedright\arraybackslash}p{0.19\textwidth} >{\raggedright\arraybackslash}p{0.38\textwidth} >{\raggedright\arraybackslash}p{0.37\textwidth}@{}}
\toprule
\textbf{Feature} & \textbf{Exact wording + timestamp} & \textbf{What it does in context} \\
\midrule
\endhead
\textbf{1. Naming} & Title: \texttt{"Yogi's Zero Tolerance For Ram Mandir Donation Money Scamsters"} \newline  \texttt{"low-level functionaries in the Ramand trust, some cashiers"} \texttt{[03:22--03:33]} vs \texttt{"79year-old prominent leader of the Vishvah Hindu Parishad"} \texttt{[20:23]} & Minimizes the 8 arrested insiders by rank (\texttt{low-level functionaries, some cashiers}) to sever upward culpability, while venerating Champat Rai via age and lifelong religious service (\texttt{79-year-old prominent leader}). \\
\addlinespace[4pt]
\textbf{2. Presupposition} & \texttt{"what hypocrisy people like Akileles and Rahul Gandhi who don't even want to visit the Ram mantr... are questioning Yogiad who's proved that there is zero tolerance for any form of corruption"} \texttt{[03:04--03:22]} & Factive verb \texttt{"who's proved"} presupposes as indisputable truth that Yogi's government has zero tolerance for corruption, while exclamative \texttt{"what hypocrisy"} presupposes opposition bad faith. \\
\addlinespace[4pt]
\textbf{3. Agency} & \texttt{"there has been theft from the coffers... But it is also true that there has been very strong and effective action... immediate action by the Yogi government, the BJP government"} \texttt{[00:54--01:12]} & Uses an agentless existential construction for the crime (\texttt{there has been theft}---no BJP/Trust agent named), immediately followed by explicit agentive credit (\texttt{by the Yogi government, the BJP government}) for the arrests. \\
\addlinespace[4pt]
\textbf{4. Antagonisation} & \texttt{"People who say let's divide Hindus... People who did not want the Ram mandr itself... GENOCIDE OF CAR SEAKS"} \texttt{[01:23--01:52]}, \texttt{[04:56]} \newline  \texttt{"are you sponsored by the Taliban... PART OF THE TALIBAN SHOOTING SQUAD?"} \texttt{[33:17--33:37]} & Constructs an existential civilizational enemy (\texttt{anti-Hindu dividers}, \texttt{kar-sevak killers}, \texttt{Taliban shooting squad}) so viewers perceive opposition critics of the scam as enemies of Hindu society itself. \\
\addlinespace[4pt]
\textbf{5. Intensification / softening} & \textbf{Softening (scam size \& scope):} \texttt{"over 80 lak rupees may have been stolen"} \texttt{[07:26]}, \texttt{"there is no political angle, there is no government angle"} \texttt{[04:08]}, \texttt{"protecting him a wee bit"} \texttt{[40:45]} \newline  \textbf{Intensification (panelist slip):} \texttt{"MY GOD, THIS IS TALIBAN... UNFORGIVABLE"} \texttt{[17:20]}, \texttt{[27:05]} & Hedges and minimizes the financial theft (\texttt{may have been}, \texttt{80 lakh}, \texttt{a wee bit}), while blowing up a panelist's rhetorical hyperbole (\texttt{publicly executed}) into a 20-minute \texttt{"Taliban"} moral panic. \\
\bottomrule
\end{longtable}

\subsection*{Short summary (Event 1)}
\begin{itemize}[leftmargin=1.5em, itemsep=3pt, topsep=4pt]
  \item \textbf{What interpretation of the event does this video encourage?} That an unfortunate \texttt{Rs.\,80+ lakh} cash-counting theft by 8 \texttt{"low-level"} employees (\texttt{[03:33]}, \texttt{[07:26]}) has already been decisively resolved by CM Yogi Adityanath's \texttt{"zero tolerance"} arrests and senior resignations (\texttt{[01:12]}, \texttt{[35:25]}), and that the \textit{real} scandal is anti-Ram opposition leaders (\texttt{Kejriwal}, \texttt{Akhilesh}, \texttt{Rahul}) hypocritically exploiting the episode ahead of UP elections (\texttt{[01:23--05:24]}, \texttt{[09:39]}).
  \item \textbf{Whose perspective is foregrounded, challenged, or missing?} \textit{Foregrounded:} Arnab Goswami's pro-Yogi prosecutorial frame, BJP spokesperson Nalin Kohli, Yogi biographer Shantanu Gupta, and pro-BJP lawyer Ishkaran Bhandari. \textit{Challenged:} AAP's Priyanka Kakkar (relentlessly hectored over \texttt{"public execution"}, \texttt{[17:09--35:00]}), SP-leaning Major Ashish Chaturvedi (\texttt{[31:25--32:58]}), and CPI/Left panelist Vivek Srivastava (ambushed about China/Soviet funds at \texttt{[38:12--39:16]}). \textit{Missing:} Unreceipted donors, whistleblower Mahipal Singh, and any scrutiny of the 2021 Ayodhya land deals (cut off at \texttt{[45:22]}).
  \item \textbf{What evidence supports your interpretation? What complicates it?} \textit{Supports:} Verifiable registration of the FIR, arrest of 8 named individuals within days (\texttt{[06:57--07:14]}), CCTV documentation of 70 cash-theft instances (\texttt{[07:26--07:34]}), and same-day resignations of Champat Rai and Anil Mishra (\texttt{[35:25]}). \textit{Complicates:} Even Arnab himself is forced at \texttt{[39:47--40:16]} to admit \texttt{"I refuse to believe that this kind of an organized theft which has happened over many months and years could have happened only by the cash counting staff"}---directly contradicting his own opening monologue (\texttt{[03:22--04:08]})!
\end{itemize}
\vspace{4pt}\noindent\rule{\textwidth}{0.4pt}\vspace{4pt}

\section*{Event2. Students' protest regarding NEET paper leak}
\textbf{Focus:} How does coverage construct the legitimacy of student protest and responsibility for the underlying problem?

\begin{itemize}[leftmargin=1.5em, itemsep=3pt, topsep=4pt]
  \item \textbf{Outlet:} Republic TV (\textit{The Debate / Arnab's Take})
  \item \textbf{Date:} \texttt{2026-07-25} \textit{(the exact evening Education Minister Dharmendra Pradhan resigned and CJP called off the protest)}
  \item \textbf{Video title:} \textit{Rahul The Biggest Loser: Arnab Unravels The Political Agenda Behind CJP's Violent Protest} (\url{https://youtu.be/HTvb6E_1Qdc?si=7aFf_zysFjpvX8af})
  \item \textbf{Timestamp(s):} Relevant Event 2 segment \texttt{[00:00--06:41]} (\texttt{100\%} direct-to-camera monologue by Arnab Goswami; total duration \texttt{06:41}).
  \item \textbf{Analyst:} Baani Solanki (\texttt{2025114016})
  \item \textbf{Language:} English with strategic Hindi insertions (\texttt{khali}, \texttt{saboot}, \texttt{sansad ukhad denge}) (Native English audio, Roman script; non-AI-dubbed)
\end{itemize}

\subsection*{Event-specific questions (Event 2)}
\begin{enumerate}[leftmargin=1.6em, itemsep=4pt, topsep=4pt]
  \item \textbf{Are participants called "students", "candidates", "protesters", "agitators", "mobs", or "political pawns"?}
  \begin{itemize}[leftmargin=1.4em, itemsep=2pt, topsep=2pt]
    \item \textbf{Systematically stripped of student identity and criminalized in extreme terms:} Arnab explicitly declares \texttt{"these were not students, you were fooled... they were not students who did this, the students had evaporated"} (\texttt{[04:58--05:26]}). Instead, he labels them \texttt{"rabid, bloodthirsty mobs"} (\texttt{[01:47]}), \texttt{"the mob"} (\texttt{[01:00]}, \texttt{[03:14]}), \texttt{"the Lynch mob of Abhijit Dipke"} (\texttt{[03:54]}), \texttt{"the goons [goods] of Abhijit Dipke"} (\texttt{[04:19]}, \texttt{[06:21]}), \texttt{"Deep State Dipke"} (\texttt{[00:23]}), and \texttt{"Sonam Wangchuk, very questionable man"} (\texttt{[00:33]}). \textit{(Note: YouTube ASR twice mishears \texttt{"CJP"} as \texttt{"BJP"} at \texttt{[00:07]}, \texttt{[00:51]}, \texttt{[06:01]}---in context Arnab clearly says CJP leadership called off the protest and left the ground.)}
  \end{itemize}
  \item \textbf{Are their stated demands explained, summarised, mocked, questioned, or left out?}
  \begin{itemize}[leftmargin=1.4em, itemsep=2pt, topsep=2pt]
    \item \textbf{Reduced to a brief transactional opening acknowledgement, then twisted into regime conspiracy:} In the first 16 seconds (\texttt{[00:00--00:15]}), Arnab notes \texttt{"Dharmendra Pradhan has resigned. That demand has been met"} plus an assurance against cases for Dipke and Saurabh Das---only to claim for the remaining 6 minutes that the real hidden agenda was \texttt{"streetocracy"} (\texttt{[00:42--00:51]}), \texttt{"sansad ukhad denge, Nepal model"} regime change (\texttt{[05:08--05:16]}), \texttt{"conversion of Hindus"}, and blocking the \texttt{FCRA}, \texttt{delimitation}, and \texttt{one nation, one poll} bills (\texttt{[02:38--02:56]})!
  \end{itemize}
  \item \textbf{Does the presenter focus on the original grievance, disruption, political affiliations, or public inconvenience?}
  \begin{itemize}[leftmargin=1.4em, itemsep=2pt, topsep=2pt]
    \item \textbf{Zero words on the NEET paper leak itself!} The word \texttt{"NEET"} or \texttt{"paper leak"} does not appear \textit{even once} in the spoken transcript (\texttt{[00:00--06:41]}). 100\% of the monologue focuses on (a) mocking Rahul Gandhi (\texttt{[00:00--02:56]}) and (b) inciting permanent nationalist fury over street violence against police/RAF (\texttt{[03:06--06:38]}).
  \end{itemize}
  \item \textbf{Who is represented as responsible: examination authorities, government, students, political parties, or others?}
  \begin{itemize}[leftmargin=1.4em, itemsep=2pt, topsep=2pt]
    \item \textbf{Zero responsibility is attributed to the NTA, Dharmendra Pradhan, or the government.} Instead, blame is placed entirely on \texttt{Rahul Gandhi}, \texttt{Arvind Kejriwal}, \texttt{Yogendra Yadav} (\texttt{[00:33--00:51]}), \texttt{"Deep State Dipke"} (\texttt{[00:23]}, \texttt{[04:19--05:43]}), and \texttt{"questionable"} \texttt{Sonam Wangchuk} (\texttt{[00:33]}).
  \end{itemize}
  \item \textbf{Are questions framed around whether students have legitimate grievances, or why they are causing disruption?}
  \begin{itemize}[leftmargin=1.4em, itemsep=2pt, topsep=2pt]
    \item Framed exclusively around \textbf{opposition frustration and punitive retribution against protest organisers}: \texttt{"what exactly is your problem, Rahul Gandhi? Why do you sound so frustrated?"} (\texttt{[01:00--01:19]}) and \texttt{"Dharmendra Pradhan has resigned, when will Abhijit Dipke be held to account now?"} (\texttt{[05:34]}).
  \end{itemize}
  \item \textbf{Are claims of violence, disorder, or political manipulation attributed to sources and substantiated?}
  \begin{itemize}[leftmargin=1.4em, itemsep=2pt, topsep=2pt]
    \item \textbf{Presented as unattributed, emotive assertions coupled with extreme geopolitical/communal leaps:} While referencing visual clips (\texttt{"Do not forget these pictures"}, \texttt{[03:26]}, \texttt{[04:03]}) and naming injured RAF officer \texttt{Sonia Sherawat} (\texttt{[03:44--04:11]}), Arnab offers \textbf{zero evidence} for \texttt{"Deep State Dipke"} (\texttt{[00:23]}), \texttt{"Nepal model"} overthrow (\texttt{[05:16]}), or the bizarre claim at \texttt{[02:38--02:48]} that protesters/opposition \texttt{"didn't want the FCRA bill to be brought in. They wanted the conversion of Hindus"}!
  \end{itemize}
  \item \textbf{How are police and government actions described compared with protesters' actions? Are similar actions evaluated differently?}
  \begin{itemize}[leftmargin=1.4em, itemsep=2pt, topsep=2pt]
    \item \textbf{Total inversion of victim and perpetrator:} The government is praised five times in a row via anaphora (\texttt{"the government outplayed them"}, \texttt{[02:06--02:56]}), and \texttt{"our police, our RAF, our security personnel"} (\texttt{[01:37--01:47]}, \texttt{[03:36--04:11]}) are portrayed purely as heroic victims of an \texttt{"attempt to lynch"}, while \textbf{100\% of police violence} (lathi-charges, 25+ tear-gas shells, pellet guns, Wangchuk's forced removal, unbadged men) is \textbf{completely erased}. Protesters' actions are characterized as \texttt{"absolute anarchy"}, \texttt{"drunken violence"}, and \texttt{"sexual perversion"} (\texttt{[01:37]}, \texttt{[04:46--04:58]}).
  \end{itemize}
  \item \textbf{Are students allowed to explain their demands in their own words, or does the presenter mainly speak for them?}
  \begin{itemize}[leftmargin=1.4em, itemsep=2pt, topsep=2pt]
    \item \textbf{Zero student voices.} Arnab Goswami speaks uninterrupted for all \texttt{06:41}, explicitly denying that genuine students were even present (\texttt{"the students had evaporated"}, \texttt{[05:26]}).
  \end{itemize}
  \item \textbf{Does the outlet's framing differ between coverage of the NEET paper-leak issue before the protest and coverage once students begin protesting?}
  \begin{itemize}[leftmargin=1.4em, itemsep=2pt, topsep=2pt]
    \item \textbf{Complete erasure of the educational issue once street protests challenge the government:} Even on the day the Union Education Minister resigns over NEET (\texttt{[00:00]}), Arnab never mentions exam leaks, student suicides, or NTA reform---turning a massive government concession (\texttt{Dharmendra Pradhan has resigned}) into a triumphalist \texttt{"Rahul The Biggest Loser"} victory lap paired with a call for street counter-mobilization (\texttt{[06:10--06:38]}).
  \end{itemize}
\end{enumerate}

\subsection*{Apply the five shared features (Event 2)}
\begin{longtable}{@{}>{\raggedright\arraybackslash}p{0.19\textwidth} >{\raggedright\arraybackslash}p{0.38\textwidth} >{\raggedright\arraybackslash}p{0.37\textwidth}@{}}
\toprule
\textbf{Feature} & \textbf{Exact wording + timestamp} & \textbf{What it does in context} \\
\midrule
\endhead
\textbf{1. Naming} & \texttt{"Deep State Dipke"} \texttt{[00:23]}; \texttt{"Sonam Wangchuk, very questionable man"} \texttt{[00:33]}; \texttt{"streetocracy"} \texttt{[00:42]}; \texttt{"rabid, bloodthirsty mobs"} \texttt{[01:47]}; \texttt{"Lynch mob / goons of Abhijit Dipke"} \texttt{[03:54]}, \texttt{[04:19]} & Dehumanizes protesters with animalistic disease imagery (\texttt{rabid, bloodthirsty}), conspiracy alliteration (\texttt{Deep State Dipke}), and anti-democratic coinage (\texttt{streetocracy}), erasing student citizenship. \\
\addlinespace[4pt]
\textbf{2. Presupposition} & \texttt{"Arnab Unravels The Political Agenda Behind CJP's Violent Protest"} \texttt{[Title]} \newline  \texttt{"Why are you, as Gen Z would say, so shook tonight?"} \texttt{[01:19]} \newline  \texttt{"these were not students, you were fooled"} \texttt{[04:58]} & Definite NP \texttt{"The Political Agenda Behind CJP's Violent Protest"} presupposes both hidden political orchestration and inherent protest violence; \texttt{"you were fooled"} presupposes viewer deception if they felt sympathy for students. \\
\addlinespace[4pt]
\textbf{3. Agency} & \texttt{"our police, our RAF, our security personnel were vehemently attacked by rabid, bloodthirsty mobs, comes a resignation"} \texttt{[01:37--01:56]} & Passivizes security forces purely as patients/victims (\texttt{were vehemently attacked}), deletes all police use of force, and nominalizes Pradhan's forced exit as an agentless motion verb (\texttt{comes a resignation}---avoiding \textit{"CJP forced the government to sack Pradhan"}!). \\
\addlinespace[4pt]
\textbf{4. Antagonisation} & \texttt{"when this violence was being carried out by the goons of Abhijit Dipke, our brave soldiers were fighting terrorists IN ANANTNAG... Nationalists, mobilize yourself... This nation is ours."} \texttt{[04:19--04:37]}, \texttt{[06:21--06:38]} & Juxtaposes Jantar Mantar protesters with \texttt{"terrorists in Anantnag"} killing Indian soldiers (compare Operation Sindoor callback at \texttt{[01:19--01:37]}), and ends with an exclusionary majoritarian call to arms (\texttt{Nationalists, mobilize yourself... This nation is ours}). \\
\addlinespace[4pt]
\textbf{5. Intensification / softening} & \textbf{Anaphoric intensification:} 5x \texttt{"the government outplayed them"} \texttt{[02:06--02:56]}; 9x \texttt{"Do not forget / Do not forgive"} \texttt{[03:06--06:21]}; \texttt{"rabid, bloodthirsty"}, \texttt{"drunken violence"}, \texttt{"sexual perversion"} \texttt{[01:47]}, \texttt{[04:46--04:58]} \newline  \textbf{Softening (govt defeat):} \texttt{"comes a resignation"} \texttt{[01:47]} & Liturgical, drumbeat repetition (\texttt{Do not forgive, do not forget}) and visceral disgust triggers (\texttt{rabid}, \texttt{drunken}, \texttt{perversion}) whip up enduring grievance, while \texttt{"comes a resignation"} + \texttt{"outplayed"} reframes ministerial capitulation into strategic genius. \\
\bottomrule
\end{longtable}

\subsection*{Short summary (Event 2)}
\begin{itemize}[leftmargin=1.5em, itemsep=3pt, topsep=4pt]
  \item \textbf{What interpretation of the event does this video encourage?} That the CJP protest had nothing to do with NEET students (\texttt{"the students had evaporated"}, \texttt{[05:26]}), was a \texttt{"Deep State"} / Rahul Gandhi \texttt{"streetocracy"} plot (\texttt{[00:23--00:51]}) to unleash \texttt{"rabid, bloodthirsty mobs"} against police/RAF and block nationalist bills like \texttt{FCRA} (\texttt{[01:47]}, \texttt{[02:38--02:56]}), that the Modi government brilliantly \texttt{"outplayed"} the opposition by ending the protest via Pradhan's resignation (\texttt{[02:06--02:56]}), and that \texttt{"Nationalists of Bharat"} must \texttt{"never forgive and never forget"} Dipke's \texttt{"goons"} (\texttt{[03:06--06:38]}).
  \item \textbf{Whose perspective is foregrounded, challenged, or missing?} \textit{Foregrounded:} Exclusively Arnab Goswami's militant nationalist perspective, aligned with the Modi government and injured RAF personnel (\texttt{Sonia Sherawat}, \texttt{[03:44--04:11]}). \textit{Challenged/Demonized:} Rahul Gandhi, Arvind Kejriwal, Yogendra Yadav, Abhijit Dipke, Sonam Wangchuk, and all Jantar Mantar protesters. \textit{Missing:} \textbf{100\% of student grievances}, NEET paper leaks, student suicides, MEA's own denial of foreign funding (cited by Sarthak \texttt{EM1GeZufO6E} \texttt{[00:40]}), and all police lathi/tear-gas/pellet violence!
  \item \textbf{What evidence supports your interpretation? What complicates it?} \textit{Supports:} Confirmed resignation of Dharmendra Pradhan and emptying of Jantar Mantar (\texttt{[00:00--00:23]}), plus real injuries sustained by police/RAF personnel during the 20 July clashes (\texttt{[03:14--04:11]}, corroborated by Aaj Tak \texttt{JXLIs0WtJAM} \texttt{[18:09]}). \textit{Complicates:} An obvious internal logical contradiction---if protesters were truly \texttt{"Deep State"} \texttt{"rabid, bloodthirsty mobs"} and \texttt{"not students"}, why did the Union Government accept their demands, remove its Education Minister, and promise \textit{not} to prosecute Dipke and Saurabh Das (\texttt{[00:00--00:15]})? Furthermore, MEA officially denied foreign funding (\texttt{EM1GeZufO6E}, \texttt{[00:40--00:48]}), refuting \texttt{"Deep State Dipke"}.
\end{itemize}
\vspace{4pt}\noindent\rule{\textwidth}{0.4pt}\vspace{4pt}

\section*{Common form: Blank excerpt record}
\textit{Copy or print this page whenever you find a passage worth analysing. Keep the exact wording and enough context to interpret it.}


\subsection*{Excerpt Record 1 (Event 1 --- Opening Monologue: Minimizing Theft \& Putting Critics "In the Dock")}
\begin{itemize}[leftmargin=1.5em, itemsep=3pt, topsep=4pt]
  \item \textbf{Outlet:} Republic TV (\textit{The Debate}) | \textbf{Date:} \texttt{2026-06-26}
  \item \textbf{Video title:} \textit{Debate With Arnab: Yogi's Zero Tolerance For Ram Mandir Donation Money Scamsters}
  \item \textbf{Timestamp(s):} \texttt{[00:54--01:52]} \& \texttt{[03:22--05:24]} | \textbf{Analyst:} Baani Solanki | \textbf{Language:} English (Roman script)
  \item \textbf{Exact transcript excerpt:}
    \begin{quote}\small\ttfamily [00:54] Yes, ladies and gentlemen, there has been theft from the coffers of the [01:02] Bhavya Raander [Ram Mandir]. It is true. It is unfortunate. But it is also true that there has been very strong and effective action, [01:12] sackings and arrests and resignations, immediate action by the Yogi government, the BJP government. But it is also true [01:23] that some people are looking at this opportunity to target the Ramander itself. People who say let's divide Hindus... [03:22] Yogiad who's proved that there is zero tolerance for any form of corruption including the corruption by low-level [03:33] functionaries in the Ramand trust some cashiers... [04:08] there is no religious angle there is no political angle there is no government angle in this controversy... [05:14] since they have politicized the matter I suggest the best response is to put them [05:24] in the dock and ask them what their record is on Hindua [Hindutva].\end{quote}
  \item \textbf{Feature(s):} $\boxtimes$ Naming $\boxtimes$ Presupposition $\boxtimes$ Agency $\boxtimes$ Antagonisation $\boxtimes$ Intensification / softening
  \item \textbf{Who is speaking? Is this the outlet's wording or a quotation?} Anchor Arnab Goswami's scripted opening editorial monologue.
  \item \textbf{What does the wording imply, foreground, soften, or leave unstated?} Uses a classic \textbf{\texttt{Yes... But...} concession-reversal structure}: concedes the undeniable theft in two short existential clauses (\texttt{"there has been theft... It is unfortunate"}), confines culprits to \texttt{"low-level functionaries, some cashiers"}, credits \texttt{"immediate action by the Yogi government"}, and then pivots the entire 54-minute agenda to prosecuting \textit{the opposition} (\texttt{"put them in the dock"}).
  \item \textbf{What context or evidence supports or complicates this interpretation?} Supported by the arrest of 8 staffers on 26 June (\texttt{[06:57--07:14]}). Complicated by the fact that allegations had circulated since 2021 (Mahipal Singh, private audit) and 7 June 2026 (Akhilesh Yadav tweet), and that Champat Rai and Anil Mishra also resigned that very day (\texttt{[35:25]})---hardly \texttt{"some cashiers"}.
  \item \textbf{How does another outlet describe the same event or actor?} \textbf{The Wire} (\texttt{7rUL8fGgRJA}, \texttt{[03:29--04:59]}) frames the exact same event as \texttt{"massive theft or dacoity"} exposing the \texttt{"moral degeneration"} of the \texttt{"entire Sangh Parivar"}, whereas \textbf{Republic TV} declares \texttt{"there is no political angle, there is no government angle"} (\texttt{[04:08]}).
  \item \textbf{My interpretation, with reasons:} Explicitly saying \texttt{"the best response is to put them in the dock"} (\texttt{[05:14--05:24]}) lays bare Republic TV's discursive function: transforming accountability \textit{of} the ruling establishment into an ideological trial \textit{of} the opposition.
\end{itemize}

\subsection*{Excerpt Record 2 (Event 1 --- Manufacturing a "Taliban" Moral Panic Over Panelist's Hyperbole)}
\begin{itemize}[leftmargin=1.5em, itemsep=3pt, topsep=4pt]
  \item \textbf{Outlet:} Republic TV (\textit{The Debate}) | \textbf{Date:} \texttt{2026-06-26}
  \item \textbf{Video title:} \textit{Debate With Arnab: Yogi's Zero Tolerance For Ram Mandir Donation Money Scamsters}
  \item \textbf{Timestamp(s):} \texttt{[16:50--17:28]}, \texttt{[20:12--21:28]} \& \texttt{[33:17--33:37]} | \textbf{Analyst:} Baani Solanki | \textbf{Language:} English (Roman script)
  \item \textbf{Exact transcript excerpt:}
    \begin{quote}\small\ttfamily [16:50] [AAP's Priyanka Kakkar:] But the issue today is why are you shielding champ [Champat Rai] and the bigger chores? Why are why should they not be publicly executed?\end{quote}
    \begin{quote}\small\ttfamily [16:59] [Arnab \& BJP's Nalin Kohli:] publicly executed. You are calling for the public execution... [17:13] breaking news for the public execution [17:20] of Champat Rai... MY GOD, THIS IS TALIBAN...\end{quote}
    \begin{quote}\small\ttfamily [20:23] [Arnab Goswami:] the Aam Aadmi Party has called for the public execution of 79year-old prominent leader of the Vishvah Hindu Parishad... [21:09] There are eight people who have been arrested but they are cash counting staff...\end{quote}
    \begin{quote}\small\ttfamily [33:17] [Arnab Goswami:] you've called for the public execution are you sponsored by [33:26] the Taliban is the Taliban sponsoring you is the Taliban your boss? Are you ideologically driven by the Taliban? ... Are you [33:37] PART OF THE TALIBAN SHOOTING SQUAD?\end{quote}
  \item \textbf{Feature(s):} $\boxtimes$ Naming $\boxtimes$ Presupposition $\square$ Agency $\boxtimes$ Antagonisation $\boxtimes$ Intensification / softening
  \item \textbf{Who is speaking? Is this the outlet's wording or a quotation?} AAP spokesperson Priyanka Kakkar (\texttt{[16:50]}), BJP spokesperson Nalin Kohli (\texttt{[17:20--20:12]}), and anchor Arnab Goswami (\texttt{[17:13]}, \texttt{[20:23--21:28]}, \texttt{[33:17--33:37]}).
  \item \textbf{What does the wording imply, foreground, soften, or leave unstated?} Just before \texttt{[16:50]}, Priyanka Kakkar asked five concrete evidentiary questions (\texttt{[14:03--15:23]}: unpublished SIT report, Brij Bhushan's silence, PM summons to Champat Rai, Raju Manwani's \texttt{200 kg} unreceipted silver). The instant she utters the reckless figurative phrase \texttt{"publicly executed"}, Arnab declares \texttt{"breaking news"}, re-names Champat Rai a vulnerable \texttt{"79year-old prominent leader"}, and hurls five consecutive \texttt{"Taliban"} questions (\texttt{[33:17--33:37]})---ensuring none of her substantive questions ever have to be answered!
  \item \textbf{What context or evidence supports or complicates this interpretation?} Calling for \texttt{"public execution"} on national TV \textit{is} genuinely unconstitutional hyperbole (as Ishkaran Bhandari soberly notes at \texttt{[36:33]}: \texttt{"We do not have Taliban style punishments... as per law all of them need to be... punished"}). Yet weaponizing that 3-word slip for 20+ minutes to accuse an Indian political party of being \texttt{"PART OF THE TALIBAN SHOOTING SQUAD"} is textbook diversionary intensification.
  \item \textbf{How does another outlet describe the same event or actor?} \textbf{Dhruv Rathee} (\texttt{oeH0fOSJuX8}, \texttt{[00:17--01:20]}) foregrounds the exact Sindhi community \texttt{210 kg} silver-brick story that Priyanka Kakkar raised at \texttt{[15:04--15:23]}, which Republic TV buried beneath the \texttt{"Taliban"} shouting match.
  \item \textbf{My interpretation, with reasons:} This exchange illustrates conversational floor-control in television debate: seizing on an opponent's tonal overreach to shift the topic from \textit{temple corruption} to \textit{opposition extremism}.
\end{itemize}

\subsection*{Excerpt Record 3 (Event 1 --- Arnab's Surprising Late-Show Pivot Questioning "Honorable Exit")}
\begin{itemize}[leftmargin=1.5em, itemsep=3pt, topsep=4pt]
  \item \textbf{Outlet:} Republic TV (\textit{The Debate}) | \textbf{Date:} \texttt{2026-06-26}
  \item \textbf{Video title:} \textit{Debate With Arnab: Yogi's Zero Tolerance For Ram Mandir Donation Money Scamsters}
  \item \textbf{Timestamp(s):} \texttt{[39:35--40:26]} | \textbf{Analyst:} Baani Solanki | \textbf{Language:} English (Roman script)
  \item \textbf{Exact transcript excerpt:}
    \begin{quote}\small\ttfamily [39:35] [Arnab Goswami to BJP's Nalin Kohli:] So Nalin, are you protecting Champatra? Question one. Allow me to complete [39:47] the highest administrative level of the trust must be blamed, not the lowest administrative level. I refuse to believe that this kind of an [39:57] organized theft which has happened over many months and years could have happened only by the cash counting staff. [40:05] Therefore, I want to know whether you will call him for questioning, whether his name will be included in the FIR or whether you are trying to give [40:16] him an honorable exit by telling him to quietly resign, right, and lay the blame on others.\end{quote}
  \item \textbf{Feature(s):} $\boxtimes$ Naming $\boxtimes$ Presupposition $\boxtimes$ Agency $\square$ Antagonisation $\boxtimes$ Intensification / softening
  \item \textbf{Who is speaking? Is this the outlet's wording or a quotation?} Anchor Arnab Goswami questioning BJP spokesperson Nalin Kohli directly (\texttt{[39:35--40:26]}).
  \item \textbf{What does the wording imply, foreground, soften, or leave unstated?} Remarkably, 40 minutes into the broadcast, Arnab briefly adopts the \textit{exact institutional critique} voiced by opposition/independent critics (\texttt{"I refuse to believe... could have happened only by the cash counting staff"}, \texttt{"honorable exit"}), before softening it seconds later (\texttt{"seen to be protecting him a wee bit"}, \texttt{[40:45]}) and letting Nalin Kohli reassure viewers that \texttt{"no one... will be spared"} (\texttt{[42:08--42:25]}).
  \item \textbf{What context or evidence supports or complicates this interpretation?} Crucial complicating evidence for our analysis! It shows why careful line-by-line transcript reading matters: even a fiercely pro-government anchor felt genuine Hindu viewer anger over Ram Mandir theft (\texttt{"people have been hurt by it"}, \texttt{[53:19]}) and could not sustain a 100\% uncritical whitewash of Champat Rai for the full hour.
  \item \textbf{How does another outlet describe the same event or actor?} Matches \textbf{Dhruv Rathee} (\texttt{oeH0fOSJuX8}, \texttt{[13:05--13:14]}: \texttt{"Is it enough to resign? ... should they also be arrested?"}) and \textbf{Sarthak Goswami} (\texttt{lb7gZnyZH0c}, \texttt{[00:46]}: \texttt{"tinnoo draivar bali ka bakra"}), though Arnab hedges it (\texttt{"a wee bit"}) and pairs it with praise for CM Yogi.
  \item \textbf{My interpretation, with reasons:} On an issue touching Hindu devotees' sacred donations, Republic TV performs "bounded accountability"---pressing BJP (\texttt{"a wee bit"}) at minute 40 to maintain credibility with devout viewers while protecting CM Yogi and PM Modi throughout.
\end{itemize}

\subsection*{Excerpt Record 4 (Event 2 --- Spinning Pradhan's Resignation as "The Government Outplayed Them")}
\begin{itemize}[leftmargin=1.5em, itemsep=3pt, topsep=4pt]
  \item \textbf{Outlet:} Republic TV | \textbf{Date:} \texttt{2026-07-25}
  \item \textbf{Video title:} \textit{Rahul The Biggest Loser: Arnab Unravels The Political Agenda Behind CJP's Violent Protest}
  \item \textbf{Timestamp(s):} \texttt{[00:00--00:51]} \& \texttt{[01:37--02:56]} | \textbf{Analyst:} Baani Solanki | \textbf{Language:} English (Roman script)
  \item \textbf{Exact transcript excerpt:}
    \begin{quote}\small\ttfamily [00:00] The biggest loser is Rahul Gandhi. Let me tell you why, ladies and gentlemen. Dharmendra Pradhan has resigned. That demand has been met. The protest has [00:07] been called off by the CJP [ASR: BJP] leadership... [00:23] Not Dipke, not Deep State Dipke, [00:33] not Sonam Wangchuk, very questionable man, but this man, 57-year-old Rahul Gandhi... because he wants streetocracy...\end{quote}
    \begin{quote}\small\ttfamily [01:37] And now, after there was absolute anarchy on the streets, our police, our RAF, our security personnel were [01:47] vehemently attacked by rabid, bloodthirsty mobs, comes a resignation, the very thing that the students [01:56] demanded. But Rahul Gandhi is not happy with that... [02:06] He has been outplayed by the government. Rahul Gandhi, you've been outplayed... [02:38] They didn't want the FCRA bill to be brought in. They wanted the conversion of Hindus, but the government outplayed them.\end{quote}
  \item \textbf{Feature(s):} $\boxtimes$ Naming $\boxtimes$ Presupposition $\boxtimes$ Agency $\boxtimes$ Antagonisation $\boxtimes$ Intensification / softening
  \item \textbf{Who is speaking? Is this the outlet's wording or a quotation?} Anchor Arnab Goswami's direct-to-camera monologue.
  \item \textbf{What does the wording imply, foreground, soften, or leave unstated?} When a Cabinet Minister is forced to resign after a 35-day street protest, normal political grammar frames the government as yielding. Arnab flips defeat into triumph via three linguistic moves: (1) agentless nominalization \texttt{"comes a resignation"} (\texttt{[01:47]}), (2) shifting the grammatical subject of defeat onto \texttt{"Rahul Gandhi"} (\texttt{"The biggest loser"}, 5x \texttt{"outplayed"}), and (3) inventing sensational communal stakes (\texttt{"They wanted the conversion of Hindus"}, \texttt{[02:38]}).
  \item \textbf{What context or evidence supports or complicates this interpretation?} Supported only by the fact that Pradhan's resignation successfully ended the Jantar Mantar sit-in before the rest of the Parliament session (\texttt{[00:07--00:23]}). Complicated by the glaring contradiction of calling protesters \texttt{"Deep State"} \texttt{"rabid, bloodthirsty mobs"} in the very same sentence that admits \texttt{"comes a resignation, the very thing that the students demanded"} (\texttt{[01:47--01:56]})!
  \item \textbf{How does another outlet describe the same event or actor?} \textbf{Sarthak Goswami} (\texttt{EM1GeZufO6E}, \texttt{[00:08--00:58]}) frames the exact same resignation the next day as \texttt{"CJP won... aakhiri mein eego todna hi pada"} [finally the government's ego had to break], and directly mocks \texttt{"doosra Goswami"} at \texttt{[22:22--23:04]}.
  \item \textbf{My interpretation, with reasons:} This passage is a textbook demonstration of \textbf{spin via Agency \& Naming}: turning a historic student victory over the government into \texttt{"Rahul The Biggest Loser"}.
\end{itemize}

\subsection*{Excerpt Record 5 (Event 2 --- Erasing Students \& Linking Jantar Mantar to Anantnag Terrorists)}
\begin{itemize}[leftmargin=1.5em, itemsep=3pt, topsep=4pt]
  \item \textbf{Outlet:} Republic TV | \textbf{Date:} \texttt{2026-07-25}
  \item \textbf{Video title:} \textit{Rahul The Biggest Loser: Arnab Unravels The Political Agenda Behind CJP's Violent Protest}
  \item \textbf{Timestamp(s):} \texttt{[03:36--05:26]} | \textbf{Analyst:} Baani Solanki | \textbf{Language:} English (Roman script)
  \item \textbf{Exact transcript excerpt:}
    \begin{quote}\small\ttfamily [03:36] Do not forget that there was an attempt to lynch an Rapid Action Force personnel... [03:44] Nationalists do not forget that RAF officer Sonia Sherawat [03:54] she was harassed she was targeted she was beaten by the Lynch mob of Abhijit Dipke... [04:19] when this violence was being carried out by the goons of Abhijit Dipke our brave soldiers were fighting terrorists IN ANANTNAG... [04:46] while you were carrying out this drunken violence a soldier risked and lost his life... [04:58] this sexual perversion shown at Jantar Mantar in the name of students these were not students you were fooled... [05:08] they say sansad ukhad denge Nepal [05:16] model they wanted to bring in India... [05:26] they were not students who did this the students had evaporated.\end{quote}
  \item \textbf{Feature(s):} $\boxtimes$ Naming $\boxtimes$ Presupposition $\boxtimes$ Agency $\boxtimes$ Antagonisation $\boxtimes$ Intensification / softening
  \item \textbf{Who is speaking? Is this the outlet's wording or a quotation?} Arnab Goswami's direct monologue.
  \item \textbf{What does the wording imply, foreground, soften, or leave unstated?} Deploys \textbf{temporal juxtaposition (\texttt{when... while...})} to morally equate domestic protesters in Delhi with Kashmir terrorists killing an Indian soldier in Anantnag (\texttt{[04:19--04:46]}), while erasing genuine NEET examinees via \texttt{"the students had evaporated"} (\texttt{[05:26]}).
  \item \textbf{What context or evidence supports or complicates this interpretation?} Injuries to police/RAF personnel (and peripheral hooliganism near CP noted even by Sarthak Goswami \texttt{EM1GeZufO6E} \texttt{[07:12--07:21]}) had a factual kernel, which Arnab generalizes to the \textit{entire} protest while omitting 100\% of police pellet-gun/lathi injuries to students and doctors (\texttt{Mohak Mangal} \texttt{v-RikSqb0Vw} \texttt{[01:21--02:59]}; \texttt{The Wire} \texttt{pYaTVvHeGoM} \texttt{[10:07--10:38]}).
  \item \textbf{How does another outlet describe the same event or actor?} \textbf{Mohak Mangal} (\texttt{v-RikSqb0Vw}, \texttt{[01:03]}, \texttt{[06:06]}, \texttt{[10:56]}) samples this exact Republic TV rhetoric (\texttt{"brutal violence carried out by the goons of the CJP"}, \texttt{"these are not students"}) and refutes it solely by letting real NEET repeaters, mothers, and teachers speak on camera.
  \item \textbf{My interpretation, with reasons:} Invoking martyred soldiers in Anantnag (\texttt{[04:19--04:46]}) and Operation Sindoor (\texttt{[01:24]}) during an exam-leak story is classic \textbf{national-security displacement}: making any sympathy for student protesters feel like treason against the armed forces.
\end{itemize}

\subsection*{Excerpt Record 6 (Event 2 --- Closing Imperative Mobilizing "Nationalists of Bharat")}
\begin{itemize}[leftmargin=1.5em, itemsep=3pt, topsep=4pt]
  \item \textbf{Outlet:} Republic TV | \textbf{Date:} \texttt{2026-07-25}
  \item \textbf{Video title:} \textit{Rahul The Biggest Loser: Arnab Unravels The Political Agenda Behind CJP's Violent Protest}
  \item \textbf{Timestamp(s):} \texttt{[05:34--06:38]} | \textbf{Analyst:} Baani Solanki | \textbf{Language:} English (Roman script)
  \item \textbf{Exact transcript excerpt:}
    \begin{quote}\small\ttfamily [05:34] Dharmendra Pradhan has resigned when will Abhijit Dipke be held to account now [05:43] and if nationalist nationalist if you forgive and forget what they have done to our country this time, [05:52] then we will have to face it a second time. If you forget these pictures, then they will do it a third time. So, [06:01] learn your lesson and don't allow this CJP [ASR: BJP] to become a model... [06:10] Come together. India cannot forgive. India cannot forget. India will not forget or forgive all the [06:21] goons [ASR: goods] of Abhijit Dipke. Nationalists, mobilize yourself. Become more visible. [06:30] Become more vocal. Become more assertive. This nation is ours.\end{quote}
  \item \textbf{Feature(s):} $\boxtimes$ Naming $\boxtimes$ Presupposition $\boxtimes$ Agency $\boxtimes$ Antagonisation $\boxtimes$ Intensification / softening
  \item \textbf{Who is speaking? Is this the outlet's wording or a quotation?} Arnab Goswami's closing peroration addressed directly to \texttt{"Nationalists"}.
  \item \textbf{What does the wording imply, foreground, soften, or leave unstated?} Abandons journalistic reporting altogether for \textbf{political mobilization}: equates his ideological in-group (\texttt{"Nationalists"}) with \texttt{"India"} itself (\texttt{"India will not forget or forgive"}, \texttt{[06:10]}), and closes with four staccato imperatives (\texttt{"mobilize yourself. Become more visible. Become more vocal. Become more assertive"}) culminating in the exclusionary possessive \texttt{"This nation is ours"} (\texttt{[06:38]})---implying dissenting youth lie outside \texttt{"our"} nation!
  \item \textbf{What context or evidence supports or complicates this interpretation?} Reveals deep establishment anxiety that a decentralized, youth-led movement (\texttt{"don't allow this CJP to become a model"}, \texttt{[06:01]}) succeeded in toppling a Cabinet Minister, prompting a call for counter-street mobilization.
  \item \textbf{How does another outlet describe the same event or actor?} Contrast \texttt{"This nation is ours"} (Republic TV \texttt{[06:38]}, exclusionary partisan ownership) with \textbf{Mohak Mangal}'s title \texttt{"We the people of India..."} (\texttt{v-RikSqb0Vw}, inclusive constitutional citizenship) and \textbf{The Wire}'s Naseeruddin Shah (\texttt{pYaTVvHeGoM}, \texttt{[25:33--25:42]}: \texttt{"That house of parliament is ours. It belongs to the kids"}).
  \item \textbf{My interpretation, with reasons:} Comparing who \texttt{"ours"} refers to in Republic TV (\texttt{"Nationalists... This nation is ours"}) versus The Wire/Mohak Mangal (\texttt{"We the people... It belongs to the kids"}) captures the entire ideological fault line of contemporary Indian media discourse.
\end{itemize}
\vspace{4pt}\noindent\rule{\textwidth}{0.4pt}\vspace{4pt}

\section*{Common form: End-of-event summary}

\subsection*{End-of-Event Summary --- Event 1 (Ram Mandir Donation Controversy)}
\begin{itemize}[leftmargin=1.5em, itemsep=3pt, topsep=4pt]
  \item \textbf{The main way this outlet framed the event was:} As an isolated \texttt{Rs.\,80+ lakh} cash-counting theft by 8 \texttt{"low-level functionaries"} (\texttt{[03:33]}, \texttt{[07:14--07:26]}) swiftly punished under CM Yogi Adityanath's \texttt{"zero tolerance"} governance (\texttt{[01:12]}, \texttt{[03:22]}), which anti-Hindu opposition leaders (\texttt{Kejriwal}, \texttt{Akhilesh}, \texttt{Rahul}) hypocritically politicized---justifying putting the opposition \texttt{"in the dock"} (\texttt{[05:14--05:24]}) and condemning AAP's \texttt{"Taliban"} rhetoric against 79-year-old VHP leader Champat Rai (\texttt{[17:20--21:28]}).
  \item \textbf{The strongest linguistic evidence was (quote + timestamp):} \texttt{"there is no religious angle, there is no political angle, there is no government angle in this controversy... since they have politicized the matter I suggest the best response is to put them in the dock and ask them what their record is on Hindutva."} (\texttt{[04:08]} \& \texttt{[05:14--05:24]}).
  \item \textbf{The shared feature(s) most clearly illustrated were:} \textbf{Naming} (\texttt{low-level functionaries, some cashiers} vs \texttt{79-year-old prominent VHP leader}), \textbf{Agency} (agentless \texttt{"there has been theft"} vs active credit \texttt{"immediate action by the Yogi government"} at \texttt{[00:54--01:12]}), and \textbf{Intensification / Antagonisation} (\texttt{GENOCIDE OF KAR SEVAKS} at \texttt{[04:56]} and \texttt{"PART OF THE TALIBAN SHOOTING SQUAD"} at \texttt{[33:37]}).
  \item \textbf{An alternative interpretation of the same evidence is:} Articulated briefly by Arnab himself at \texttt{[39:47--40:16]} and by opposition panelists (\texttt{[14:03--15:23]}, \texttt{[27:34]}): that multi-year organized theft across 70+ CCTV instances and unreceipted silver gifts could not occur without senior Trust negligence or patronage, making quiet resignations without FIR naming an \texttt{"honorable exit"} rather than true \texttt{"zero tolerance"}.
  \item \textbf{What I cannot conclude from this transcript alone:} What the unpublished SIT report concluded regarding senior trustees' oversight, whether Raju Manwani's \texttt{200 kg} silver bricks or 2021 land-deal allegations were investigated, or whether custodial interrogation of the 8 clerks (\texttt{[50:13--50:22]}) implicated higher-ups.
\end{itemize}

\subsection*{End-of-Event Summary --- Event 2 (Students' Protest Regarding NEET Paper Leak)}
\begin{itemize}[leftmargin=1.5em, itemsep=3pt, topsep=4pt]
  \item \textbf{The main way this outlet framed the event was:} By \textbf{completely erasing NEET paper leaks and genuine students} (\texttt{"the students had evaporated"}, \texttt{[05:26]}), portraying the Jantar Mantar movement as a \texttt{"Deep State"} / Rahul Gandhi \texttt{"streetocracy"} conspiracy (\texttt{[00:23--00:51]}) of \texttt{"rabid, bloodthirsty mobs"} attacking RAF/police while soldiers died fighting terrorists in Anantnag (\texttt{[01:47]}, \texttt{[04:19--04:46]}), spinning Dharmendra Pradhan's resignation as a masterstroke that \texttt{"outplayed"} the opposition (\texttt{[02:06--02:56]}), and exhorting \texttt{"Nationalists"} never to forgive or forget (\texttt{[05:43--06:38]}).
  \item \textbf{The strongest linguistic evidence was (quote + timestamp):} \texttt{"our police, our RAF, our security personnel were vehemently attacked by rabid, bloodthirsty mobs, comes a resignation... these were not students, you were fooled... Nationalists, mobilize yourself... This nation is ours."} (\texttt{[01:37--01:47]}, \texttt{[04:58]} \& \texttt{[06:21--06:38]}).
  \item \textbf{The shared feature(s) most clearly illustrated were:} \textbf{Naming} (\texttt{Deep State Dipke}, \texttt{rabid, bloodthirsty mobs}, \texttt{Lynch mob}, \texttt{streetocracy}), \textbf{Agency} (deleting 100\% of police violence and hiding government capitulation inside agentless \texttt{"comes a resignation"} + 5x \texttt{"the government outplayed them"}), and \textbf{Antagonisation} (pitting \texttt{"Nationalists"} \& \texttt{"brave soldiers in Anantnag"} against protesters, closing with \texttt{"This nation is ours"}).
  \item \textbf{An alternative interpretation of the same evidence is:} That Dharmendra Pradhan resigned and the government dropped cases against Dipke and Saurabh Das (\texttt{[00:00--00:15]}) precisely because the movement represented massive, legitimate student anger over exam leaks---not a foreign \texttt{"Deep State"} mob (which MEA itself denied, \texttt{EM1GeZufO6E} \texttt{[00:40]}).
  \item \textbf{What I cannot conclude from this transcript alone:} Anything at all about the NEET examination crisis (never mentioned once!), student casualties from lathis/tear-gas/pellet guns on 20 July, or the terms of CJP's settlement beyond Pradhan's resignation and non-prosecution assurance (\texttt{[00:00--00:15]}).
\end{itemize}
\textit{Research reminder: interpret wording in context, distinguish allegations from established facts, and support claims with exact excerpts. Mixed or contradictory evidence is worth recording.}


\clearpage

\begin{center}
    {\large\bfseries LANGUAGE \& SOCIETY PROJECT}\\[3pt]
    {\LARGE\bfseries Transcript Analysis Worksheet}\\[3pt]
    {\normalsize Qualitative analysis of political news discourse}
\end{center}
\vspace{2pt}
\noindent\rule{\textwidth}{0.4pt}
\vspace{4pt}


\section*{Creator 6: The Wire (Independent Left-Liberal Long-Form Digital Journalism)}
\begin{itemize}[leftmargin=1.5em, itemsep=3pt, topsep=4pt]
  \item \textbf{Analyst:} Baani Solanki (\texttt{2025114016})
  \item \textbf{Channel / Outlet:} The Wire (\texttt{@TheWire}, Foundation for Independent Journalism)
  \item \textbf{Ideological / Editorial Positioning:} Independent, critical left-liberal English digital journalism; uses long-form one-on-one elite interviews (with senior historians/authors and cultural figures) to situate breaking controversies within structural critiques of Hindutva/Sangh Parivar institutional history, PMO centralization, and democratic backsliding.
  \item \textbf{Transcript Files Referenced:}
  \begin{itemize}[leftmargin=1.4em, itemsep=2pt, topsep=2pt]
    \item Event 1: \texttt{transcripts/event1\_ram\_mandir/06\_the\_wire\_7rUL8fGgRJA.txt}
    \item Event 2: \texttt{transcripts/event2\_neet\_protest/06\_the\_wire\_pYaTVvHeGoM.txt}
  \end{itemize}
\end{itemize}
\vspace{4pt}\noindent\rule{\textwidth}{0.4pt}\vspace{4pt}

\section*{Guidelines: Five features to look for in every event}
\begin{longtable}{@{}>{\raggedright\arraybackslash}p{0.17\textwidth} >{\raggedright\arraybackslash}p{0.20\textwidth} >{\raggedright\arraybackslash}p{0.24\textwidth} >{\raggedright\arraybackslash}p{0.32\textwidth}@{}}
\toprule
\textbf{Feature} & \textbf{What to look for} & \textbf{Illustrative example} & \textbf{Questions to answer} \\
\midrule
\endhead
\textbf{1. Naming} & What are people, groups, actions, or events called? & "Students", "protesters", "agitators"; "irregularity", "scam", "controversy". & What exact label is used? Who or what does it refer to? Who uses it? What alternative label could have been used? \\
\addlinespace[4pt]
\textbf{2. Presupposition} & What is treated as already true or taken for granted? & "Why are students creating chaos again?" may presuppose earlier disruption. & What does the wording assume? Is that assumption supported, disputed, or left unexamined? Does it appear in a question or words such as "again", "still", or "finally"? \\
\addlinespace[4pt]
\textbf{3. Agency} & Who is shown doing something? Who is acted upon? Who disappears? & "Police injured protesters" vs. "Protesters were injured." & Who is the actor? Who experiences the action? Is responsibility named, backgrounded, or omitted? \\
\addlinespace[4pt]
\textbf{4. Antagonisation} & Who is positioned against whom? & "Ordinary citizens" vs. "political troublemakers". & Which groups are contrasted? Is one presented as more legitimate, moral, patriotic, dangerous, or disloyal? Is an "us vs. them" division implied? \\
\addlinespace[4pt]
\textbf{5. Intensification / softening} & What is exaggerated, minimised, softened, or euphemised? & "Violent mob" vs. "peaceful demonstrators"; "incident" instead of a direct description. & Which words heighten emotion or seriousness? Which minimise it? Are metaphors, dramatic adjectives, hedges, or euphemisms used? Does the wording match the evidence? \\
\bottomrule
\end{longtable}
\textit{Reminder: A label or implication is not automatically false or manipulative. Examine the surrounding context and evidence.}

\vspace{4pt}\noindent\rule{\textwidth}{0.4pt}\vspace{4pt}

\section*{Event1. Ram Mandir donation controversy}
\textbf{Focus:} How is the controversy framed---as financial accountability, religious sentiment, political motivation, or something else? Be precise about the specific allegation or episode covered.

\begin{itemize}[leftmargin=1.5em, itemsep=3pt, topsep=4pt]
  \item \textbf{Outlet:} The Wire (Special Discussion hosted by M.K. Venu / Venugopal with senior journalist \& Ayodhya historian Nilanjan Mukhopadhyay)
  \item \textbf{Date:} \texttt{2026-07-01}
  \item \textbf{Video title:} \textit{Ram Temple Donations Loot: Biggest Credibility Test For Modi-RSS} (\url{https://youtu.be/7rUL8fGgRJA?si=9HV3jAYEipsTwvAI})
  \item \textbf{Timestamp(s):} Relevant Event 1 segment \texttt{[00:00--39:06]} (\texttt{[00:00--02:48]} cold-open highlights reel + \texttt{[03:00--39:06]} full conversation; \texttt{100\%} focused on Event 1; total duration \texttt{39:18}).
  \item \textbf{Analyst:} Baani Solanki (\texttt{2025114016})
  \item \textbf{Language:} English (Native English long-form discussion, Roman script; non-AI-dubbed)
\end{itemize}

\subsection*{Event-specific questions (Event 1)}
\begin{enumerate}[leftmargin=1.6em, itemsep=4pt, topsep=4pt]
  \item \textbf{What terms are used: "scam", "alleged irregularity", "corruption", "political controversy", or something else?}
  \begin{itemize}[leftmargin=1.4em, itemsep=2pt, topsep=2pt]
    \item Uses high-level structural, legal, and moral terminology: \texttt{"Ram Temple Donations Loot"} (\texttt{[Title]}), \texttt{"the Ram Temple scandal"} (\texttt{[03:10]}), \texttt{"massive theft or decoity [dacoity], as Nripendra Misra... has described it"} (\texttt{[00:09--00:17]}, \texttt{[04:23--04:32]}), \texttt{"a land scam of this magnitude involving potentially massive money laundering via land"} (\texttt{[00:17--00:26]}, \texttt{[04:32--04:40]}), \texttt{"hawala"} (\texttt{[26:42]}), \texttt{"moral degeneration"} (\texttt{[00:44]}, \texttt{[09:08]}), \texttt{"moral and ethical morass"} (\texttt{[28:12--28:21]}), and \texttt{"a platform of corruption, theft and lies"} (\texttt{[37:48]}).
  \end{itemize}
  \item \textbf{How are the temple, trust, donors, and people making allegations described?}
  \begin{itemize}[leftmargin=1.4em, itemsep=2pt, topsep=2pt]
    \item \textbf{Temple \& Movement:} Analysed historically not as an eternal pilgrimage centre, but as \texttt{"the Sangha Parivar's, the BJP's... ideological mission, project"} since \texttt{1949}/\texttt{1983} (\texttt{[03:37--03:48]}, \texttt{[05:29--05:38]}) and \texttt{"Ayodhya the Hindu Vatican... actually the Sangh Parivar's creation"} (\texttt{[35:04--35:13]}).
    \item \textbf{Trust \& Leadership:} Described as a government-nominated \texttt{"corporate bureaucratic structure"} (\texttt{[11:00--11:39]}, \texttt{[23:19]}) staffed by \texttt{"political loyalists"} (\texttt{[25:39--25:48]}) from an RSS culture of \texttt{"shadow organizations"} (\texttt{[22:24]}) and unaccounted \textit{Guru Dakshina} (\texttt{[17:35--18:40]}). Champat Rai is characterized via his nickname \texttt{"Patwari"} (adept at land/revenue records rather than transparent corporate governance, \texttt{[24:18--24:36]}), while Construction Chair Nripendra Misra is identified as \texttt{"very much a part of the core team of Mr. Modi"} who \texttt{"surely knew"} and kept \texttt{"superiors [PMO] in the loop"} (\texttt{[11:48--14:55]}).
    \item \textbf{Donors:} Portrayed with sociological nuance---on one hand sincere devotees from the 1989 \textit{Ram Shila Yatras} onward who \texttt{"put their body and soul into this movement"} only to find \texttt{"the darshan has been hollowed out, that it has been stolen"} (\texttt{[36:28--37:08]}), and on the other hand rational-losing \texttt{"bhakts"} viewing Modi as \texttt{"not just a biological being"} who might \texttt{"overlook anything"} (\texttt{[02:18--02:27]}, \texttt{[35:43--36:20]}).
    \item \textbf{People Making Allegations:} Highlights \textbf{internal Hindu/movement critics}---specifically dissenting orders like \texttt{Nirmohi Akhada} (\texttt{[ASR: Nirmala Akada]}, \texttt{[29:11--29:38]}) calling Sangh Parivar trustees \texttt{"all thieves"} after being excluded from the Trust, alongside historical references to honest 1990 Income Tax Deputy Commissioner \texttt{Vishwa Bandhu Gupta} (\texttt{[ASR: Vishva Hindu Gupta]}, \texttt{[15:33--17:15]}) and Karnataka minister \texttt{Priyank Kharge} (\texttt{[18:59--19:09]}).
  \end{itemize}
  \item \textbf{Is the issue framed as financial accountability, an attack on Hindu faith, a political dispute, or a matter of public trust?}
  \begin{itemize}[leftmargin=1.4em, itemsep=2pt, topsep=2pt]
    \item Framed as an \textbf{existential ideological and moral credibility crisis (\texttt{"Biggest Credibility Test For Modi-RSS"}, \texttt{"biggest challenge for the RSS after the assassination of Mahatma Gandhi"}, \texttt{[08:23--08:31]})} rooted in the Sangh Parivar's 40-year institutional refusal to maintain transparent, audited financial accounts (\texttt{[17:24--22:32]}).
  \end{itemize}
  \item \textbf{Are questions about donations presented as legitimate scrutiny or as an attempt to malign the temple or its supporters?}
  \begin{itemize}[leftmargin=1.4em, itemsep=2pt, topsep=2pt]
    \item \textbf{Unequivocally legitimate structural and forensic scrutiny}---calling for Enforcement Directorate (\texttt{ED}) investigation into land-sale \textit{hawala} trails (\texttt{[26:42--27:18]}) and auditing of unregistered Sangh affiliates (\texttt{[17:24--21:30]}).
  \end{itemize}
  \item \textbf{Are allegations attributed to specific people and supported with evidence, or presented as established facts?}
  \begin{itemize}[leftmargin=1.4em, itemsep=2pt, topsep=2pt]
    \item \textbf{Carefully hedged on financial destinations, deductive on PMO awareness, and richly documented on institutional history:} Host M.K. Venu attributes \texttt{"massive theft or dacoity"} directly to Nripendra Misra (\texttt{[04:23]}), hedges money-laundering beneficiaries with epistemic modals (\texttt{"potentially massive money laundering... We don't know who are the beneficiaries... They could be political coffers, they could be Sangh Parivar associates"}, \texttt{[04:40--04:59]}), and explicitly notes \texttt{"We do not have any empirical data as yet"} on voter outrage (\texttt{[37:08--37:19]}). Meanwhile, Nilanjan Mukhopadhyay uses \textbf{deductive bureaucratic reasoning} (\texttt{"I find it very hard to believe... on that I have absolutely no doubt"}, \texttt{[13:04--14:55]}) to infer PMO knowledge via Nripendra Misra, backed by deep archival knowledge of \texttt{1983--1990} VHP/\texttt{Mukti Yagna Samiti}/\texttt{Ram Janmabhoomi Nyas} structures (\texttt{[15:25--22:50]}).
  \end{itemize}
  \item \textbf{How are denials and counterarguments represented: explained, dismissed, challenged, or omitted?}
  \begin{itemize}[leftmargin=1.4em, itemsep=2pt, topsep=2pt]
    \item \textbf{Analysed politically rather than hosted live} (as there is no BJP/RSS guest in this two-journalist format): UP CM Yogi Adityanath's public assurance is dismissed as \texttt{"having to put up a brave front"} (\texttt{[04:14]}, \texttt{[29:00]}) amid \texttt{"a major fissure between Yogi Adityanath and the central leadership"} (\texttt{[30:04--30:24]}), and the RSS's traditional defense---playing state-persecuted victimhood as during the Emergency or 1992 ban---is dismantled because \texttt{"Today, they are in power. So, they can't play victims"} (\texttt{[10:14--10:49]}).
  \end{itemize}
  \item \textbf{Does the presenter distinguish criticism of the trust or its financial practices from criticism of the temple or religion itself?}
  \begin{itemize}[leftmargin=1.4em, itemsep=2pt, topsep=2pt]
    \item \textbf{Distinguishes classical Hinduism from Sangh Parivar political Hindutva:} Mukhopadhyay emphasizes that established Hindu spiritual heads (\texttt{Shankaracharyas} / traditional \textit{acharyas} and \textit{akhadas} in Varanasi, Mathura, Haridwar) were deliberately \texttt{"kept out"} by the Sangh Parivar when constructing \texttt{"Ayodhya the Hindu Vatican"} (\texttt{[34:07--35:22]}), framing the scandal as a crisis of \textit{Sangh Parivar real-estate/corporate Hindutva} (\texttt{[26:14--26:32]}), not traditional Hindu spirituality.
  \end{itemize}
  \item \textbf{Who is portrayed as protecting donors' interests, and who as exploiting religious sentiment?}
  \begin{itemize}[leftmargin=1.4em, itemsep=2pt, topsep=2pt]
    \item \textbf{Exploiting religious sentiment:} The BJP-RSS-VHP \textit{"Sangh Parivar"}, which mobilized grassroots \textit{Ram Shila} devotees (\texttt{[36:28--37:08]}), sidelined traditional Hindu \textit{akhadas} (\texttt{[29:20--29:38]}, \texttt{[34:07--35:22]}), and deliberately avoided \texttt{"leak-proof"} accounting structures (\texttt{[24:44--25:01]}) while getting \texttt{"addicted to power"} and plush Delhi headquarters (\texttt{[31:22--32:27]}).
    \item \textbf{Protecting public accountability:} Transferred 1990 tax officer Vishwa Bandhu Gupta (\texttt{[15:33--17:15]}) and independent journalists/opposition actors if they can \texttt{"think creatively"} ahead of UP elections (\texttt{[37:28--38:24]}).
  \end{itemize}
\end{enumerate}

\subsection*{Apply the five shared features (Event 1)}
\begin{longtable}{@{}>{\raggedright\arraybackslash}p{0.19\textwidth} >{\raggedright\arraybackslash}p{0.38\textwidth} >{\raggedright\arraybackslash}p{0.37\textwidth}@{}}
\toprule
\textbf{Feature} & \textbf{Exact wording + timestamp} & \textbf{What it does in context} \\
\midrule
\endhead
\textbf{1. Naming} & \texttt{"shadow banks, you know, shadow organizations"} \texttt{[22:24]}; \texttt{"corporate bureaucratic structure"} staffed by \texttt{"political loyalists"} \texttt{[23:19]}, \texttt{[25:48]}; \texttt{"Ayodhya the Hindu Vatican"} \texttt{[35:13]} & Replaces devotional labels with critical political-economy and ecclesiastical terms (\texttt{shadow organizations}, \texttt{corporate bureaucratic structure}, \texttt{Hindu Vatican}), demystifying the Trust as a centralized Sangh power apparatus. \\
\addlinespace[4pt]
\textbf{2. Presupposition} & \texttt{"within 2 years of that, the unraveling of the Ram temple project in moral terms, in ethical terms"} \texttt{[06:15--06:26]} \newline  \texttt{"Sangha Parivar values got obviously got completely contaminated by the real estate game"} \texttt{[26:14]} & Nominalization (\texttt{the unraveling of the Ram temple project}) and adverb \texttt{"obviously"} presuppose that the temple project has already morally collapsed and been corrupted by real-estate profiteering. \\
\addlinespace[4pt]
\textbf{3. Agency} & \texttt{"Setting up structures was not done. It was probably not done because... they wanted that it should not be a leak-proof system. So you keep a certain amount of ambiguity, a certain amount of fluidity."} \texttt{[24:36--25:01]} & Converts what other outlets call passive "oversight" or "informality" (\texttt{anaupcharik} on Aaj Tak) into \textbf{deliberate, strategic agency} (\texttt{they wanted that it should not be a leak-proof system}). \\
\addlinespace[4pt]
\textbf{4. Antagonisation} & \texttt{"RSS chief described it as the real freedom India has attained as opposed to what was attained in 1947... and within two years... the entire edifice was built on a platform of corruption, theft and lies."} \texttt{[05:57--06:15]}, \texttt{[37:48]} & Contrasts RSS Chief Mohan Bhagwat's grand civilizational claim (\texttt{real freedom} vs \texttt{1947}) against sordid financial crime (\texttt{hawala}, \texttt{corruption, theft and lies}), framing the Sangh Parivar's moral claim over India as hollow. \\
\addlinespace[4pt]
\textbf{5. Intensification / softening} & \textbf{Intensification:} \texttt{"surrounded by 100 questions today, maybe 1,000 questions"} \texttt{[00:00]}, \texttt{[04:05]}; \texttt{"biggest challenge for the RSS after the assassination of Mahatma Gandhi"} \texttt{[08:23--08:31]} \newline  \textbf{Epistemic Softening:} \texttt{"potentially massive money laundering... We don't know who are the beneficiaries... They could be..."} \texttt{[04:40--04:59]}; \texttt{"We do not have any empirical data as yet"} \texttt{[37:19]} & Pairs sweeping historical intensification (comparing the crisis to 1948 post-Gandhi assassination) with careful journalistic hedging (\texttt{potentially}, \texttt{could be}, \texttt{no empirical data}) wherever hard forensic proof is still missing. \\
\bottomrule
\end{longtable}

\subsection*{Short summary (Event 1)}
\begin{itemize}[leftmargin=1.5em, itemsep=3pt, topsep=4pt]
  \item \textbf{What interpretation of the event does this video encourage?} That the Ayodhya donation \texttt{"dacoity"} and \texttt{5x--10x} land-price inflation (\texttt{[04:23]}, \texttt{[26:42]}) are not rogue clerk anomalies, but the direct institutional consequence of the Sangh Parivar's decades-long ethos of opaque, unaudited \texttt{"shadow organizations"} (\texttt{[17:24--22:32]}, \texttt{[24:44--25:01]}) colliding with massive post-2020 state patronage under PMO-linked appointees like Nripendra Misra and \texttt{"Patwari"} Champat Rai (\texttt{[11:48--14:55]}, \texttt{[24:18]}).
  \item \textbf{Whose perspective is foregrounded, challenged, or missing?} \textit{Foregrounded:} Secular-historical political analysis (M.K. Venu \& Nilanjan Mukhopadhyay), sidelined Hindu monastic orders (\texttt{Nirmohi Akhada}, traditional spiritual heads at \texttt{[29:11--29:38]}, \texttt{[34:07--35:22]}), and ordinary 1989 \textit{Ram Shila} devotees whose darshan was \texttt{"hollowed out"} (\texttt{[36:28--37:08]}). \textit{Challenged:} Sangh Parivar moral supremacy (\texttt{[05:57--06:35]}), RSS financial opacity (\texttt{[17:24--19:35]}), and PMO/Yogi plausible deniability (\texttt{[12:08--14:55]}, \texttt{[28:31--30:24]}). \textit{Missing:} Any Trust/RSS/BJP representative or ground reportage from Ayodhya's counting room.
  \item \textbf{What evidence supports your interpretation? What complicates it?} \textit{Supports:} Documented historical precedents---\texttt{1990} IT notice by Vishwa Bandhu Gupta over \texttt{Rs.\,700 crore} VHP collections and his punitive transfer (\texttt{[15:33--17:15]}), dual registered/unregistered 1984 \texttt{Mukti Yagna Samiti} vs \texttt{Nyas} architecture (\texttt{[20:36--22:24]}), \texttt{2019} Supreme Court trust mandate (\texttt{[11:08--11:39]}), and Nripendra Misra's own word \texttt{"dacoity"} (\texttt{[04:23]}). \textit{Complicates:} Both speakers candidly admit they lack direct proof of money-laundering recipients (\texttt{[04:48]}), ED investigation trails (\texttt{[26:51--27:18]}), or empirical data on whether \texttt{"bhakts"} will actually punish the BJP electorally in UP (\texttt{[36:01--37:28]}).
\end{itemize}
\vspace{4pt}\noindent\rule{\textwidth}{0.4pt}\vspace{4pt}

\section*{Event2. Students' protest regarding NEET paper leak}
\textbf{Focus:} How does coverage construct the legitimacy of student protest and responsibility for the underlying problem?

\begin{itemize}[leftmargin=1.5em, itemsep=3pt, topsep=4pt]
  \item \textbf{Outlet:} The Wire (\textit{The Interview with Karan Thapar} --- Guest: veteran actor \& educator Naseeruddin Shah)
  \item \textbf{Date:} \texttt{2026-07-22} \textit{(two days after the 20 July "Black Monday" crackdown; three days before Dharmendra Pradhan's resignation)}
  \item \textbf{Video title:} \textit{I stand with the protesting students; this is not the India I love; I don't take Modi seriously} (\url{https://youtu.be/pYaTVvHeGoM?si=TPLfYLT3eB6UHh5o})
  \item \textbf{Timestamp(s):} Relevant Event 2 segments \texttt{[00:00--02:45]} (cold-open preview) and \texttt{[03:59--31:05]} (main interview; total duration \texttt{32:26}; viewer-donation appeals at \texttt{[02:46--03:58]} and \texttt{[31:06--32:26]} excluded per project guidelines; note \texttt{[04:25--06:15]} has sparse ASR during playback of Shah's 1m50s video clip).
  \item \textbf{Analyst:} Baani Solanki (\texttt{2025114016})
  \item \textbf{Language:} English with occasional Hindi/Urdu lexical items (\texttt{dandas}, \texttt{bachchas}, \texttt{Pradhan Sevak}, \texttt{Chowkidar}) (Native English audio, Roman script; non-AI-dubbed)
\end{itemize}

\subsection*{Event-specific questions (Event 2)}
\begin{enumerate}[leftmargin=1.6em, itemsep=4pt, topsep=4pt]
  \item \textbf{Are participants called "students", "candidates", "protesters", "agitators", "mobs", or "political pawns"?}
  \begin{itemize}[leftmargin=1.4em, itemsep=2pt, topsep=2pt]
    \item Warmly humanized and idealized throughout by both Karan Thapar and Naseeruddin Shah as \texttt{"protesting students"} (\texttt{[Title]}, \texttt{[04:09]}), \texttt{"people in their teens or their early 20s"} (\texttt{[06:33--06:41]}), \texttt{"our future"} / \texttt{"the children of our future"} (\texttt{[07:42]}, \texttt{[30:19]}), \texttt{"young children"} / \texttt{"these kids"} (\texttt{[08:38]}, \texttt{[09:46]}, \texttt{[20:15]}, \texttt{[25:33]}), and affectionately \texttt{"these bachas [bachchas]"} (\texttt{"These bachas know what they're doing... these bachas are much smarter than us"}, \texttt{[23:18--23:28]}).
  \end{itemize}
  \item \textbf{Are their stated demands explained, summarised, mocked, questioned, or left out?}
  \begin{itemize}[leftmargin=1.4em, itemsep=2pt, topsep=2pt]
    \item \textbf{Clearly explained and morally endorsed, with a nuanced tactical disagreement on Pradhan's resignation:} Karan Thapar summarises their demands---speaking to MPs about exam leaks and the education system (\texttt{[06:23]}), and Education Minister Dharmendra Pradhan's resignation accepting moral responsibility after multiple leaks and suicides (\texttt{[11:55--12:14]}, \texttt{[22:31--22:40]}). Interestingly, while Shah agrees Pradhan \texttt{"should have done so on his own initiative"} (\texttt{[12:14]}), he calls demanding it \texttt{"pointless"} because \texttt{"no way is he going to resign"} and urges pressing systemic demands more urgent than \texttt{"having this idiot to resign"} (\texttt{[12:23--12:50]}) \textit{(a prediction dramatically overturned 3 days later on 25 July!)}.
  \end{itemize}
  \item \textbf{Does the presenter focus on the original grievance, disruption, political affiliations, or public inconvenience?}
  \begin{itemize}[leftmargin=1.4em, itemsep=2pt, topsep=2pt]
    \item Focuses on \textbf{(a) the moral/democratic betrayal of state violence against peaceful youth (\texttt{[06:23--11:55]}), (b) broader syllabus degradation \& ministerial accountability (\texttt{[12:04--14:57]}, \texttt{[19:35--23:37]}), and (c) a scathing character indictment of PM Narendra Modi and authoritarian fear (\texttt{[00:00--02:45]}, \texttt{[15:07--19:01]}, \texttt{[23:47--30:40]})}---explicitly rejecting any disruption frame (\texttt{"It's not as if they were burning effigies or kicking things or destroying stuff"}, \texttt{[25:42--25:51]}).
  \end{itemize}
  \item \textbf{Who is represented as responsible: examination authorities, government, students, political parties, or others?}
  \begin{itemize}[leftmargin=1.4em, itemsep=2pt, topsep=2pt]
    \item \textbf{Prime Minister Narendra Modi} (\texttt{"arrogant autocrats"}, \texttt{[23:47]}; \texttt{"cowardly prime minister"}, \texttt{[26:02]}), \textbf{Education Minister Dharmendra Pradhan} (\texttt{"this idiot"}, Orwellian \texttt{"classic double speak"}, \texttt{[12:42--13:51]}), \textbf{Health Minister J.P. Nadda} (for theatrical hospital photo-ops instead of meeting hunger-striker Sonam Wangchuk, \texttt{[21:03--21:44]}), \textbf{Delhi Police / RAF / privately hired unbadged \texttt{"lumpens"}} (\texttt{[07:15--09:28]}, \texttt{[11:26--11:37]}), and \textbf{silent Indian elites/celebrities} stifling their conscience for state rewards (\texttt{"a dog with a bone in its mouth can't bark"}, \texttt{[29:31--30:00]}).
  \end{itemize}
  \item \textbf{Are questions framed around whether students have legitimate grievances, or why they are causing disruption?}
  \begin{itemize}[leftmargin=1.4em, itemsep=2pt, topsep=2pt]
    \item Karan Thapar's interview questions are \textbf{heavily presuppositional leading questions affirming student legitimacy and state guilt}: e.g., \texttt{"By its response, has the government embarrassed and betrayed India's cherished democracy?"} (\texttt{[06:41--06:49]}), \texttt{"Doesn't that suggest that our political value system has degraded?"} (\texttt{[22:40--22:48]}), and \texttt{"Do peaceful nonviolent Indian protesters protesting in a just cause have the right to walk to parliament...?"} (\texttt{[25:02--25:23]}).
  \end{itemize}
  \item \textbf{Are claims of violence, disorder, or political manipulation attributed to sources and substantiated?}
  \begin{itemize}[leftmargin=1.4em, itemsep=2pt, topsep=2pt]
    \item Karan Thapar anchors key factual claims to \textbf{named mainstream print sources and legal experts}: \textit{Hindustan Times} on hundreds injured (\texttt{[06:33]}), \textit{The Hindu} on pellet guns and shock batons injuring 25-year-old \texttt{Sheikh Irshad Mansuri} with pellets lodged near a neck blood vessel requiring multiple surgeries (\texttt{[10:07--10:38]}), and Senior Advocate \texttt{Kapil Sibal}'s press conference warning that unnamed FIRs could escalate to draconian \texttt{UAPA} / \texttt{BNS Section 195} charges as in the \texttt{CAA} protests (\texttt{[26:41--27:37]}). Conversely, Shah's claims that plainclothes attackers are \texttt{"uneducated... lumpens who've been recruited by being paid"} (\texttt{[11:26--11:37]}) or that PM Modi \texttt{"left by the back gate"} (\texttt{[26:02--26:13]}) rely on viral video impressions and hearsay (\texttt{"witnesses have said"}, \texttt{[21:23]}).
  \end{itemize}
  \item \textbf{How are police and government actions described compared with protesters' actions? Are similar actions evaluated differently?}
  \begin{itemize}[leftmargin=1.4em, itemsep=2pt, topsep=2pt]
    \item Police/state attackers are described as \texttt{"sadistically beating up these students"} (\texttt{[07:15]}), \texttt{"let loose an army of goons on our future"} (\texttt{[07:42]}), \texttt{"splitting young children's heads open, women being attacked and molested"} (\texttt{[08:38--08:47]}), and firing guns reminiscent of \texttt{"Jallianwala Bagh"} (\texttt{[10:48--11:08]}). When discussing \textbf{police casualties}, Shah openly mocks official injury figures (\texttt{"supposedly more policemen injured than students... astounding... something like 600 policemen and about 60 students"}, \texttt{[06:49--07:07]}) and justifies student pushback via a cornered-victim proverb: \texttt{"the students retaliating. It's not surprising. I mean, even a rat fights back when it's pushed against the wall. And these kids are justified"} (\texttt{[09:36--09:46]}).
  \end{itemize}
  \item \textbf{Are students allowed to explain their demands in their own words, or does the presenter mainly speak for them?}
  \begin{itemize}[leftmargin=1.4em, itemsep=2pt, topsep=2pt]
    \item \textbf{Elite proxy advocacy rather than direct student voices:} Because the format is a studio Zoom interview between veteran journalist Karan Thapar and 75-year-old actor Naseeruddin Shah, \textbf{no student speaks directly}; instead, two prominent public intellectuals speak \textit{in solidarity with} (\texttt{"I stand with them"}, \texttt{[20:25]}; \texttt{"We owe it to the children"}, \texttt{[30:19]}) and interpret the youth movement for an English-speaking civic audience.
  \end{itemize}
  \item \textbf{Does the outlet's framing differ between coverage of the NEET paper-leak issue before the protest and coverage once students begin protesting?}
  \begin{itemize}[leftmargin=1.4em, itemsep=2pt, topsep=2pt]
    \item Rather than delegitimizing students once they march to Parliament, The Wire treats the \textbf{20 July march and students' courageous return to Jantar Mantar in bandages on 21 July (\texttt{[19:09--19:18]})} as a historic moral watershed (\texttt{"bigger than anything that's happened before"}, \texttt{[30:09--30:19]}) exposing 12 years of autocratic governance (\texttt{[00:00--00:08]}, \texttt{[23:47--24:07]}).
  \end{itemize}
\end{enumerate}

\subsection*{Apply the five shared features (Event 2)}
\begin{longtable}{@{}>{\raggedright\arraybackslash}p{0.19\textwidth} >{\raggedright\arraybackslash}p{0.38\textwidth} >{\raggedright\arraybackslash}p{0.37\textwidth}@{}}
\toprule
\textbf{Feature} & \textbf{Exact wording + timestamp} & \textbf{What it does in context} \\
\midrule
\endhead
\textbf{1. Naming} & Attackers: \texttt{"an army of goons"} \texttt{[07:42]}, \texttt{"uneducated... lumpens"} \texttt{[11:26]} \newline  Leaders: \texttt{"this idiot"} \texttt{[12:42]}, \texttt{"arrogant autocrats"} \texttt{[23:47]}, \texttt{"rabblerouser"} \texttt{[24:43]}, \texttt{"cowardly prime minister"} \texttt{[26:02]} \newline  Students: \texttt{"our future"}, \texttt{"these kids"}, \texttt{"these bachas"} \texttt{[07:42]}, \texttt{[23:18]} & Inverts Republic TV's naming matrix: \texttt{"goons"} and \texttt{"lumpens"} are applied to unbadged state-aligned attackers, harsh moral-intellectual epithets target ministers/PM, and familial Hindi endearment (\texttt{bachas}) humanizes protesters. \\
\addlinespace[4pt]
\textbf{2. Presupposition} & Karan Thapar \texttt{[13:10]}: \texttt{"Accountable is the one thing the minister refuses to be"} \newline  \texttt{[22:40]}: \texttt{"Doesn't that suggest that our political value system has degraded?"} \newline  \texttt{[23:47]}: \texttt{"both he and his ministers are behaving like arrogant autocrats"} & Negative-interrogative (\texttt{Doesn't that suggest...}) and categorical declaratives embedded inside interview "questions" presuppose ministerial guilt and autocracy as shared consensus between interviewer and guest. \\
\addlinespace[4pt]
\textbf{3. Agency} & \texttt{"to let loose an army of goons on our future"} \texttt{[07:42]}; \texttt{"wrapping him [Wangchuk] up in bed sheets and dragging him away"} \texttt{[21:14--21:23]} vs \texttt{"the students retaliating... even a rat fights back when it's pushed against the wall"} \texttt{[09:36--09:46]} & Assigns sadistic, unprovoked initiatory agency to the state (\texttt{let loose}, \texttt{sadistically beating}, \texttt{dragging him away}), while framing student force strictly as cornered self-defense (\texttt{retaliating}, \texttt{pushed against the wall}). \\
\addlinespace[4pt]
\textbf{4. Antagonisation} & \texttt{"we are living in 1984... 2 minutes of hate has been transformed into 24 hours of hate"} \texttt{[14:02--14:11]} \newline  \texttt{"That house of parliament is ours. It belongs to the kids... This is not the India I was taught to love."} \texttt{[25:33--26:32]} & Pits an Orwellian, anti-intellectual, autocratic regime (\texttt{1984}, \texttt{24 hours of hate}, \texttt{15th century} syllabus at \texttt{[19:53]}) and complicit silent elites (\texttt{dog with a bone}, \texttt{[29:31]}) against courageous, smart youth reclaiming citizens' Parliament. \\
\addlinespace[4pt]
\textbf{5. Intensification / softening} & \textbf{Intensification:} \texttt{"Lie upon lie upon lie upon lie"} \texttt{[01:05]}, \texttt{[16:34]}; \texttt{"broke my heart and made my blood boil"} \texttt{[08:19]}; \texttt{"deaf, dumb and blind"} \texttt{[09:57]}; \texttt{"classic double speak, Newspeak"} \texttt{[13:21--13:33]} \newline  \textbf{Skeptical distancing:} \texttt{"supposedly more policemen injured"} \texttt{[06:49]} & Repetition (\texttt{lie upon lie}), visceral somatic idioms (\texttt{made my blood boil}), and literary allusions (\textit{1984}, \textit{Jallianwala Bagh}, Churchill/JFK/MLK speeches at \texttt{[24:35--24:43]}) intensify moral outrage, while \texttt{"supposedly"} casts doubt on police casualty reports. \\
\bottomrule
\end{longtable}

\subsection*{Short summary (Event 2)}
\begin{itemize}[leftmargin=1.5em, itemsep=3pt, topsep=4pt]
  \item \textbf{What interpretation of the event does this video encourage?} That the 20 July lathi/tear-gas/pellet-gun crackdown on teenagers and 20-somethings marching to Parliament---after 25 days of total ministerial indifference (\texttt{[20:35--20:54]}) and Sonam Wangchuk's forced removal (\texttt{[21:14--21:23]})---is an Orwellian betrayal of Indian democracy (\texttt{"living in 1984"}, \texttt{"not the India I was taught to love"}, \texttt{[14:02]}, \texttt{[26:32]}) driven by \texttt{"arrogant autocrats"} (\texttt{[23:47]}) who fear educated youth (\texttt{[18:23--20:25]}).
  \item \textbf{Whose perspective is foregrounded, challenged, or missing?} \textit{Foregrounded:} Dissenting liberal-constitutionalist intelligentsia (Karan Thapar \& Naseeruddin Shah) standing in moral solidarity with injured NEET protesters (\texttt{Sheikh Irshad Mansuri}, \texttt{[10:19]}) and jailed student activists (\texttt{CAA} parallel via Kapil Sibal, \texttt{[26:41--27:37]}). \textit{Challenged:} PM Modi's \texttt{"Pradhan Sevak"} / \texttt{"non-biological"} persona (\texttt{[00:28]}, \texttt{[23:47--24:54]}), Dharmendra Pradhan's X post promising \texttt{"answers, reforms and accountability"} (\texttt{[12:50--13:51]}), Delhi Police casualty claims (\texttt{[06:49--07:07]}), and silent Bollywood/elite peers (\texttt{[29:18--30:00]}). \textit{Missing:} Direct voices of CJP students on the ground, government/NTA spokespersons, and injured police officers.
  \item \textbf{What evidence supports your interpretation? What complicates it?} \textit{Supports:} Named \textit{Hindu} report on 25-year-old Sheikh Irshad Mansuri's neck pellet surgery (\texttt{[10:07--10:38]}), \textit{Hindustan Times} injury reports (\texttt{[06:33]}), historical cabinet resignation precedents (\texttt{Lal Bahadur Shastri 1956}, \texttt{Madhavrao Scindia 1992}, \texttt{Nitish Kumar 1999} cited at \texttt{[22:11--22:31]}), Dharmendra Pradhan's verbatim X post (\texttt{[13:00]}), and Kapil Sibal's legal analysis of unnamed FIRs (\texttt{[26:41--27:18]}). \textit{Complicates:} Shah's confident assertion that \texttt{"no way is he [Pradhan] going to resign"} (\texttt{[12:23--12:33]}) was disproven just 72 hours later; several anti-Modi personal attacks (\texttt{"patently fake degree"}, \texttt{"uneducated himself"}, marital/family remarks at \texttt{[18:23--19:01]}) drift into ad hominem polemic far beyond NEET; and Karan Thapar never challenges his guest once, producing an echo-chamber interview rather than adversarial questioning.
\end{itemize}
\vspace{4pt}\noindent\rule{\textwidth}{0.4pt}\vspace{4pt}

\section*{Common form: Blank excerpt record}
\textit{Copy or print this page whenever you find a passage worth analysing. Keep the exact wording and enough context to interpret it.}


\subsection*{Excerpt Record 1 (Event 1 --- Connecting PMO \& Nripendra Misra to Ayodhya Land/Donation Scandal)}
\begin{itemize}[leftmargin=1.5em, itemsep=3pt, topsep=4pt]
  \item \textbf{Outlet:} The Wire | \textbf{Date:} \texttt{2026-07-01}
  \item \textbf{Video title:} \textit{Ram Temple Donations Loot: Biggest Credibility Test For Modi-RSS}
  \item \textbf{Timestamp(s):} \texttt{[11:48--13:42]} | \textbf{Analyst:} Baani Solanki | \textbf{Language:} English (Roman script)
  \item \textbf{Exact transcript excerpt:}
    \begin{quote}\small\ttfamily [11:48] [Nilanjan Mukhopadhyay:] The government of India also decided to place the gentleman who had been the principal secretary of Mr. Modi in the years between 2014 and 2019, Mr. Nripendra [11:58] Misra, to personally oversee the construction of the Ram temple... [12:17] He was very much a part of the core team of Mr. Modi. Besides that, as the person who was [12:27] supervising the construction, he would have virtually gone through and decided, played a very decisive role in deciding where the material is going to be [12:36] bought, which land is going to be bought from whom, at what price... [12:55] artificial inflation of money so that people could personally benefit on it, Mr. Misra would have been in the know of [13:04] that. And I find it very hard to believe that if he was in the know of something, the others who had actually played a decisive role in appointing him, [13:15] you know, did not know what all was happening.\end{quote}
  \item \textbf{Feature(s):} $\boxtimes$ Naming $\square$ Presupposition $\boxtimes$ Agency $\square$ Antagonisation $\boxtimes$ Intensification / softening
  \item \textbf{Who is speaking? Is this the outlet's wording or a quotation?} Senior journalist and author Nilanjan Mukhopadhyay, prompted and affirmed by The Wire host M.K. Venu (\texttt{[13:24--13:42]}).
  \item \textbf{What does the wording imply, foreground, soften, or leave unstated?} Uses deductive modal chains (\texttt{"would have virtually gone through"}, \texttt{"would have been in the know"}, \texttt{"find it very hard to believe... did not know"}) to trace institutional agency from Ayodhya land/material procurement directly into \texttt{"the core team of Mr. Modi"} at the PMO---softened grammatically as logical inference (\texttt{would have}) rather than direct documentary proof.
  \item \textbf{What context or evidence supports or complicates this interpretation?} Supported by Nripendra Misra's undisputed status as ex-Principal Secretary to PM Modi (\texttt{2014--2019}) appointed to chair the Construction Committee (\texttt{[11:48--11:58]}). Complicated by \textbf{Aaj Tak} (\texttt{B4L1pz0o4dc}, \texttt{[04:35--06:05]}, \texttt{[10:57--11:24]}), where Misra himself publicly blew the whistle on counting-room informality (\texttt{"dacoity"}) and Anjana Om Kashyap argued he oversaw \textit{construction}, not \textit{Trust treasury/counting}.
  \item \textbf{How does another outlet describe the same event or actor?} \textbf{Aaj Tak} (\texttt{B4L1pz0o4dc}, \texttt{[10:57--11:24]}) explicitly exonerates Nripendra Misra (\texttt{"nirman alag hai, prabandhan alag hai"}), whereas \textbf{The Wire} and \textbf{Dhruv Rathee} (\texttt{oeH0fOSJuX8}, \texttt{[04:28--04:46]}) treat Misra's appointment as proof that the Trust was run from Delhi/PMO.
  \item \textbf{My interpretation, with reasons:} This passage exemplifies The Wire's analytical register: instead of shouting \texttt{"chor"} (like TV panelists), Mukhopadhyay uses bureaucratic command-chain logic to pierce PMO plausible deniability.
\end{itemize}

\subsection*{Excerpt Record 2 (Event 1 --- Sangh Parivar "Shadow Organizations" \& Strategic Ambiguity)}
\begin{itemize}[leftmargin=1.5em, itemsep=3pt, topsep=4pt]
  \item \textbf{Outlet:} The Wire | \textbf{Date:} \texttt{2026-07-01}
  \item \textbf{Video title:} \textit{Ram Temple Donations Loot: Biggest Credibility Test For Modi-RSS}
  \item \textbf{Timestamp(s):} \texttt{[17:24--18:40]} \& \texttt{[24:36--25:01]} | \textbf{Analyst:} Baani Solanki | \textbf{Language:} English (Roman script)
  \item \textbf{Exact transcript excerpt:}
    \begin{quote}\small\ttfamily [17:24] [M.K. Venu:] RSS always believed in an unstructured existence... [17:43] whatever money comes to the Sangh, RSS's Guru Dakshina, therefore it is non-taxable and it can be [17:51] spent in whatever way they want to, and they need not keep any accounts. And they need not even register themselves... [18:30] would you say that this is a result of this history where the RSS or the Sangh Parivar, they refuse to be accountable?\end{quote}
    \begin{quote}\small\ttfamily [24:36] [Nilanjan Mukhopadhyay:] Setting up structures was not done. It [24:44] was probably not done because they wanted that it should not be a leak-proof system. The moment you structure something, then it becomes [24:53] leak-proof. So you keep a certain amount of ambiguity, a certain amount of fluidity.\end{quote}
  \item \textbf{Feature(s):} $\boxtimes$ Naming $\boxtimes$ Presupposition $\boxtimes$ Agency $\boxtimes$ Antagonisation $\boxtimes$ Intensification / softening
  \item \textbf{Who is speaking? Is this the outlet's wording or a quotation?} Host M.K. Venu (\texttt{[17:24--18:40]}) and guest Nilanjan Mukhopadhyay (\texttt{[24:36--25:01]}).
  \item \textbf{What does the wording imply, foreground, soften, or leave unstated?} Reframes administrative "lack of SOPs" from accidental incompetence into \textbf{purposeful institutional design} (\texttt{"they wanted that it should not be a leak-proof system"}, \texttt{"keep a certain amount of ambiguity"}), rooted in a 40-year history of unregistered Sangh bodies (\texttt{Guru Dakshina}, \texttt{Mukti Yagna Samiti}, \texttt{[17:43]}, \texttt{[20:46--21:20]}). Note the careful epistemic hedge \texttt{"probably"} (\texttt{[24:44]}) softening an otherwise explosive claim of intentionality.
  \item \textbf{What context or evidence supports or complicates this interpretation?} Supported by the 1990 \texttt{Rs.\,700 crore} VHP tax-notice episode (\texttt{[15:33--17:15]}) and the November 2020 private audit warning that went unimplemented for 6 years (cited in Dhruv Rathee \texttt{oeH0fOSJuX8} \texttt{[10:12--10:32]} and Republic TV panelist Ashish Chaturvedi \texttt{0H-bruoEkaY} \texttt{[27:34--27:53]}).
  \item \textbf{How does another outlet describe the same event or actor?} \textbf{Aaj Tak}'s VHP spokesperson (\texttt{B4L1pz0o4dc}, \texttt{[34:21]}) calls the exact same lack of structure harmless \texttt{"kuch avyavastha"} [some disarray], while \textbf{Republic TV} (\texttt{0H-bruoEkaY}, \texttt{[26:56]}) attributes it merely to 79-year-old Champat Rai being \texttt{"not skilled administratively"}.
  \item \textbf{My interpretation, with reasons:} Comparing \texttt{"kuch avyavastha"} (Aaj Tak/VHP) with \texttt{"they wanted that it should not be a leak-proof system"} (The Wire) shows how two outlets observe the \textit{exact same absence of accounting controls} and assign opposite moral agency.
\end{itemize}

\subsection*{Excerpt Record 3 (Event 1 --- "Ayodhya the Hindu Vatican" \& Hollowed-Out Devotee Faith)}
\begin{itemize}[leftmargin=1.5em, itemsep=3pt, topsep=4pt]
  \item \textbf{Outlet:} The Wire | \textbf{Date:} \texttt{2026-07-01}
  \item \textbf{Video title:} \textit{Ram Temple Donations Loot: Biggest Credibility Test For Modi-RSS}
  \item \textbf{Timestamp(s):} \texttt{[35:04--37:08]} | \textbf{Analyst:} Baani Solanki | \textbf{Language:} English (Roman script)
  \item \textbf{Exact transcript excerpt:}
    \begin{quote}\small\ttfamily [35:04] [Nilanjan Mukhopadhyay:] Ayodhya in a way, Ayodhya the [35:13] temple, Ayodhya the Hindu Vatican is actually the Sangh Parivar's creation. So, within the Hindu Vatican, it does [35:22] not give any space to these established, you know, our temporal [spiritual] heads... [36:01] the moment people lose their rationality, they become bhakts. [36:11] We really do not know to what extent people have become bhakts. That they are willing to [36:20] overlook anything that happens... [36:28] they put their body and soul into this movement. You remember the Ram Shila Yatras of 1989... [36:59] And now you find that the darshan has been hollowed out, that it has been stolen.\end{quote}
  \item \textbf{Feature(s):} $\boxtimes$ Naming $\boxtimes$ Presupposition $\square$ Agency $\boxtimes$ Antagonisation $\boxtimes$ Intensification / softening
  \item \textbf{Who is speaking? Is this the outlet's wording or a quotation?} Nilanjan Mukhopadhyay speaking in dialogue with M.K. Venu.
  \item \textbf{What does the wording imply, foreground, soften, or leave unstated?} Coins the striking metaphor \texttt{"Ayodhya the Hindu Vatican"} (\texttt{[35:13]}) to argue that the Sangh Parivar centralized a pluralistic Hindu tradition (historically centred in Varanasi, Mathura, Haridwar) into a politically controlled headquarters while sidelining traditional \textit{acharyas} (\texttt{[34:07--35:22]}). Simultaneously distinguishes blind \texttt{"bhakts"} (\texttt{"lose their rationality"}, \texttt{[36:01]}) from sincere grassroots \texttt{1989} \textit{Ram Shila} devotees whose sacred \texttt{"darshan has been hollowed out"} (\texttt{[36:59]}).
  \item \textbf{What context or evidence supports or complicates this interpretation?} Supported by public criticism from traditional \textit{akhadas} (\texttt{Nirmohi Akhada} at \texttt{[29:20]}) and \textit{Shankaracharyas} (also referenced by SP's Anurag Bhadauria on Aaj Tak \texttt{B4L1pz0o4dc} \texttt{[15:17]}). Complicated by Mukhopadhyay's own candid admission (\texttt{[37:08--37:19]}) that \texttt{"We do not have any empirical data as yet"} on whether devotees will actually turn against the BJP.
  \item \textbf{How does another outlet describe the same event or actor?} \textbf{Sarthak Goswami} (\texttt{lb7gZnyZH0c}, \texttt{[06:26]}) uses a remarkably parallel Hindi formulation---\texttt{"Ram ka chanda nahin Ram ka charitra chori ho chuka hai"}---to express the same "hollowed-out sacredness" in popular Hinglish rather than academic English.
  \item \textbf{My interpretation, with reasons:} By separating pluralistic Hinduism and sincere \textit{Ram Shila} pilgrims from \texttt{"the Hindu Vatican"} of the Sangh Parivar, The Wire critiques Hindutva institutional power while empathizing with ordinary believers' sense of betrayal.
\end{itemize}

\subsection*{Excerpt Record 4 (Event 2 --- Unbadged "Army of Goons", Neck Pellets \& Student Retaliation)}
\begin{itemize}[leftmargin=1.5em, itemsep=3pt, topsep=4pt]
  \item \textbf{Outlet:} The Wire (\textit{The Interview}) | \textbf{Date:} \texttt{2026-07-22}
  \item \textbf{Video title:} \textit{I stand with the protesting students; this is not the India I love; I don't take Modi seriously}
  \item \textbf{Timestamp(s):} \texttt{[07:15--07:52]}, \texttt{[09:36--09:57]} \& \texttt{[10:07--10:38]} | \textbf{Analyst:} Baani Solanki | \textbf{Language:} English (Roman script)
  \item \textbf{Exact transcript excerpt:}
    \begin{quote}\small\ttfamily [07:15] [Naseeruddin Shah:] these people who are sadistically beating up these students who are asking [07:25] for their rights are young people themselves most of the time. A lot of them don't look like policemen... [07:34] they're not carrying name badges, they're not carrying ranks, they're not even in proper uniform. Some of them are even admitting that they are private [07:42] privately hired. So to let loose an army of goons on our future...\end{quote}
    \begin{quote}\small\ttfamily [09:36] And what I watched this morning is the students retaliating. It's not surprising. I mean, even a rat fights back when it's [09:46] pushed against the wall. And these kids are justified in what they're asking. Their anger is absolutely valid...\end{quote}
    \begin{quote}\small\ttfamily [10:07] [Karan Thapar:] today's Hindu reports that pellet guns and shock batons were used on some of the students. [10:19] It reports that Sheikh Irshad Mansuri, who's only 25 years old, had to undergo surgery for pellets lodged in the upper part of his body and his neck to be removed.\end{quote}
  \item \textbf{Feature(s):} $\boxtimes$ Naming $\boxtimes$ Presupposition $\boxtimes$ Agency $\boxtimes$ Antagonisation $\boxtimes$ Intensification / softening
  \item \textbf{Who is speaking? Is this the outlet's wording or a quotation?} Guest Naseeruddin Shah (\texttt{[07:15--09:57]}) and The Wire interviewer Karan Thapar citing \textit{The Hindu} (\texttt{[10:07--10:38]}).
  \item \textbf{What does the wording imply, foreground, soften, or leave unstated?} Combines Shah's moral-emotional intensification (\texttt{"sadistically beating"}, \texttt{"an army of goons on our future"}) with Karan Thapar's hard newspaper citation (\textit{The Hindu} on 25-year-old \texttt{Sheikh Irshad Mansuri}'s neck pellet surgery). Crucially, student counter-violence is explicitly softened and morally excused via \texttt{"retaliating"} and the cornered-victim idiom \texttt{"even a rat fights back when it's pushed against the wall"} (\texttt{[09:36--09:46]}).
  \item \textbf{What context or evidence supports or complicates this interpretation?} \textit{The Hindu}'s medical detail on Sheikh Irshad Mansuri provides concrete verification of pellet-gun injuries cited more generally by \textbf{Sarthak Goswami} (\texttt{EM1GeZufO6E}, \texttt{[03:39--03:47]}). Meanwhile, Shah's dismissal of police injuries at \texttt{[06:49--07:07]} (\texttt{"supposedly more policemen injured... something like 600"}) exaggerates/garbles the official \texttt{110--118} figure documented with X-rayed fractures on \textbf{Aaj Tak} (\texttt{JXLIs0WtJAM}, \texttt{[18:09]}, \texttt{[21:12--23:23]}).
  \item \textbf{How does another outlet describe the same event or actor?} Notice the exact lexical mirror image with \textbf{Republic TV} (\texttt{HTvb6E\_1Qdc}, \texttt{[04:19]}): Arnab calls protesters \texttt{"the goons of Abhijit Dipke"} attacking uniformed officers, whereas Naseeruddin Shah calls unbadged state-aligned assailants \texttt{"an army of goons"} let loose on \texttt{"our future"}!
  \item \textbf{My interpretation, with reasons:} Pairing Karan Thapar's print-journalism verification (\textit{The Hindu}) with Shah's visceral moral outrage gives The Wire both factual anchoring and emotional force, while clearly asymmetrically excusing student \texttt{"retaliating"}.
\end{itemize}

\subsection*{Excerpt Record 5 (Event 2 --- Deconstructing Dharmendra Pradhan's X Post as Orwellian "Double Speak")}
\begin{itemize}[leftmargin=1.5em, itemsep=3pt, topsep=4pt]
  \item \textbf{Outlet:} The Wire (\textit{The Interview}) | \textbf{Date:} \texttt{2026-07-22}
  \item \textbf{Video title:} \textit{I stand with the protesting students; this is not the India I love; I don't take Modi seriously}
  \item \textbf{Timestamp(s):} \texttt{[12:50--14:20]} | \textbf{Analyst:} Baani Solanki | \textbf{Language:} English (Roman script)
  \item \textbf{Exact transcript excerpt:}
    \begin{quote}\small\ttfamily [12:50] [Karan Thapar:] Yesterday the education minister posted on X. This is one of the [13:00] things he said. Speaking to the students, he wrote: "We owe them answers, reforms and accountability. This is exactly what our government [13:10] remains committed to delivering." Accountable is the one thing the minister refuses to be and his silence suggests he refuses to answer their [13:21] questions. How do you view his statement?\end{quote}
    \begin{quote}\small\ttfamily [Naseeruddin Shah:] This is classic double speak. What it was called in 1984, the book, [13:33] classic double speak, Newspeak... He's sticking to the vocabulary which he hopes will mollify some... [14:02] we are living in 1984 where there's surveillance on everybody. 2 minutes of hate has been transformed [14:11] into 24 hours of hate. And it's just completely turning [14:20] into a dystopian scenario here.\end{quote}
  \item \textbf{Feature(s):} $\boxtimes$ Naming $\boxtimes$ Presupposition $\square$ Agency $\boxtimes$ Antagonisation $\boxtimes$ Intensification / softening
  \item \textbf{Who is speaking? Is this the outlet's wording or a quotation?} Karan Thapar quoting Dharmendra Pradhan's official X post (\texttt{[13:00--13:10]}) and adding his own editorial commentary (\texttt{[13:10--13:21]}), followed by Naseeruddin Shah's analysis (\texttt{[13:21--14:20]}).
  \item \textbf{What does the wording imply, foreground, soften, or leave unstated?} Notice how Karan Thapar pre-frames the quote \textit{before} asking his question (\texttt{"Accountable is the one thing the minister refuses to be"}---leaving zero interpretive neutrality). Shah then performs explicit \textbf{critical discourse analysis} using George Orwell's \textit{1984} (\texttt{"classic double speak, Newspeak"}, \texttt{"24 hours of hate"}), arguing that state actors deploy hollow democratic buzzwords (\texttt{answers, reforms, accountability}) while practising surveillance and repression.
  \item \textbf{What context or evidence supports or complicates this interpretation?} Supported by the glaring 30-day gap during which Dharmendra Pradhan never once visited protesting students at Jantar Mantar (\texttt{[20:35--20:54]}, corroborated by citizen interviews in Mohak Mangal \texttt{v-RikSqb0Vw} \texttt{[05:31--05:38]}). Complicated by the fact that 3 days after this interview (\texttt{25 July}), Pradhan \textit{did} resign (\texttt{EM1GeZufO6E} \texttt{[00:08]}), showing democratic pressure still functioned beyond a total \textit{1984} dystopia.
  \item \textbf{How does another outlet describe the same event or actor?} \textbf{Mohak Mangal} (\texttt{v-RikSqb0Vw}, \texttt{[00:00--01:03]}) critiques Dharmendra Pradhan not through Western literary theory (\textit{Orwell's 1984}), but through Indian biographical irony (his \texttt{1997} Odisha ABVP lathi-charge injuries)---appealing to a much broader Hindi/Hinglish youth audience.
  \item \textbf{My interpretation, with reasons:} This excerpt showcases The Wire's cultural register---using cosmopolitan English literary touchstones (\textit{Orwell's 1984}, \texttt{Newspeak}) and pre-loaded interviewer prompts to indict official government communication.
\end{itemize}

\subsection*{Excerpt Record 6 (Event 2 --- Citizens' Right to Parliament \& Kapil Sibal's UAPA Warning)}
\begin{itemize}[leftmargin=1.5em, itemsep=3pt, topsep=4pt]
  \item \textbf{Outlet:} The Wire (\textit{The Interview}) | \textbf{Date:} \texttt{2026-07-22}
  \item \textbf{Video title:} \textit{I stand with the protesting students; this is not the India I love; I don't take Modi seriously}
  \item \textbf{Timestamp(s):} \texttt{[25:02--27:37]} | \textbf{Analyst:} Baani Solanki | \textbf{Language:} English (Roman script)
  \item \textbf{Exact transcript excerpt:}
    \begin{quote}\small\ttfamily [25:02] [Karan Thapar:] Do peaceful nonviolent Indian protesters protesting in a just cause have the [25:13] right to walk to parliament and demand the attention of their MPs and government? The government's answer is clearly they don't. In which case, how [25:23] meaningful is our democracy?\end{quote}
    \begin{quote}\small\ttfamily [Naseeruddin Shah:] That question answers itself. Karan, this parliament is run with our [25:33] taxpayer money. These people are our representatives. That house of parliament is ours. It belongs to the kids. It belongs to the [25:42] future. They have an absolute right to enter it... [26:32] This is not the India I was taught to love.\end{quote}
    \begin{quote}\small\ttfamily [26:41] [Karan Thapar:] The police have already filed FIRs against unnamed persons. Kapil Sibal fears that in due course [26:58] they will find names and possibly charge people under UAPA or section 152 of the old penal code [27:07] which is now section 195 of BNS... that's exactly what happened 6 years ago when [27:18] the CAA protests occurred. And those young people from the CAA days are still rotting [ASR: pushing] in jail.\end{quote}
  \item \textbf{Feature(s):} $\boxtimes$ Naming $\boxtimes$ Presupposition $\boxtimes$ Agency $\boxtimes$ Antagonisation $\boxtimes$ Intensification / softening
  \item \textbf{Who is speaking? Is this the outlet's wording or a quotation?} Karan Thapar (\texttt{[25:02--25:23]}, \texttt{[26:41--27:27]}) and Naseeruddin Shah (\texttt{[25:23--26:32]}).
  \item \textbf{What does the wording imply, foreground, soften, or leave unstated?} Reclaims democratic ownership of physical state space (\texttt{"That house of parliament is ours. It belongs to the kids"}, \texttt{[25:33--25:42]}) against security-state exclusion, and links the 5 open-ended Delhi Police FIRs (reported positively on Aaj Tak \texttt{JXLIs0WtJAM} \texttt{[19:25--20:22]}) to a chilling pattern of post-protest criminalization (\texttt{UAPA} / \texttt{BNS 195} pre-trial detention since the 2020 CAA protests).
  \item \textbf{What context or evidence supports or complicates this interpretation?} Supported by the actual filing of 5 FIRs with unnamed accused (\texttt{Aaj Tak} \texttt{[19:25--20:22]}) and by CJP making \textit{withdrawal/non-filing of police cases} one of its core settlement conditions before ending the protest on 25 July (\texttt{Republic TV} \texttt{HTvb6E\_1Qdc} \texttt{[00:07--00:15]}; \texttt{Sarthak Goswami} \texttt{EM1GeZufO6E} \texttt{[08:25]})! Complicated legally by the fact that even in constitutional democracies, unrestricted mass entry into an active Parliament building (\texttt{"an absolute right to enter it"}, \texttt{[25:42]}) is subject to reasonable security regulation.
  \item \textbf{How does another outlet describe the same event or actor?} \textbf{Aaj Tak} (\texttt{JXLIs0WtJAM}, \texttt{[19:25--20:41]}) presents the exact same 5 FIRs and Parliament cordons neutrally as routine \texttt{"ligal ekshan"} around \texttt{"vaital instoleshanans"}, whereas \textbf{The Wire} exposes how open-ended FIRs can become multi-year \texttt{UAPA} traps for students.
  \item \textbf{My interpretation, with reasons:} Highlighting the danger of unnamed FIRs proved prescient---explaining why CJP refused to vacate Jantar Mantar even after Pradhan resigned until the government explicitly guaranteed no cases would be pursued (\texttt{HTvb6E\_1Qdc}, \texttt{[00:07--00:15]}).
\end{itemize}
\vspace{4pt}\noindent\rule{\textwidth}{0.4pt}\vspace{4pt}

\section*{Common form: End-of-event summary}

\subsection*{End-of-Event Summary --- Event 1 (Ram Mandir Donation Controversy)}
\begin{itemize}[leftmargin=1.5em, itemsep=3pt, topsep=4pt]
  \item \textbf{The main way this outlet framed the event was:} As a historic moral and institutional unraveling of the Sangh Parivar's core 40-year project (\texttt{"Biggest Credibility Test For Modi-RSS"}, \texttt{"biggest challenge... after the assassination of Mahatma Gandhi"}, \texttt{[08:23--08:31]}), arguing that massive donation \texttt{"dacoity"} and \texttt{5x--10x} land-price \textit{hawala} inflation stemmed directly from RSS's unaccountable \texttt{"shadow organization"} culture (\texttt{[17:24--22:32]}), intentional avoidance of \texttt{"leak-proof"} structures (\texttt{[24:44--25:01]}), and PMO-linked oversight under Nripendra Misra (\texttt{[11:48--14:55]}).
  \item \textbf{The strongest linguistic evidence was (quote + timestamp):} \texttt{"Setting up structures was not done. It was probably not done because they wanted that it should not be a leak-proof system... Ayodhya the Hindu Vatican is actually the Sangh Parivar's creation... And now you find that the darshan has been hollowed out, that it has been stolen."} (\texttt{[24:36--24:53]}, \texttt{[35:13]} \& \texttt{[36:59]}).
  \item \textbf{The shared feature(s) most clearly illustrated were:} \textbf{Naming} (\texttt{shadow organizations}, \texttt{corporate bureaucratic structure}, \texttt{political loyalists}, \texttt{Ayodhya the Hindu Vatican}), \textbf{Agency} (attributing lack of accounting controls to deliberate Sangh design rather than clerical oversight at \texttt{[24:44--25:01]}), and \textbf{Epistemic Hedging vs Historical Intensification} (\texttt{potentially}, \texttt{no empirical data} paired with \texttt{100, maybe 1,000 questions} and \texttt{1948} parallels).
  \item \textbf{An alternative interpretation of the same evidence is:} That transitioning from an agitational movement to managing millions of pilgrims created genuine early administrative bottlenecks rather than intentional \texttt{"non-leak-proof"} conspiracy, and that Nripendra Misra's public admission of \texttt{"dacoity"} alongside 8 SIT arrests and senior resignations reflected corrective institutional action rather than PMO complicity.
  \item \textbf{What I cannot conclude from this transcript alone:} Who the ultimate financial beneficiaries of the inflated land transactions were (\texttt{[04:48]}), whether Nripendra Misra actually briefed PM Modi on donation irregularities, or whether voter outrage will translate into electoral losses in UP (\texttt{[37:08--37:28]}).
\end{itemize}

\subsection*{End-of-Event Summary --- Event 2 (Students' Protest Regarding NEET Paper Leak)}
\begin{itemize}[leftmargin=1.5em, itemsep=3pt, topsep=4pt]
  \item \textbf{The main way this outlet framed the event was:} As an agonizing democratic and moral crisis (\texttt{"this is not the India I was taught to love"}, \texttt{"living in 1984"}, \texttt{[14:02]}, \texttt{[26:32]}) in which \texttt{"arrogant autocrats"} (\texttt{[23:47]}) ignored wronged students for 25 days, offered Orwellian \texttt{"double speak"} (\texttt{[13:21]}), and unleashed tear gas, neck-lodging pellet guns (\texttt{Sheikh Irshad Mansuri}, \texttt{[10:19]}), and unbadged \texttt{"lumpens"} / \texttt{"an army of goons"} (\texttt{[07:42]}, \texttt{[11:26]}) on peaceful teenagers rightfully marching toward citizens' Parliament.
  \item \textbf{The strongest linguistic evidence was (quote + timestamp):} \texttt{"to let loose an army of goons on our future... This is classic double speak. What it was called in 1984... That house of parliament is ours. It belongs to the kids. It belongs to the future."} (\texttt{[07:42]}, \texttt{[13:21--13:33]} \& \texttt{[25:33--25:42]}).
  \item \textbf{The shared feature(s) most clearly illustrated were:} \textbf{Presupposition} (Karan Thapar's heavily loaded negative-interrogative and declarative prompts at \texttt{[06:41]}, \texttt{[13:10]}, \texttt{[22:40]}, \texttt{[23:47]}), \textbf{Naming} (\texttt{army of goons} and \texttt{lumpens} for state-aligned attackers vs \texttt{these kids} / \texttt{these bachas} / \texttt{our future} for protesters), and \textbf{Agency} (framing student clashes strictly as justified cornered self-defense---\texttt{"even a rat fights back when it's pushed against the wall"}, \texttt{[09:36--09:46]}).
  \item \textbf{An alternative interpretation of the same evidence is:} That security cordons prevented a chaotic physical breach of Parliament during session, that over \texttt{110} uniformed police officers suffered genuine fractures/injuries from crowd stone-pelting (dismissed skeptically as \texttt{"supposedly"} by Shah at \texttt{[06:49]}), and that Indian democratic accountability was far less closed than \textit{1984}---proven when Dharmendra Pradhan resigned just 3 days after this interview.
  \item \textbf{What I cannot conclude from this transcript alone:} What protesting students on the ground prioritized in negotiations (since no student was interviewed), whether the unbadged men were formally police personnel or privately paid recruits (\texttt{[07:34--07:42]}), or why the government reversed course to accept CJP's demands on 25 July.
\end{itemize}
\textit{Research reminder: interpret wording in context, distinguish allegations from established facts, and support claims with exact excerpts. Mixed or contradictory evidence is worth recording.}


\end{document}
